\input{preamble}

% OK, start here.
%
\begin{document}

\title{Résolution des surfaces}


\maketitle

\phantomsection
\label{section-phantom}

\tableofcontents

\section{Introduction} \label{section-introduction}

\noindent Ce chapitre traite de la résolution des singularités des surfaces en suivant Lipman \cite{Lipman} et, pour l'essentiel, l'exposé d'Artin dans \cite{Artin-Lipman}. Le résultat principal (théorème \ref{theorem-resolve}) affirme qu'un schéma noethérien de dimension $2$, noté $Y$, admet une résolution des singularités lorsqu'il possède une normalisation finie $Y^\nu \to Y$ ayant un nombre fini de points singuliers $y_i \in Y^\nu$ et que, pour tout $i$, le complété $\mathcal{O}_{Y^\nu, y_i}^\wedge$ est normal.

\medskip\noindent En particulier, si $Y$ est un schéma de dimension $2$ de type fini sur un anneau de base quasi-excellent $R$ (par exemple un corps ou un anneau de Dedekind dont le corps des fractions est de caractéristique $0$, tel que $\mathbf{Z}$), alors la normalisation de $Y$ est finie, possède un nombre fini de points singuliers et les complétés de ses anneaux locaux sont normaux. Voir la discussion dans Compléments d'algèbre, sections \ref{more-algebra-section-singular-locus}, \ref{more-algebra-section-G-ring} et \ref{more-algebra-section-excellent}, ainsi que Compléments d'algèbre, lemme \ref{more-algebra-lemma-normal-goes-up}. Un tel $Y$ admet donc une résolution des singularités.

\medskip\noindent Voici les grandes lignes de la démonstration. Soit $A$ un anneau local noethérien intègre de dimension $2$. Les étapes de la démonstration sont les suivantes : \begin{enumerate} \item[N] remplacer $A$ par son normalisé ; \item[V] démontrer le théorème de Grauert-Riemenschneider ; \item[B] montrer qu'il existe un maximum $g$ des longueurs de $H^1(X, \mathcal{O}_X)$ lorsque l'on parcourt toutes les modifications normales $X \to \Spec(A)$, et se ramener au cas $g = 0$ ; \item[R] on dit que $A$ définit une singularité rationnelle si $g = 0$ ; dans ce cas, après un nombre fini d'éclatements, on peut supposer que $A$ est de Gorenstein et que $g = 0$ ; \item[D] on dit que $A$ définit un point double rationnel si $g = 0$ et si $A$ est de Gorenstein ; dans ce cas, on résout explicitement les singularités. \end{enumerate} Chacune de ces étapes nécessite des hypothèses sur l'anneau $A$. Nous les examinerons successivement.

\medskip\noindent Étape N : il faut ici supposer que $A$ possède une normalisation finie (ce qui n'est pas automatique). Dans la majeure partie de ce chapitre, nous supposerons que notre schéma est de Nagata lorsqu'il nous faudra savoir qu'une normalisation est finie. Toutefois, être de Nagata est une condition légèrement plus forte que celle qui figure dans l'énoncé du théorème. Une solution à ce léger problème aurait consisté à utiliser le fait que notre anneau $A$ est formellement non ramifié (c'est-à-dire que son complété est réduit) et à appliquer le lemme \ref{lemma-formally-unramified}. Cependant, avec la démonstration adoptée ici, il s'avère plus simple d'utiliser le lemme \ref{lemma-normalization-completion} pour déduire de la finitude de la normalisation au-dessus du complété celle de la normalisation au-dessus de $A$.

\medskip\noindent Étape V : il s'agit de la proposition \ref{proposition-Grauert-Riemenschneider}, qui affirme, en substance, que pour une modification normale $f : X \to \Spec(A)$ on a $R^1f_*\omega_X = 0$, où $\omega_X$ est le module dualisant de $X/A$ (remarque \ref{remark-dualizing-setup}). En fait, par dualité, ce résultat équivaut à un énoncé (lemme \ref{lemma-R1-injective}) concernant l'objet $Rf_*\mathcal{O}_X$ de la {\it catégorie dérivée} $D(A)$. Cela dit, la démonstration utilise le fait classique que les composantes de la fibre spéciale ont des faisceaux conormaux positifs (lemme \ref{lemma-nontrivial-normal-bundle}).

\medskip\noindent Étape B : c'est, en un sens, la partie la plus délicate de la démonstration. Nous n'aurons finalement besoin du résultat de cette étape que lorsque $A$ est un anneau local noethérien complet, bien que l'exposé soit un peu plus général. La terminologie est fixée dans la définition \ref{definition-reduce-to-rational}. Si $g$ (défini ci-dessus) est borné, un argument simple montre que l'on peut trouver une modification normale $X \to \Spec(A)$ telle que tous les points singuliers de $X$ soient des singularités rationnelles ; voir le lemme \ref{lemma-reduce-to-rational}. Nous montrons que, pour une extension finie $A \subset B$, l'invariant $g$ est borné pour $B$ s'il l'est pour $A$ dans les deux cas suivants : (1) lorsque l'extension des corps des fractions est séparable, voir le lemme \ref{lemma-reduce-to-rational} ; (2) lorsque l'extension des corps des fractions est de degré $p$, que la caractéristique est $p$ et que $A$ est régulier et complet, voir le lemme \ref{lemma-go-up-degree-p}.

\medskip\noindent Étape R : nous ramenons ici le cas $g = 0$ au cas de Gorenstein. Le fait remarquable qui permet de mener à bien la démonstration est que l'éclatement d'une singularité rationnelle de surface est normal ; voir le lemme \ref{lemma-blow-up-normal-rational}.

\medskip\noindent Étape D : la résolution des points doubles rationnels se fait par des calculs assez explicites ; voir la section \ref{section-rational-double-points}. Un point double rationnel est une singularité d'hypersurface (c'est vrai, mais nous ne le démontrons pas, car nous n'en avons pas besoin). Son équation locale est de la forme $$
a_{11} x_1^2 + a_{12} x_1x_2 + a_{13}x_1x_3 + a_{22} x_2^2 +
a_{23} x_2x_3 + a_{33} x_3^2 =
\sum a_{ijk} x_ix_jx_k
$$ Le fait que la partie quadratique ne puisse être nulle, puisque la multiplicité vaut $2$ et reste égale à $2$ après tout éclatement, joint à la normalité de chaque éclatement, permet d'obtenir rapidement une résolution. Une particularité est que nous ne disposons pas d'un invariant décroissant à chaque éclatement ; nous nous appuyons sur les résultats relatifs aux arcs formels de la section \ref{section-arcs} pour démontrer que le procédé s'arrête.

\medskip\noindent Pour réunir tous ces éléments, un travail supplémentaire est nécessaire. La difficulté principale consiste à relever une résolution du complété $A^\wedge$ en une résolution de $\Spec(A)$. Pour cela, nous montrons d'abord que, s'il existe une résolution, il en existe une par éclatements normalisés (lemme \ref{lemma-existence-implies-existence-by-normalized-blowing-ups}). Une suite d'éclatements normalisés peut être relevée à partir du complété grâce au lemme \ref{lemma-normalized-blowup-completion}. Nous utilisons ensuite ce fait jusque dans la démonstration de la résolution des anneaux locaux complets $A$, car notre méthode procède par récurrence sur le degré d'une inclusion finie $A_0 \subset A$ avec $A_0$ régulier ; voir le lemme \ref{lemma-resolve-complete}. Un résultat plus fort à l'étape B (tel que celui démontré dans l'article de Lipman) permettrait d'éviter cette étape.




\section{Une application trace en caractéristique positive} \label{section-trace}

\noindent Certains résultats de cette section se déduisent de l'étude beaucoup plus générale des traces sur les formes différentielles dans Cohomologie de de Rham, section \ref{derham-section-trace}. Voir la remarque \ref{remark-compare-Garel} pour une discussion.

\medskip\noindent Fixons un nombre premier $p$. Soit $R$ une $\mathbf{F}_p$-algèbre. Étant donné $a \in R$, posons $S = R[x]/(x^p - a)$. Définissons une application $R$-linéaire $$
\text{Tr}_x : \Omega_{S/R} \longrightarrow \Omega_R
$$ par la règle $$
x^i\text{d}x \longmapsto
\left\{
\begin{matrix}
0 & \text{si} & 0 \leq i \leq p - 2, \\
\text{d}a & \text{si} & i = p - 1
\end{matrix}
\right.
$$ Cela a un sens, car $\Omega_{S/R}$ est un $R$-module libre de base $x^i\text{d}x$, $0 \leq i \leq p - 1$. Le lemme suivant implique que l'application trace est bien définie, c'est-à-dire indépendante du choix de la coordonnée $x$.

\begin{lemma} \label{lemma-trace-well-defined} Soit $\varphi : R[x]/(x^p - a) \to R[y]/(y^p - b)$ un homomorphisme de $R$-algèbres. Alors $\text{Tr}_x = \text{Tr}_y \circ \varphi$. \end{lemma}

\begin{proof} Écrivons $\varphi(x) = \lambda_0 + \lambda_1 y + \ldots + \lambda_{p - 1}y^{p - 1}$ avec $\lambda_i \in R$. Dire que l'envoi de $x$ sur $\lambda_0 + \lambda_1 y + \ldots + \lambda_{p - 1}y^{p - 1}$ induit un homomorphisme de $R$-algèbres $R[x]/(x^p - a) \to R[y]/(y^p - b)$ équivaut à la condition $$
a = \lambda_0^p + \lambda_1^p b + \ldots + \lambda_{p - 1}^pb^{p - 1}
$$ dans l'anneau $R$. Considérons l'anneau de polynômes $$
R_{univ} = \mathbf{F}_p[b, \lambda_0, \ldots, \lambda_{p - 1}]
$$ muni de l'élément $a = \lambda_0^p + \lambda_1^p b + \ldots + \lambda_{p - 1}^pb^{p - 1}$ Considérons l'homomorphisme universel d'algèbres $\varphi_{univ} : R_{univ}[x]/(x^p - a) \to R_{univ}[y]/(y^p - b)$ défini par l'envoi de $x$ sur $\lambda_0 + \lambda_1 y + \ldots + \lambda_{p - 1}y^{p - 1}$. On obtient une application canonique $$
R_{univ} \longrightarrow R
$$ envoyant $b, \lambda_i$ sur $b, \lambda_i$. Par construction, on a un diagramme commutatif $$
\xymatrix{
R_{univ}[x]/(x^p - a) \ar[r] \ar[d]_{\varphi_{univ}} &
R[x]/(x^p - a) \ar[d]^\varphi \\
R_{univ}[y]/(y^p - b) \ar[r] & R[y]/(y^p - b)
}
$$ dont les flèches horizontales sont compatibles aux applications trace. Il suffit donc de démontrer le lemme pour l'application $\varphi_{univ}$. On peut ainsi supposer que $R = \mathbf{F}_p[b, \lambda_0, \ldots, \lambda_{p - 1}]$ est un anneau de polynômes. Nous vérifierons le lemme dans ce cas en calculant $\text{Tr}_y(\varphi(x)^i\text{d}\varphi(x))$ pour $i = 0 , \ldots, p - 1$.

\medskip\noindent Cas $0 \leq i \leq p - 2$. Développons $$
(\lambda_0 + \lambda_1 y + \ldots + \lambda_{p - 1}y^{p - 1})^i
(\lambda_1 + 2 \lambda_2 y + \ldots + (p - 1)\lambda_{p - 1}y^{p - 2})
$$ dans l'anneau $R[y]/(y^p - b)$. Il faut montrer que le coefficient de $y^{p - 1}$ est nul. Pour cela, il suffit de montrer que, dans l'expression ci-dessus considérée comme polynôme en $y$, les coefficients des puissances $y^{pk - 1}$ sont nuls. Écrivons alors notre polynôme sous la forme $$
\frac{\text{d}}{(i + 1)\text{d}y}
(\lambda_0 + \lambda_1 y + \ldots + \lambda_{p - 1}y^{p - 1})^{i + 1}
$$ Les coefficients de $y^{kp - 1}$ sont effectivement tous nuls.

\medskip\noindent Cas $i = p - 1$. Développons $$
(\lambda_0 + \lambda_1 y + \ldots + \lambda_{p - 1}y^{p - 1})^{p - 1}
(\lambda_1 + 2 \lambda_2 y + \ldots + (p - 1)\lambda_{p - 1}y^{p - 2})
$$ dans l'anneau $R[y]/(y^p - b)$. Pour achever la démonstration, il faut montrer que le coefficient de $y^{p - 1}$ multiplié par $\text{d}b$ vaut $\text{d}a$. Nous utilisons ici le fait que $R$ est égal à $S/pS$, où $S = \mathbf{Z}[b, \lambda_0, \ldots, \lambda_{p - 1}]$. L'expression ci-dessus, considérée comme polynôme en $y$, vaut alors $$
\frac{\text{d}}{p\text{d}y}
(\lambda_0 + \lambda_1 y + \ldots + \lambda_{p - 1}y^{p - 1})^p
$$ Puisque $\frac{\text{d}}{\text{d}y}(y^{pk}) = pk y^{pk - 1}$, il suffit de déterminer les coefficients de $y^{pk}$ dans le polynôme $(\lambda_0 + \lambda_1 y + \ldots + \lambda_{p - 1}y^{p - 1})^p$ modulo $p$. La somme de ces termes donne $$
\lambda_0^p + \lambda_1^py^p + \ldots + \lambda_{p - 1}^py^{p(p - 1)}
\bmod p
$$ On voit donc qu'après application de l'opérateur $\frac{\text{d}}{p\text{d}y}$ et réduction modulo $y^p - b$, on obtient la valeur $$
\lambda_1^p + 2\lambda_2^pb + \ldots + (p - 1)\lambda_{p - 1}b^{p - 2}
$$ pour le coefficient de $y^{p - 1}$ que nous voulions calculer. Or, puisque $a = \lambda_0^p + \lambda_1^p b + \ldots + \lambda_{p - 1}^pb^{p - 1}$ dans $R$, on obtient $$
\text{d}a = (\lambda_1^p  + 2 \lambda_2^p b + \ldots +
(p - 1) \lambda_{p - 1}^p b^{p - 2}) \text{d}b
$$ dans $R$. Le coefficient de $y^{p - 1}$ est donc bien celui que l'on cherchait. \end{proof}

\begin{lemma} \label{lemma-trace-higher} Soient $\mathbf{F}_p \subset \Lambda \subset R \subset S$ des extensions d'anneaux ; supposons que $S$ soit isomorphe à $R[x]/(x^p - a)$ pour un certain $a \in R$. Il existe alors des applications canoniques $R$-linéaires $$
\text{Tr} :
\Omega^{t + 1}_{S/\Lambda}
\longrightarrow
\Omega_{R/\Lambda}^{t + 1}
$$ pour $t \geq 0$, telles que $$
\eta_1 \wedge \ldots \wedge \eta_t \wedge x^i\text{d}x
\longmapsto
\left\{
\begin{matrix}
0 & \text{si} & 0 \leq i \leq p - 2, \\
\eta_1 \wedge \ldots \wedge \eta_t \wedge \text{d}a & \text{si} & i = p - 1
\end{matrix}
\right.
$$ pour $\eta_i \in \Omega_{R/\Lambda}$ et telles que $\text{Tr}$ annule l'image de $S \otimes_R \Omega_{R/\Lambda}^{t + 1} \to \Omega_{S/\Lambda}^{t + 1}$. \end{lemma}

\begin{proof} Pour $t = 0$, on utilise le composé $$
\Omega_{S/\Lambda} \to \Omega_{S/R} \to \Omega_R \to \Omega_{R/\Lambda}
$$ où la deuxième application est celle du lemme \ref{lemma-trace-well-defined}. On a une suite exacte $$
H_1(L_{S/R}) \xrightarrow{\delta} \Omega_{R/\Lambda} \otimes_R S \to
\Omega_{S/\Lambda} \to \Omega_{S/R} \to 0
$$ (Algèbre, lemme \ref{algebra-lemma-exact-sequence-NL}). Le module $\Omega_{S/R}$ est libre sur $S$ de base $\text{d}x$, et le module $H_1(L_{S/R})$ est libre sur $S$ de base $x^p - a$, que $\delta$ envoie sur $-\text{d}a \otimes 1$ dans $\Omega_{R/\Lambda} \otimes_R S$. En particulier, si l'on pose $$
M = \Coker(R \to \Omega_{R/\Lambda}, 1 \mapsto -\text{d}a)
$$, on voit que $\Coker(\delta) = M \otimes_R S$. On obtient une application canonique $$
\Omega^{t + 1}_{S/\Lambda} \to
\wedge_S^t(\Coker(\delta)) \otimes_S \Omega_{S/R} =
\wedge^t_R(M) \otimes_R \Omega_{S/R}
$$ Or l'image de l'application $\text{Tr} : \Omega_{S/R} \to \Omega_{R/\Lambda}$ du lemme \ref{lemma-trace-well-defined} est contenue dans $R\text{d}a$ ; le produit extérieur par un élément de cette image annule donc $\text{d}a$. Il existe par conséquent une application canonique $$
\wedge^t_R(M) \otimes_R \Omega_{S/R} \to \Omega_{R/\Lambda}^{t + 1}
$$ envoyant $\overline{\eta}_1 \wedge \ldots \wedge \overline{\eta}_t \wedge \omega$ sur $\eta_1 \wedge \ldots \wedge \eta_t \wedge \text{Tr}(\omega)$. \end{proof}

\begin{remark} \label{remark-compare-Garel} Soient $\mathbf{F}_p \subset \Lambda \subset R \subset S$ et $\text{Tr}$ comme dans le lemme \ref{lemma-trace-higher}. D'après Cohomologie de de Rham, proposition \ref{derham-proposition-Garel}, il existe un morphisme canonique de complexes $$
\Theta_{S/R} :
\Omega_{S/\Lambda}^\bullet
\longrightarrow
\Omega_{R/\Lambda}^\bullet
$$ Le calcul de Cohomologie de de Rham, exemple \ref{derham-example-Garel}, montre que $\Theta_{S/R}(x^i \text{d}x) = \text{Tr}_x(x^i\text{d}x)$ pour tout $i$. Puisque $\text{Trace}_{S/R} = \Theta^0_{S/R}$ est identiquement nul et que $$
\Theta_{S/R}(a \wedge b) = a \wedge \Theta_{S/R}(b)
$$ pour $a \in \Omega^i_{R/\Lambda}$ et $b \in \Omega^j_{S/\Lambda}$, on en déduit $\text{Tr} = \Theta_{S/R}$. L'avantage de $\text{Tr}$ est que sa construction est nettement plus élémentaire. \end{remark}

\begin{lemma} \label{lemma-trace-extends} Soit $S$ un schéma sur $\mathbf{F}_p$. Soit $f : Y \to X$ un morphisme fini de schémas noethériens normaux intègres sur $S$. Supposons que \begin{enumerate} \item l'extension des corps de fonctions soit purement inséparable de degré $p$, et que \item $\Omega_{X/S}$ soit un $\mathcal{O}_X$-module cohérent (par exemple lorsque $X$ est de type fini sur $S$). \end{enumerate} Pour $i \geq 1$, il existe une application canonique $$
\text{Tr} : f_*\Omega^i_{Y/S} \longrightarrow (\Omega_{X/S}^i)^{**}
$$ dont la fibre au point générique de $X$ est l'application trace du lemme \ref{lemma-trace-higher}. \end{lemma}

\begin{proof} La suite exacte $f^*\Omega_{X/S} \to \Omega_{Y/S} \to \Omega_{Y/X} \to 0$ montre que $\Omega_{Y/S}$, et donc $f_*\Omega_{Y/S}$, sont également des modules cohérents. Il suffit donc de prouver que l'application trace au point générique se prolonge aux fibres en $x \in X$ tels que $\dim(\mathcal{O}_{X, x}) = 1$ ; voir Diviseurs, lemme \ref{divisors-lemma-describe-reflexive-hull}. On se ramène ainsi au cas traité dans le paragraphe suivant.

\medskip\noindent Supposons que $X = \Spec(A)$ et $Y = \Spec(B)$, où $A$ est un anneau de valuation discrète et $B$ est fini sur $A$. Puisque l'extension induite $L/K$ des corps des fractions est purement inséparable, $B$ est lui aussi local. Ainsi $B$ est également un anneau de valuation discrète. Alors, ou bien \begin{enumerate} \item $B/A$ a pour indice de ramification $p$, et donc $B = A[x]/(x^p - a)$, où $a \in A$ est une uniformisante ; ou bien \item $\mathfrak m_B = \mathfrak m_A B$ et le corps résiduel $B/\mathfrak m_A B$ est purement inséparable de degré $p$ sur $\kappa_A = A/\mathfrak m_A$. Choisissons un élément quelconque $x \in B$ dont la classe résiduelle n'appartienne pas à $\kappa_A$ ; on aura alors $B = A[x]/(x^p - a)$, où $a \in A$ est inversible. \end{enumerate} Soit $\Spec(\Lambda) \subset S$ un ouvert affine tel que $X$ ait son image dans $\Spec(\Lambda)$. On peut alors appliquer le lemme \ref{lemma-trace-higher} pour voir que l'application trace se prolonge en $\Omega^i_{B/\Lambda} \to \Omega^i_{A/\Lambda}$ pour tout $i \geq 1$. \end{proof}


















\section{Transformations quadratiques} \label{section-quadratic}

\noindent Dans cette section, nous étudions ce qui se produit lorsqu'on éclate un point non singulier d'une surface. Nous hésitons à définir formellement un tel morphisme comme une {\it transformation quadratique}, car, d'une part, on emploie souvent d'autres noms et, d'autre part, l'expression ``transformation quadratique'' est parfois utilisée dans un autre sens.

\begin{lemma} \label{lemma-blowup} Soit $(A, \mathfrak m, \kappa)$ un anneau local régulier de dimension $2$. Soit $f : X \to S = \Spec(A)$ l'éclatement de $A$ en $\mathfrak m$, de diviseur exceptionnel $E$. Il existe une immersion fermée $$
r : X \longrightarrow \mathbf{P}^1_S
$$ au-dessus de $S$ telle que \begin{enumerate} \item $r|_E : E \to \mathbf{P}^1_\kappa$ soit un isomorphisme, \item $\mathcal{O}_X(E) = \mathcal{O}_X(-1) =
r^*\mathcal{O}_{\mathbf{P}^1}(-1)$, et \item $\mathcal{C}_{E/X} = (r|_E)^*\mathcal{O}_{\mathbf{P}^1}(1)$ et $\mathcal{N}_{E/X} = (r|_E)^*\mathcal{O}_{\mathbf{P}^1}(-1)$. \end{enumerate} \end{lemma}

\begin{proof} Puisque $A$ est régulier de dimension $2$, on peut écrire $\mathfrak m = (x, y)$. Alors $x$ et $y$, placés en degré $1$, engendrent l'algèbre de Rees $\bigoplus_{n \geq 0} \mathfrak m^n$ sur $A$. Rappelons que $X = \text{Proj}(\bigoplus_{n \geq 0} \mathfrak m^n)$ ; voir Diviseurs, lemme \ref{divisors-lemma-blowing-up-affine}. La surjection $$
A[T_0, T_1] \longrightarrow \bigoplus\nolimits_{n \geq 0} \mathfrak m^n,
\quad
T_0 \mapsto x,\ T_1 \mapsto y
$$
\begingroup\emergencystretch=2em
de $A$-algèbres graduées induit donc une immersion fermée $r : X \to \mathbf{P}^1_S = \text{Proj}(A[T_0, T_1])$ telle que $\mathcal{O}_X(1) = r^*\mathcal{O}_{\mathbf{P}^1_S}(1)$ ; voir Constructions, lemme \ref{constructions-lemma-surjective-graded-rings-generated-degree-1-map-proj}. Cela prouve (2), car $\mathcal{O}_X(E) = \mathcal{O}_X(-1)$ d'après Diviseurs, lemme \ref{divisors-lemma-blowing-up-gives-effective-Cartier-divisor}.\par\endgroup

\medskip\noindent Pour prouver (1), remarquons que $$
\left(\bigoplus\nolimits_{n \geq 0} \mathfrak m^n\right) \otimes_A \kappa =
\bigoplus\nolimits_{n \geq 0} \mathfrak m^n/\mathfrak m^{n + 1} \cong
\kappa[\overline{x}, \overline{y}]
$$ est une algèbre de polynômes ; voir Algèbre, lemme \ref{algebra-lemma-regular-graded}. Cela prouve que la fibre de $X \to S$ au-dessus de $\Spec(\kappa)$ est égale à $\text{Proj}(\kappa[\overline{x}, \overline{y}]) = \mathbf{P}^1_\kappa$ ; voir Constructions, lemme \ref{constructions-lemma-base-change-map-proj}. Rappelons que $E$ est le sous-schéma fermé de $X$ défini par $\mathfrak m\mathcal{O}_X$, c'est-à-dire $E = X_\kappa$. Par le choix du morphisme $r$, on voit que $r|_E$ réalise en fait l'identification de $E = X_\kappa$ avec la fibre spéciale de $\mathbf{P}^1_S \to S$.

\medskip\noindent L'assertion (3) résulte de (1), de (2) et de Diviseurs, lemme \ref{divisors-lemma-conormal-effective-Cartier-divisor}. \end{proof}

\begin{lemma} \label{lemma-blowup-regular} Soit $(A, \mathfrak m, \kappa)$ un anneau local régulier de dimension $2$. Soit $f : X \to S = \Spec(A)$ l'éclatement de $A$ en $\mathfrak m$. Alors $X$ est un schéma régulier irréductible. \end{lemma}

\begin{proof} Remarquons que $X$ est intègre d'après Diviseurs, lemme \ref{divisors-lemma-blow-up-integral-scheme}, et Algèbre, lemme \ref{algebra-lemma-regular-domain}. Pour voir que $X$ est régulier, il suffit de vérifier que $\mathcal{O}_{X, x}$ est régulier aux points fermés $x \in X$ ; voir Propriétés, lemme \ref{properties-lemma-characterize-regular}. Soit $x \in X$ un point fermé. Puisque $f$ est propre, $x$ a pour image $\mathfrak m$ ; autrement dit, $x$ appartient au diviseur exceptionnel $E$. Alors $E$ est un diviseur de Cartier effectif et $E \cong \mathbf{P}^1_\kappa$. Ainsi, si $g \in \mathfrak m_x \subset \mathcal{O}_{X, x}$ est une équation locale de $E$, on a $\mathcal{O}_{X, x}/(g) \cong \mathcal{O}_{\mathbf{P}^1_\kappa, x}$. Puisque $\mathbf{P}^1_\kappa$ est recouvert par deux ouverts affines qui sont les spectres d'anneaux de polynômes sur $\kappa$, l'anneau $\mathcal{O}_{\mathbf{P}^1_\kappa, x}$ est régulier d'après Algèbre, lemme \ref{algebra-lemma-dim-affine-space}. On conclut par Algèbre, lemme \ref{algebra-lemma-regular-mod-x}. \end{proof}

\begin{lemma} \label{lemma-blowup-pic} Soit $(A, \mathfrak m, \kappa)$ un anneau local régulier de dimension $2$. Soit $f : X \to S = \Spec(A)$ l'éclatement de $A$ en $\mathfrak m$. Alors $\Pic(X) = \mathbf{Z}$, avec pour générateur $\mathcal{O}_X(E)$. \end{lemma}

\begin{proof} Rappelons que $E = \mathbf{P}^1_\kappa$ a pour groupe de Picard $\mathbf{Z}$, avec pour générateur $\mathcal{O}(1)$ ; voir Diviseurs, lemme \ref{divisors-lemma-Pic-projective-space-UFD}. D'après le lemme \ref{lemma-blowup}, le $\mathcal{O}_X$-module inversible $\mathcal{O}_X(E)$ se restreint en $\mathcal{O}(-1)$. Ainsi $\mathcal{O}_X(E)$ engendre un groupe cyclique infini dans $\Pic(X)$. Puisque $A$ est régulier, il est factoriel ; voir Compléments d'algèbre, lemme \ref{more-algebra-lemma-regular-local-UFD}. Le spectre épointé $U = S \setminus \{\mathfrak m\} = X \setminus E$ a donc un groupe de Picard trivial ; voir Diviseurs, lemme \ref{divisors-lemma-open-subscheme-UFD}. Par conséquent, pour tout $\mathcal{O}_X$-module inversible $\mathcal{L}$, il existe un isomorphisme $s : \mathcal{O}_U \to \mathcal{L}|_U$. Alors $s$ est une section méromorphe régulière de $\mathcal{L}$ et l'on voit que $\text{div}_\mathcal{L}(s) = nE$ pour un certain $n \in \mathbf{Z}$ (Diviseurs, définition \ref{divisors-definition-divisor-invertible-sheaf}). D'après Diviseurs, lemme \ref{divisors-lemma-normal-c1-injective} (et la normalité de $X$ donnée par le lemme \ref{lemma-blowup-regular}), on conclut que $\mathcal{L} = \mathcal{O}_X(nE)$. \end{proof}

\begin{lemma} \label{lemma-cohomology-of-blowup} Soit $(A, \mathfrak m, \kappa)$ un anneau local régulier de dimension $2$. Soit $f : X \to S = \Spec(A)$ l'éclatement de $A$ en $\mathfrak m$. Soit $\mathcal{F}$ un $\mathcal{O}_X$-module quasi-cohérent. \begin{enumerate} \item $H^p(X, \mathcal{F}) = 0$ pour $p \not \in \{0, 1\}$ ; \item $H^1(X, \mathcal{O}_X(n)) = 0$ pour $n \geq -1$ ; \item $H^1(X, \mathcal{F}) = 0$ si $\mathcal{F}$ ou $\mathcal{F}(1)$ est engendré par ses sections globales ; \item $H^0(X, \mathcal{O}_X(n)) = \mathfrak m^{\max(0, n)}$ ; \item $\text{length}_A H^1(X, \mathcal{O}_X(n)) = -n(-n - 1)/2$ si $n < 0$. \end{enumerate} \end{lemma}

\begin{proof} Si $\mathfrak m = (x, y)$, alors $X$ est recouvert par les spectres des algèbres d'éclatement affines $A[\frac{\mathfrak m}{x}]$ et $A[\frac{\mathfrak m}{y}]$, car $x$ et $y$, placés en degré $1$, engendrent l'algèbre de Rees $\bigoplus \mathfrak m^n$ sur $A$. Voir Diviseurs, lemme \ref{divisors-lemma-blowing-up-affine}, et Constructions, lemme \ref{constructions-lemma-proj-quasi-compact}. Puisque $X$ est séparé d'après Constructions, lemme \ref{constructions-lemma-proj-separated}, la cohomologie des faisceaux quasi-cohérents s'annule en degrés $\geq 2$ d'après Cohomologie des schémas, lemme \ref{coherent-lemma-vanishing-nr-affines}.

\medskip\noindent Soit $i : E \to X$ le diviseur exceptionnel ; voir Diviseurs, définition \ref{divisors-definition-blow-up}. Rappelons que $\mathcal{O}_X(-E) = \mathcal{O}_X(1)$ est relativement ample pour $f$ ; voir Diviseurs, lemme \ref{divisors-lemma-blowing-up-gives-effective-Cartier-divisor}. On sait donc que $H^1(X, \mathcal{O}_X(-nE)) = 0$ pour un certain $n > 0$ ; voir Cohomologie des schémas, lemme \ref{coherent-lemma-kill-by-twisting}. Considérons la filtration $$
\mathcal{O}_X(-nE) \subset \mathcal{O}_X(-(n - 1)E) \subset
\ldots \subset \mathcal{O}_X(-E) \subset \mathcal{O}_X \subset \mathcal{O}_X(E)
$$ Ses quotients successifs sont les faisceaux $$
\mathcal{O}_X(-t E)/\mathcal{O}_X(-(t + 1)E) =
\mathcal{O}_X(t)/\mathcal{I}(t) =
i_*\mathcal{O}_E(t)
$$ où $\mathcal{I} = \mathcal{O}_X(-E)$ est le faisceau d'idéaux de $E$. D'après le lemme \ref{lemma-blowup}, on a $E = \mathbf{P}^1_\kappa$, et $\mathcal{O}_E(1)$ correspond bien à la torsion de Serre usuelle du faisceau structural sur $\mathbf{P}^1$. La cohomologie de $\mathcal{O}_E(t)$ s'annule donc en degré $1$ pour $t \geq -1$ ; voir Cohomologie des schémas, lemme \ref{coherent-lemma-cohomology-projective-space-over-ring}. Comme elle est égale à $H^1(X, i_*\mathcal{O}_E(t))$ (d'après Cohomologie des schémas, lemme \ref{coherent-lemma-relative-affine-cohomology}), on en déduit que $H^1(X, \mathcal{O}_X(-(t + 1)E)) \to H^1(X, \mathcal{O}_X(-tE))$ est surjective pour $t \geq -1$. Ainsi $$
0 = H^1(X, \mathcal{O}_X(-nE))
\longrightarrow
H^1(X, \mathcal{O}_X(-tE)) = H^1(X, \mathcal{O}_X(t))
$$ est surjective pour $t \geq -1$, ce qui prouve (2).

\medskip\noindent Supposons $\mathcal{F}$ engendré par ses sections globales. Il existe donc une suite exacte courte $$
0 \to \mathcal{G} \to \bigoplus\nolimits_{i \in I} \mathcal{O}_X
\to \mathcal{F} \to 0
$$
\begingroup\emergencystretch=1em
Remarquons que $H^1(X, \bigoplus_{i \in I} \mathcal{O}_X) =
\bigoplus_{i \in I} H^1(X, \mathcal{O}_X)$ d'après Cohomologie, lemme \ref{cohomology-lemma-quasi-separated-cohomology-colimit}. Par (2), on a $H^1(X, \mathcal{O}_X) = 0$. Si $\mathcal{F}(1)$ est engendré par ses sections globales, on peut trouver une surjection $\bigoplus_{i \in I} \mathcal{O}_X(-1) \to \mathcal{F}$ et raisonner de même. Autrement dit, (3) résulte de (2).
\par\endgroup

\medskip\noindent Pour (4), remarquons que, pour tout $n$ assez grand, on a $\Gamma(X, \mathcal{O}_X(n)) = \mathfrak m^n$ ; voir Cohomologie des schémas, lemme \ref{coherent-lemma-recover-tail-graded-module}. Si $n \geq 0$, la suite exacte courte $$
0 \to \mathcal{O}_X(n) \to \mathcal{O}_X(n - 1) \to
i_*\mathcal{O}_E(n - 1) \to 0
$$ et l'annulation de $H^1$ pour le faisceau de gauche donnent un diagramme commutatif
\begingroup\small
$$
\xymatrix{
0 \ar[r] &
\mathfrak m^{\max(0, n)} \ar[r] \ar[d] &
\mathfrak m^{\max(0, n - 1)} \ar[r] \ar[d] &
\mathfrak m^{\max(0, n)}/\mathfrak m^{\max(0, n - 1)} \ar[r] \ar[d] & 0\\
0 \ar[r] &
\Gamma(X, \mathcal{O}_X(n)) \ar[r] &
\Gamma(X, \mathcal{O}_X(n - 1)) \ar[r] &
\Gamma(E, \mathcal{O}_E(n - 1)) \ar[r] & 0
}
$$
\endgroup
à lignes exactes. En fait, les lignes sont également exactes pour $n < 0$, car les groupes de droite sont alors nuls. Dans la démonstration du lemme \ref{lemma-blowup}, nous avons vu que la flèche verticale de droite est un isomorphisme (détails omis). Si celle de gauche est un isomorphisme, il en est donc de même de celle du milieu. On obtient ainsi (4) par récurrence descendante sur $n$.

\medskip\noindent Enfin, on prouve (5) par récurrence descendante sur $n$ à l'aide des suites $$
0 \to \mathcal{O}_X(n) \to \mathcal{O}_X(n - 1) \to
i_*\mathcal{O}_E(n - 1) \to 0
$$ En effet, pour $n \geq -1$, on sait déjà que $H^1(X, \mathcal{O}_X(n)) = 0$. Puisque $$
H^1(X, i_*\mathcal{O}_E(-2)) =
H^1(E, \mathcal{O}_E(-2)) =
H^1(\mathbf{P}^1_\kappa, \mathcal{O}(-2)) \cong \kappa
$$ d'après Cohomologie des schémas, lemme \ref{coherent-lemma-cohomology-projective-space-over-ring}, et que ce module est de longueur $1$ sur $A$, la suite exacte longue de cohomologie montre que (5) est vrai pour $n = -2$. On poursuit de même. \end{proof}

\begin{lemma} \label{lemma-blowup-improve} Soit $(A, \mathfrak m)$ un anneau local régulier de dimension $2$. Soit $f : X \to S = \Spec(A)$ l'éclatement de $A$ en $\mathfrak m$. Soit $\mathfrak m^n \subset I \subset \mathfrak m$ un idéal. Soit $d \geq 0$ le plus grand entier tel que $$
I \mathcal{O}_X \subset \mathcal{O}_X(-dE)
$$, où $E$ est le diviseur exceptionnel. Posons $\mathcal{I}' = I\mathcal{O}_X(dE) \subset \mathcal{O}_X$. Alors $d > 0$, le faisceau $\mathcal{O}_X/\mathcal{I}'$ a pour support un ensemble fini de points fermés $x_1, \ldots, x_r$ de $X$, et \begin{align*}
\text{length}_A(A/I)
& >
\text{length}_A \Gamma(X, \mathcal{O}_X/\mathcal{I}') \\
& \geq
\sum\nolimits_{i = 1, \ldots, r}
\text{length}_{\mathcal{O}_{X, x_i}}
(\mathcal{O}_{X, x_i}/\mathcal{I}'_{x_i})
\end{align*} \end{lemma}

\begin{proof} Puisque $I \subset \mathfrak m$, tout élément de $I$ s'annule sur $E$. On a donc $d \geq 1$. D'autre part, puisque $\mathfrak m^n \subset I$, on a $d \leq n$. Considérons la suite exacte courte $$
0 \to I\mathcal{O}_X \to \mathcal{O}_X \to \mathcal{O}_X/I\mathcal{O}_X \to 0
$$ Puisque $I\mathcal{O}_X$ est engendré par ses sections globales, on a $H^1(X, I\mathcal{O}_X) = 0$ d'après le lemme \ref{lemma-cohomology-of-blowup}. On obtient donc une surjection $A/I \to \Gamma(X, \mathcal{O}_X/I\mathcal{O}_X)$. Considérons la suite exacte courte $$
0 \to
\mathcal{O}_X(-dE)/I\mathcal{O}_X \to
\mathcal{O}_X/I\mathcal{O}_X \to
\mathcal{O}_X/\mathcal{O}_X(-dE) \to 0
$$ D'après Diviseurs, lemme \ref{divisors-lemma-codim-1-part}, le faisceau $\mathcal{O}_X(-dE)/I\mathcal{O}_X$ a pour support un ensemble fini de points fermés de $X$. En particulier, ses groupes de cohomologie supérieure sont nuls (détail omis). Dans le diagramme suivant
\begingroup\footnotesize
$$
\xymatrix{
& & A/I \ar[d] \\
0 \ar[r] &
\Gamma(X, \mathcal{O}_X(-dE)/I\mathcal{O}_X) \ar[r] &
\Gamma(X, \mathcal{O}_X/I\mathcal{O}_X) \ar[r] &
\Gamma(X, \mathcal{O}_X/\mathcal{O}_X(-dE)) \ar[r] & 0
}
$$\endgroup, la ligne inférieure est donc exacte et la flèche verticale surjective. On a $$
\text{length}_A \Gamma(X, \mathcal{O}_X(-dE)/I\mathcal{O}_X) <
\text{length}_A(A/I)
$$, car $\Gamma(X, \mathcal{O}_X/\mathcal{O}_X(-dE))$ est non nul. En effet, l'image de $1 \in \Gamma(X, \mathcal{O}_X)$ est non nulle puisque $d > 0$.

\medskip\noindent Pour achever la démonstration, traduisons les résultats ci-dessus en les assertions du lemme. Puisque $\mathcal{O}_X(dE)$ est inversible, on a $$
\mathcal{O}_X/\mathcal{I}' =
\mathcal{O}_X(-dE)/I\mathcal{O}_X \otimes_{\mathcal{O}_X} \mathcal{O}_X(dE).
$$ Ainsi $\mathcal{O}_X/\mathcal{I}'$ et $\mathcal{O}_X(-dE)/I\mathcal{O}_X$ ont pour support le même ensemble fini de points fermés, notés $x_1, \ldots, x_r \in E \subset X$. De plus, on obtient $$
\Gamma(X, \mathcal{O}_X(-dE)/I\mathcal{O}_X) =
\bigoplus \mathcal{O}_X(-dE)_{x_i}/I\mathcal{O}_{X, x_i}
\cong
\bigoplus \mathcal{O}_{X, x_i}/\mathcal{I}'_{x_i} =
\Gamma(X, \mathcal{O}_X/\mathcal{I}')
$$, car un module inversible sur un anneau local est trivial. On en déduit l'inégalité stricte. On obtient également la seconde inégalité, car $$
\text{length}_A(\mathcal{O}_{X, x_i}/\mathcal{I}'_{x_i}) \geq
\text{length}_{\mathcal{O}_{X, x_i}}(\mathcal{O}_{X, x_i}/\mathcal{I}'_{x_i})
$$, comme il résulte immédiatement de la définition de la longueur. \end{proof}

\begin{lemma} \label{lemma-differentials-of-blowup} Soit $(A, \mathfrak m, \kappa)$ un anneau local régulier de dimension $2$. Soit $f : X \to S = \Spec(A)$ l'éclatement de $A$ en $\mathfrak m$. Alors $\Omega_{X/S} = i_*\Omega_{E/\kappa}$, où $i : E \to X$ est l'immersion du diviseur exceptionnel. \end{lemma}

\begin{proof} Écrivons $\mathbf{P}^1 = \mathbf{P}^1_S$ et soit $r : X \to \mathbf{P}^1$ comme dans le lemme \ref{lemma-blowup}. On a alors une suite exacte $$
\mathcal{C}_{X/\mathbf{P}^1} \to r^*\Omega_{\mathbf{P}^1/S} \to
\Omega_{X/S} \to 0
$$ ; voir Morphismes, lemme \ref{morphisms-lemma-differentials-relative-immersion}. Puisque $\Omega_{\mathbf{P}^1/S}|_E = \Omega_{E/\kappa}$ d'après Morphismes, lemme \ref{morphisms-lemma-base-change-differentials}, il suffit de voir que la première flèche définit une surjection sur le noyau de l'application canonique $r^*\Omega_{\mathbf{P}^1/S} \to i_*\Omega_{E/\kappa}$. Cela se vérifie localement. Avec les notations de la démonstration du lemme \ref{lemma-blowup}, sur un ouvert affine de $X$, le morphisme $f$ correspond à l'homomorphisme d'anneaux $$
A \to A[t]/(xt - y)
$$, où $x, y \in \mathfrak m$ sont des générateurs. Ainsi $\text{d}(xt - y) = x\text{d}t$ et $y\text{d}t = t \cdot x \text{d}t$, ce qui prouve l'assertion. \end{proof}



\section{Domination par des transformations quadratiques} \label{section-dominating-by-quadratic}

\noindent Le résultat précédent permet de prouver que les éclatements de points dominent toute modification d'un schéma régulier de dimension $2$.

\medskip\noindent Soit $X$ un schéma. Soit $x \in X$ un point fermé. Comme d'habitude, on considère $i : x = \Spec(\kappa(x)) \to X$ comme un sous-schéma fermé. L'{\it éclatement $X' \to X$ de $X$ au point $x$} est l'éclatement de $X$ suivant le sous-schéma fermé $x \subset X$. Remarquons que, si $X$ est localement noethérien, alors $X' \to X$ est projectif (donc propre) d'après Diviseurs, lemme \ref{divisors-lemma-blowing-up-projective}.

\begin{lemma} \label{lemma-make-ideal-principal} Soit $X$ un schéma noethérien. Soit $T \subset X$ un ensemble fini de points fermés $x$ tels que $\mathcal{O}_{X, x}$ soit régulier de dimension $2$ pour $x \in T$. Soit $\mathcal{I} \subset \mathcal{O}_X$ un faisceau quasi-cohérent d'idéaux tel que $\mathcal{O}_X/\mathcal{I}$ ait son support dans $T$. Il existe alors une suite $$
X_n \to X_{n - 1} \to \ldots \to X_1 \to X_0 = X
$$, où $X_{i + 1} \to X_i$ est l'éclatement de $X_i$ en un point fermé situé au-dessus d'un point de $T$, telle que $\mathcal{I}\mathcal{O}_{X_n}$ soit un faisceau d'idéaux inversible. \end{lemma}

\begin{proof} Écrivons $T = \{x_1, \ldots, x_r\}$. Notons $I_i$ la fibre de $\mathcal{I}$ en $x_i$. Posons $$
n_i = \text{length}_{\mathcal{O}_{X, x_i}}(\mathcal{O}_{X, x_i}/I_i)
$$ Cet entier est fini, car $\mathcal{O}_X/\mathcal{I}$ a son support dans $T$ et le support de $\mathcal{O}_{X, x_i}/I_i$ est donc égal à $\{\mathfrak m_{x_i}\}$ (voir Algèbre, lemme \ref{algebra-lemma-support-point}). Raisonnons par récurrence sur $\sum n_i$. Si $n_i = 0$ pour tout $i$, alors $\mathcal{I} = \mathcal{O}_X$ et la démonstration est terminée.

\medskip\noindent Supposons que $n_i > 0$. Soit $X' \to X$ l'éclatement de $X$ en $x_i$ (voir la discussion précédant le lemme). Puisque $\Spec(\mathcal{O}_{X, x_i}) \to X$ est plat, $X' \times_X \Spec(\mathcal{O}_{X, x_i})$ est l'éclatement de l'anneau $\mathcal{O}_{X, x_i}$ suivant son idéal maximal ; voir Diviseurs, lemme \ref{divisors-lemma-flat-base-change-blowing-up}. Le carré du diagramme commutatif $$
\xymatrix{
\text{Proj}(\bigoplus\nolimits_{d \geq 0} \mathfrak m_{x_i}^d) \ar[r] \ar[d] &
X' \ar[d] \\
\Spec(\mathcal{O}_{X, x_i}) \ar[r] & X
}
$$ est donc cartésien. Soient $E \subset X'$ et $E' \subset \text{Proj}(\bigoplus\nolimits_{d \geq 0} \mathfrak m_{x_i}^d)$ les diviseurs exceptionnels. Soit $d \geq 1$ l'entier donné par le lemme \ref{lemma-blowup-improve} pour l'idéal $\mathcal{I}_i \subset \mathcal{O}_{X, x_i}$. Puisque les flèches horizontales du diagramme sont plates, que $E' \to E$ est surjectif et que $E'$ est l'image réciproque de $E$, on voit que $$
\mathcal{I}\mathcal{O}_{X'} \subset \mathcal{O}_{X'}(-dE)
$$ (certains détails sont omis). Posons $\mathcal{I}' = \mathcal{I}\mathcal{O}_{X'}(dE) \subset \mathcal{O}_{X'}$. Alors $\mathcal{O}_{X'}/\mathcal{I}'$ a pour support un ensemble fini de points fermés $T' \subset |X'|$, car c'est le cas au-dessus de $X \setminus \{x_i\}$ et après image réciproque sur $\text{Proj}(\bigoplus\nolimits_{d \geq 0} \mathfrak m_{x_i}^d)$. La dernière assertion du lemme \ref{lemma-blowup-improve} indique que la somme des longueurs des fibres $\mathcal{O}_{X', x'}/\mathcal{I}'\mathcal{O}_{X', x'}$, pour $x'$ situé au-dessus de $x_i$, est $< n_i$. La somme des longueurs a donc diminué.

\medskip\noindent Par l'hypothèse de récurrence, il existe une suite $$
X'_n \to \ldots \to X'_1 \to X'
$$ d'éclatements en des points fermés situés au-dessus de $T'$ telle que $\mathcal{I}'\mathcal{O}_{X'_n}$ soit inversible. Puisque $\mathcal{I}'\mathcal{O}_{X'}(-dE) = \mathcal{I}\mathcal{O}_{X'}$, on voit que $\mathcal{I}\mathcal{O}_{X'_n} =
\mathcal{I}'\mathcal{O}_{X'_n}(-d(f')^{-1}E)$, où $f' : X'_n \to X'$ est le composé. Remarquons que $(f')^{-1}E$ est un diviseur de Cartier effectif d'après Diviseurs, lemme \ref{divisors-lemma-blow-up-pullback-effective-Cartier}. On conclut donc par Diviseurs, lemme \ref{divisors-lemma-sum-effective-Cartier-divisors}. \end{proof}

\begin{lemma} \label{lemma-dominate-by-blowing-up-in-points} Soit $X$ un schéma noethérien. Soit $T \subset X$ un ensemble fini de points fermés $x$ tels que $\mathcal{O}_{X, x}$ soit un anneau local régulier de dimension $2$. Soit $f : Y \to X$ un morphisme propre de schémas qui soit un isomorphisme au-dessus de $U = X \setminus T$. Il existe alors une suite $$
X_n \to X_{n - 1} \to \ldots \to X_1 \to X_0 = X
$$, où $X_{i + 1} \to X_i$ est l'éclatement de $X_i$ en un point fermé $x_i$ situé au-dessus d'un point de $T$, ainsi qu'une factorisation $X_n \to Y \to X$ du composé. \end{lemma}

\begin{proof} D'après Compléments de platitude, lemme \ref{flat-lemma-dominate-modification-by-blowup}, il existe un éclatement $U$-admissible $X' \to X$ qui domine $Y \to X$. On peut donc supposer qu'il existe un faisceau d'idéaux $\mathcal{I} \subset \mathcal{O}_X$ tel que $\mathcal{O}_X/\mathcal{I}$ ait son support dans $T$ et que $Y$ soit l'éclatement de $X$ suivant $\mathcal{I}$. D'après le lemme \ref{lemma-make-ideal-principal}, il existe une suite $$
X_n \to X_{n - 1} \to \ldots \to X_1 \to X_0 = X
$$, où $X_{i + 1} \to X_i$ est l'éclatement de $X_i$ en un point fermé $x_i$ situé au-dessus d'un point de $T$, telle que $\mathcal{I}\mathcal{O}_{X_n}$ soit un faisceau d'idéaux inversible. La propriété universelle de l'éclatement (Diviseurs, lemme \ref{divisors-lemma-universal-property-blowing-up}) fournit la factorisation cherchée. \end{proof}

\begin{lemma} \label{lemma-extend-rational-map-blowing-up} Soit $S$ un schéma. Soit $X$ un schéma sur $S$, régulier et de dimension $2$. Soit $Y$ un schéma propre sur $S$. Étant donnée une application $S$-rationnelle $f : U \to Y$ de $X$ vers $Y$, il existe une suite $$
X_n \to X_{n - 1} \to \ldots \to X_1 \to X_0 = X
$$ et un $S$-morphisme $f_n : X_n \to Y$ tels que $X_{i + 1} \to X_i$ soit l'éclatement de $X_i$ en un point fermé non situé au-dessus de $U$ et que $f_n$ et $f$ coïncident. \end{lemma}

\begin{proof} On peut supposer que $U$ contient tous les points de codimension $1$ ; voir Morphismes, lemme \ref{morphisms-lemma-extend-across}. Le complémentaire $T \subset X$ de $U$ est donc un ensemble fini de points fermés dont les anneaux locaux sont réguliers de dimension $2$. En appliquant Diviseurs, lemme \ref{divisors-lemma-extend-rational-map-after-modification}, on trouve un morphisme propre $p : X' \to X$ qui est un isomorphisme au-dessus de $U$ et un morphisme $f' : X' \to Y$ coïncidant avec $f$ au-dessus de $U$. Appliquons le lemme \ref{lemma-dominate-by-blowing-up-in-points} au morphisme $p : X' \to X$. Le composé $X_n \to X' \to Y$ est le morphisme cherché. \end{proof}






\section{Domination par des éclatements normalisés} \label{section-normalized-blowups}

\noindent Dans cette section, nous prouvons qu'une modification d'une surface peut être dominée par une suite d'éclatements normalisés de points.

\begin{definition} \label{definition-normalized-blowup} Soit $X$ un schéma dont tout ouvert quasi-compact possède un nombre fini de composantes irréductibles. Soit $x \in X$ un point fermé. L'{\it éclatement normalisé de $X$ en $x$} est le composé $X'' \to X' \to X$, où $X' \to X$ est l'éclatement de $X$ en $x$ et $X'' \to X'$ est la normalisation de $X'$. \end{definition}

\noindent La normalisation $X'' \to X'$ est ici définie, car le schéma $X'$ possède un recouvrement par des ouverts ayant un nombre fini de composantes irréductibles d'après Diviseurs, lemme \ref{divisors-lemma-blow-up-and-irreducible-components}. Voir Morphismes, définition \ref{morphisms-definition-normalization}, pour la définition de la normalisation.

\medskip\noindent En général, l'éclatement normalisé n'est pas nécessairement propre, même si $X$ est noethérien. Rappelons qu'un schéma est de Nagata s'il possède un recouvrement par des ouverts affines qui sont des spectres d'anneaux de Nagata (Propriétés, définition \ref{properties-definition-nagata}).

\begin{lemma} \label{lemma-Nagata-normalized-blowup} Dans la définition \ref{definition-normalized-blowup}, si $X$ est de Nagata, alors l'éclatement normalisé de $X$ en $x$ est normal, de Nagata et propre sur $X$. \end{lemma}

\begin{proof} Le morphisme d'éclatement $X' \to X$ est propre (puisque $X$ est localement noethérien, on peut appliquer Diviseurs, lemme \ref{divisors-lemma-blowing-up-projective}). Ainsi $X'$ est de Nagata (Morphismes, lemme \ref{morphisms-lemma-finite-type-nagata}). La normalisation $X'' \to X'$ est donc finie (Morphismes, lemme \ref{morphisms-lemma-nagata-normalization}), d'où la propreté de $X'' \to X$ (Morphismes, lemmes \ref{morphisms-lemma-finite-proper} et \ref{morphisms-lemma-composition-proper}). Il en résulte que l'éclatement normalisé est un espace algébrique normal (Morphismes, lemme \ref{morphisms-lemma-normalization-normal}) et de Nagata. \end{proof}

\noindent Dans le lemme suivant, il faut supposer que $X$ est noethérien pour assurer qu'il possède un nombre fini de composantes irréductibles. La propreté de $f : Y \to X$ assure alors que $Y$ possède lui aussi un nombre fini de composantes irréductibles ; on peut donc demander que $f$ soit birationnel (Morphismes, définition \ref{morphisms-definition-birational}).

\begin{lemma} \label{lemma-dominate-by-normalized-blowing-up} Soit $X$ un schéma noethérien de Nagata de dimension $2$. Soit $f : Y \to X$ un morphisme propre birationnel. Il existe alors un diagramme commutatif $$
\xymatrix{
X_n \ar[r] \ar[d] &
X_{n - 1} \ar[r] &
\ldots \ar[r] &
X_1 \ar[r] &
X_0 \ar[d] \\
Y \ar[rrrr]  & & & & X
}
$$, où $X_0 \to X$ est la normalisation et où $X_{i + 1} \to X_i$ est l'éclatement normalisé de $X_i$ en un point fermé. \end{lemma}

\begin{proof} Nous utiliserons sans autre mention les résultats de Morphismes, sections \ref{morphisms-section-nagata}, \ref{morphisms-section-dimension-formula} et \ref{morphisms-section-normalization}. On peut remplacer $Y$ par sa normalisation. Soit $X_0 \to X$ la normalisation. Le morphisme $Y \to X$ se factorise par $X_0$. On peut donc supposer que $X$ et $Y$ sont tous deux normaux.

\medskip\noindent Supposons $X$ et $Y$ normaux. Le morphisme $f : Y \to X$ est un isomorphisme au-dessus d'un ouvert contenant tous les points de codimension $0$ et $1$ de $Y$, ainsi que tout point de $Y$ au-dessus duquel la fibre est finie ; voir Variétés, lemme \ref{varieties-lemma-modification-normal-iso-over-codimension-1}. Il existe donc un ensemble fini de points fermés $T \subset X$ tel que $f$ soit un isomorphisme au-dessus de $X \setminus T$. Pour tout $x \in T$, la fibre $Y_x$ est un schéma propre géométriquement connexe de dimension $1$ sur $\kappa(x)$ ; voir Compléments sur les morphismes, lemme \ref{more-morphisms-lemma-geometrically-connected-fibres-towards-normal}. Ainsi $$
BadCurves(f) = \{C \subset Y\text{ fermé} \mid
\dim(C) = 1, f(C) = \text{un point}\}
$$ est un ensemble fini. Nous prouverons le lemme par récurrence sur le nombre d'éléments de $BadCurves(f)$. Le cas initial est celui où $BadCurves(f)$ est vide ; alors $f$ est un isomorphisme.

\medskip\noindent Fixons $x \in T$. Soit $X' \to X$ l'éclatement normalisé de $X$ en $x$ et soit $Y'$ la normalisation de $Y \times_X X'$. On a le diagramme $$
\xymatrix{
Y' \ar[r]_{f'} \ar[d] & X' \ar[d] \\
Y \ar[r]^f & X
}
$$ Soit $x' \in X'$ un point fermé situé au-dessus de $x$ tel que la fibre $Y'_{x'}$ soit de dimension $\geq 1$. Soit $C' \subset Y'$ une composante irréductible de $Y'_{x'}$, c'est-à-dire $C' \in BadCurves(f')$. Puisque $Y' \to Y \times_X X'$ est fini, $C'$ a nécessairement pour image une composante irréductible $C \subset Y_x$. Il est clair que $C \in BadCurves(f)$. Puisque $Y' \to Y$ est birationnel, donc un isomorphisme au-dessus des points de codimension $1$ de $Y$, on obtient une application injective $$
BadCurves(f') \longrightarrow BadCurves(f)
$$ Il suffit donc de montrer qu'après un nombre fini de ces éclatements normalisés, on élimine au moins une des courbes exceptionnelles considérées, c'est-à-dire que l'application affichée n'est pas surjective.

\medskip\noindent Nous éliminerons une telle courbe au moyen d'un argument dû à Zariski. Choisissons $C \in BadCurves(f)$ situé au-dessus du point $x$ fixé. Notons $\mathcal{O}_{Y, C}$ l'anneau local de $Y$ au point générique de $C$. Choisissons un élément $u \in \mathcal{O}_{X, C}$ dont l'image dans le corps résiduel $R(C)$ soit transcendante sur $\kappa(x)$ (cela est possible, car $R(C)$ est de degré de transcendance $1$ sur $\kappa(x)$ d'après Variétés, lemme \ref{varieties-lemma-dimension-locally-algebraic}). On peut écrire $u = a/b$ avec $a, b \in \mathcal{O}_{X, x}$, car $\mathcal{O}_{Y, C}$ et $\mathcal{O}_{X, x}$ ont le même corps des fractions. Le choix de $u$ impose $a, b \in \mathfrak m_x$. Ainsi $$
N_{u, a, b} = \min
\{\text{ord}_{\mathcal{O}_{Y, C}}(a), \text{ord}_{\mathcal{O}_{Y, C}}(b)\} > 0
$$ On peut donc raisonner par récurrence descendante sur cet entier. Soit $X' \to X$ l'éclatement normalisé de $x$ et soit $Y'$ la normalisation de $X' \times_X Y$, comme ci-dessus. Nous montrerons que, si $C$ est l'image d'une courbe exceptionnelle $C' \subset Y'$ située au-dessus de $x' \in X'$, il existe un choix de $a', b' \mathcal{O}_{X', x'}$ tel que $N_{u, a', b'} < N_{u, a, b}$. Cela achèvera la démonstration. En effet, puisque $X' \to X$ se factorise par l'éclatement, il existe un élément non nul $d \in \mathfrak m_{x'}$ tel que $a = a' d$ et $b = b' d$ (prendre pour $d$ une équation locale du diviseur exceptionnel de l'éclatement). Puisque $Y' \to Y$ est un isomorphisme au-dessus d'un ouvert contenant le point générique de $C$ (comme on l'a vu), on a $\mathcal{O}_{Y', C'} = \mathcal{O}_{Y, C}$. Donc $$
\text{ord}_{\mathcal{O}_{Y, C}}(a) =
\text{ord}_{\mathcal{O}_{Y', C'}}(a' d) =
\text{ord}_{\mathcal{O}_{Y', C'}}(a') +
\text{ord}_{\mathcal{O}_{Y', C'}}(d) >
\text{ord}_{\mathcal{O}_{Y', C'}}(a')
$$ On raisonne de même pour $b$, et la démonstration est terminée. \end{proof}

\begin{lemma} \label{lemma-extend-rational-map-normalized-blowing-up} Soit $S$ un schéma. Soit $X$ un schéma sur $S$, noethérien, de Nagata et de dimension $2$. Soit $Y$ un schéma propre sur $S$. Étant donnée une application $S$-rationnelle $f : U \to Y$ de $X$ vers $Y$, il existe une suite $$
X_n \to X_{n - 1} \to \ldots \to X_1 \to X_0 \to X
$$ et un $S$-morphisme $f_n : X_n \to Y$ tels que $X_0 \to X$ soit la normalisation, que $X_{i + 1} \to X_i$ soit l'éclatement normalisé de $X_i$ en un point fermé et que $f_n$ et $f$ coïncident. \end{lemma}

\begin{proof} En appliquant Diviseurs, lemme \ref{divisors-lemma-extend-rational-map-after-modification}, on trouve un morphisme propre $p : X' \to X$ qui est un isomorphisme au-dessus de $U$ et un morphisme $f' : X' \to Y$ coïncidant avec $f$ au-dessus de $U$. Appliquons le lemme \ref{lemma-dominate-by-normalized-blowing-up} au morphisme $p : X' \to X$. Le composé $X_n \to X' \to Y$ est le morphisme cherché. \end{proof}





\section{Modifications au-dessus d'anneaux locaux} \label{section-modifications}

\noindent Soit $S$ un schéma. Soient $s_1, \ldots, s_n \in S$ des points fermés deux à deux distincts. Supposons que l'immersion ouverte $$
U = S \setminus \{s_1, \ldots, s_n\} \longrightarrow S
$$ soit quasi-compacte. Notons $FP_{S, \{s_1, \ldots, s_n\}}$ la catégorie des morphismes $f : X \to S$ de présentation finie qui induisent un isomorphisme $f^{-1}(U) \to U$. Les morphismes sont les morphismes de schémas au-dessus de $S$. Pour tout $i$, posons $S_i = \Spec(\mathcal{O}_{S, s_i})$ et soit $V_i = S_i \setminus \{s_i\}$. Notons $FP_{S_i, s_i}$ la catégorie des morphismes $g_i : Y_i \to S_i$ de présentation finie qui induisent un isomorphisme $g_i^{-1}(V_i) \to V_i$. Les morphismes sont les morphismes au-dessus de $S_i$. Le changement de base définit un foncteur \begin{equation}
\label{equation-equivalence}
F :
FP_{S, \{s_1, \ldots, s_n\}}
\longrightarrow
FP_{S_1, s_1} \times \ldots \times FP_{S_n, s_n}
\end{equation} Pour ramener au moins certains des problèmes de ce chapitre au cas des anneaux locaux, on dispose du lemme suivant.

\begin{lemma} \label{lemma-equivalence} Le foncteur $F$ (\ref{equation-equivalence}) est une équivalence. \end{lemma}

\begin{proof} Pour $n = 1$, c'est Limites, lemme \ref{limits-lemma-modifications}. Pour $n > 1$, on peut démontrer le lemme exactement de la même manière, ou le déduire de ce cas. Par exemple, supposons que $g_i : Y_i \to S_i$ soient des objets de $FP_{S_i, s_i}$. Le cas $n = 1$ fournit des morphismes $f'_i : X'_i \to S$ de présentation finie, isomorphes au-dessus de $S \setminus \{s_i\}$, dont le changement de base à $S_i$ est $g_i$. On peut alors poser $$
f : X = X'_1 \times_S \ldots \times_S X'_n \to S
$$ C'est un objet de $FP_{S, \{s_1, \ldots, s_n\}}$ dont le changement de base par $S_i \to S$ redonne $g_i$. Le foncteur est donc essentiellement surjectif. Nous omettons la preuve de sa pleine fidélité. \end{proof}

\begin{lemma} \label{lemma-equivalence-properties} Soit $S, s_i, S_i$ comme dans (\ref{equation-equivalence}). Si $f : X \to S$ correspond à $g_i : Y_i \to S_i$ par $F$, alors $f$ est séparé, propre, fini, si et seulement si $g_i$ l'est pour $i = 1, \ldots, n$. \end{lemma}

\begin{proof} Cela résulte de Limites, lemme \ref{limits-lemma-modifications-properties}. \end{proof}

\begin{lemma} \label{lemma-equivalence-fibre} Soit $S, s_i, S_i$ comme dans (\ref{equation-equivalence}). Si $f : X \to S$ correspond à $g_i : Y_i \to S_i$ par $F$, alors $X_{s_i} \cong (Y_i)_{s_i}$ en tant que schémas sur $\kappa(s_i)$. \end{lemma}

\begin{proof} C'est clair. \end{proof}

\begin{lemma} \label{lemma-equivalence-sequence-blowups} Soit $S, s_i, S_i$ comme dans (\ref{equation-equivalence}) ; supposons que $f : X \to S$ corresponde à $g_i : Y_i \to S_i$ par $F$. Il existe alors une factorisation $$
X = Z_m \to Z_{m - 1} \to \ldots \to Z_1 \to Z_0 = S
$$ de $f$, où $Z_{j + 1} \to Z_j$ est l'éclatement de $Z_j$ en un point fermé $z_j$ situé au-dessus de $\{s_1, \ldots, s_n\}$, si et seulement si, pour tout $i$, il existe une factorisation $$
Y_i = Z_{i, m_i} \to Z_{i, m_i - 1} \to \ldots \to Z_{i, 1} \to Z_{i, 0} = S_i
$$ de $g_i$, où $Z_{i, j + 1} \to Z_{i, j}$ est l'éclatement de $Z_{i, j}$ en un point fermé $z_{i, j}$ situé au-dessus de $s_i$. \end{lemma}

\begin{proof} Partons d'une suite d'éclatements $Z_m \to Z_{m - 1} \to \ldots \to Z_1 \to Z_0 = S$. Le premier morphisme $Z_1 \to S$ est obtenu en éclatant l'un des $s_i$, disons $s_1$. En appliquant $F$ à $Z_1 \to S$, on obtient l'éclatement $Z_{1, 1} \to S_1$ en $s_1$ ; il s'agit bien de l'éclatement en $s_1$, tandis que $Z_{i, 0} = S_i$ pour $i > 1$. À l'étape suivante, soit on éclate l'un des $s_i$, $i \geq 2$, sur $Z_1$, soit on choisit un point fermé $z_1$ de la fibre de $Z_1 \to S$ au-dessus de $s_1$. Dans le premier cas, la construction est claire ; dans le second, on utilise $(Z_1)_{s_1} \cong (Z_{1, 1})_{s_1}$ (lemme \ref{lemma-equivalence-fibre}) pour obtenir un point fermé $z_{1, 1} \in Z_{1, 1}$ correspondant à $z_1$. On prend alors pour $Z_{1, 2} \to Z_{1, 1}$ l'éclatement en $z_{1, 1}$. En poursuivant ainsi, on construit les factorisations de chaque $g_i$.

\medskip\noindent Réciproquement, étant données des suites d'éclatements $Z_{i, m_i} \to Z_{i, m_i - 1} \to \ldots \to Z_{i, 1} \to Z_{i, 0} = S_i$, on construit la suite d'éclatements de $S$ exactement de la même manière. \end{proof}

\noindent Voici l'analogue du lemme \ref{lemma-equivalence-sequence-blowups} pour les éclatements normalisés.

\begin{lemma} \label{lemma-equivalence-sequence-normalized-blowups} Soit $S, s_i, S_i$ comme dans (\ref{equation-equivalence}) ; supposons que $f : X \to S$ corresponde à $g_i : Y_i \to S_i$ par $F$. Supposons que tout ouvert quasi-compact de $S$ possède un nombre fini de composantes irréductibles. Il existe alors une factorisation $$
X = Z_m \to Z_{m - 1} \to \ldots \to Z_1 \to Z_0 = S
$$ de $f$, où $Z_{j + 1} \to Z_j$ est l'éclatement normalisé de $Z_j$ en un point fermé $z_j$ situé au-dessus de $\{x_1, \ldots, x_n\}$, si et seulement si, pour tout $i$, il existe une factorisation $$
Y_i = Z_{i, m_i} \to Z_{i, m_i - 1} \to \ldots \to Z_{i, 1} \to Z_{i, 0} = S_i
$$ de $g_i$, où $Z_{i, j + 1} \to Z_{i, j}$ est l'éclatement normalisé de $Z_{i, j}$ en un point fermé $z_{i, j}$ situé au-dessus de $s_i$. \end{lemma}

\begin{proof} L'hypothèse sur $S$ sert à assurer, de proche en proche, que les schémas que l'on normalise possèdent localement un nombre fini de composantes irréductibles, de sorte que l'énoncé a un sens. Cela dit, le lemme se déduit exactement du même argument que celui utilisé pour démontrer le lemme \ref{lemma-equivalence-sequence-blowups}. \end{proof}








\section{Annulation} \label{section-vanishing}

\noindent Dans cette section, nous travaillerons souvent dans le cadre suivant. Rappelons qu'une modification est un morphisme propre birationnel entre schémas intègres (Morphismes, définition \ref{morphisms-definition-modification}).

\begin{situation} \label{situation-vanishing} Ici, $(A, \mathfrak m, \kappa)$ est un anneau local noethérien normal intègre de dimension $2$. Soit $s$ le point fermé de $S = \Spec(A)$ et soit $U = S \setminus \{s\}$. Soit $f : X \to S$ une modification. Nous notons $C_1, \ldots, C_r$ les composantes irréductibles de la fibre spéciale $X_s$ de $f$. \end{situation}

\noindent D'après Variétés, lemme \ref{varieties-lemma-modification-normal-iso-over-codimension-1}, le morphisme $f$ définit un isomorphisme $f^{-1}(U) \to U$. La fibre spéciale $X_s$ est propre sur $\Spec(\kappa)$ et de dimension au plus $1$ d'après Variétés, lemme \ref{varieties-lemma-dimension-fibre-in-higher-codimension}. Par factorisation de Stein (Compléments sur les morphismes, lemme \ref{more-morphisms-lemma-geometrically-connected-fibres-towards-normal}), on a $f_*\mathcal{O}_X = \mathcal{O}_S$ et la fibre spéciale $X_s$ est géométriquement connexe sur $\kappa$. Si $X_s$ est de dimension $0$, alors $f$ est fini (Compléments sur les morphismes, lemme \ref{more-morphisms-lemma-proper-finite-fibre-finite-in-neighbourhood}), donc un isomorphisme (Morphismes, lemme \ref{morphisms-lemma-finite-birational-over-normal}). Nous écartons ce cas sans intérêt et concluons que $\dim(C_i) = 1$ pour $i = 1, \ldots, r$.

\begin{lemma} \label{lemma-dominate-by-scheme-modification} Dans la situation \ref{situation-vanishing}, il existe un éclatement $U$-admissible $X' \to S$ qui domine $X$. \end{lemma}

\begin{proof} C'est un cas particulier de Compléments sur la platitude, lemme \ref{flat-lemma-dominate-modification-by-blowup}. \end{proof}

\begin{lemma} \label{lemma-nice-meromorphic-function} Dans la situation \ref{situation-vanishing}, il existe un élément non nul $f \in \mathfrak m$ tel que, pour tout $i = 1, \ldots, r$, il existe \begin{enumerate} \item un point fermé $x_i \in C_i$ vérifiant $x_i \not \in C_j$ pour $j \not = i$, \item une factorisation $f = g_i f_i$ de $f$ dans $\mathcal{O}_{X, x_i}$ telle que $g_i \in \mathfrak m_{x_i}$ ait une image non nulle dans $\mathcal{O}_{C_i, x_i}$. \end{enumerate} \end{lemma}

\begin{proof} Nous utiliserons sans autre mention les observations faites après la situation \ref{situation-vanishing}. Choisissons un point fermé $x_i \in C_i$ qui n'appartienne à aucun $C_j$ pour $j \not = i$. Choisissons $g_i \in \mathfrak m_{x_i}$ dont l'image dans $\mathcal{O}_{C_i, x_i}$ soit non nulle. Puisque le corps des fractions de $A$ est celui de $\mathcal{O}_{X_i, x_i}$, on peut écrire $g_i = a_i/b_i$ pour certains $a_i, b_i \in A$. Prenons $f = \prod a_i$. \end{proof}

\begin{lemma} \label{lemma-nontrivial-normal-bundle} Dans la situation \ref{situation-vanishing}, supposons $X$ normal. Soit $Z \subset X$ un diviseur de Cartier effectif non vide tel que $Z \subset X_s$ ensemblistement. Alors le faisceau conormal de $Z$ n'est pas trivial. Plus précisément, il existe un $i$ tel que $C_i \subset Z$ et $\deg(\mathcal{C}_{Z/X}|_{C_i}) > 0$. \end{lemma}

\begin{proof} Nous utiliserons sans autre mention les observations faites après la situation \ref{situation-vanishing}. Soit $f$ une fonction comme dans le lemme \ref{lemma-nice-meromorphic-function}. Soit $\xi_i \in C_i$ le point générique. Soit $\mathcal{O}_i$ l'anneau local de $X$ en $\xi_i$. Alors $\mathcal{O}_i$ est un anneau de valuation discrète. Soit $e_i$ la valuation de $f$ dans $\mathcal{O}_i$, de sorte que $e_i > 0$. Soit $h_i \in \mathcal{O}_i$ une équation locale de $Z$ et soit $d_i$ sa valuation. Alors $d_i \geq 0$. Choisissons et fixons $i$ de sorte que $d_i/e_i$ soit maximal (alors $d_i > 0$ puisque $Z$ n'est pas vide). Remplaçons $f$ par $f^{d_i}$ et $Z$ par $e_iZ$. Cela est permis en vertu de la relation $\mathcal{O}_X(e_i Z) = \mathcal{O}_X(Z)^{\otimes e_i}$, de la relation entre le faisceau conormal et $\mathcal{O}_X(Z)$ (voir Diviseurs, lemmes \ref{divisors-lemma-invertible-sheaf-sum-effective-Cartier-divisors} et \ref{divisors-lemma-conormal-effective-Cartier-divisor}, et du fait que le degré est multiplié par $e_i$, voir Variétés, lemme \ref{varieties-lemma-degree-tensor-product}. Soit $\mathcal{I}$ le faisceau d'idéaux de $Z$, de sorte que $\mathcal{C}_{Z/X} = \mathcal{I}|_Z$. Considérons l'image $\overline{f}$ de $f$ dans $\Gamma(Z, \mathcal{O}_Z)$. D'après les choix effectués ci-dessus, $\overline{f}$ s'annule aux points génériques des composantes irréductibles de $Z$ (ce sont tous des points génériques de $C_j$ puisque $Z$ est contenu dans la fibre spéciale). D'autre part, $Z$ est $(S_1)$ d'après Diviseurs, lemme \ref{divisors-lemma-normal-effective-Cartier-divisor-S1}. Ainsi le schéma $Z$ n'a pas de point associé immergé et nous concluons que $\overline{f} = 0$ (Diviseurs, lemmes \ref{divisors-lemma-S1-no-embedded} et \ref{divisors-lemma-restriction-injective-open-contains-weakly-ass}). Par conséquent, $f$ est une section globale de $\mathcal{I}$ qui engendre $\mathcal{I}_{\xi_i}$ par construction. L'image $s_i$ de $f$ dans $\Gamma(C_i, \mathcal{I}|_{C_i})$ est donc non nulle. Toutefois, notre choix de $f$ garantit que $s_i$ possède un zéro en $x_i$. Le degré de $\mathcal{I}|_{C_i}$ est donc $> 0$ d'après Variétés, lemme \ref{varieties-lemma-check-invertible-sheaf-trivial}. \end{proof}

\begin{lemma} \label{lemma-H1-injective} Dans la situation \ref{situation-vanishing}, supposons $X$ normal et $A$ de Nagata. L'application $$
H^1(X, \mathcal{O}_X) \longrightarrow H^1(f^{-1}(U), \mathcal{O}_X)
$$ est injective. \end{lemma}

\begin{proof} Soit $0 \to \mathcal{O}_X \to \mathcal{E} \to \mathcal{O}_X \to 0$ l'extension correspondant à un élément non trivial $\xi$ de $H^1(X, \mathcal{O}_X)$ (Cohomologie, lemme \ref{cohomology-lemma-h1-extensions}). Soit $\pi : P = \mathbf{P}(\mathcal{E}) \to X$ le fibré projectif associé à $\mathcal{E}$. La surjection $\mathcal{E} \to \mathcal{O}_X$ définit une section $\sigma : X \to P$ dont le faisceau conormal est isomorphe à $\mathcal{O}_X$ (Diviseurs, lemme \ref{divisors-lemma-conormal-sheaf-section-projective-bundle}). Si la restriction de $\xi$ à $f^{-1}(U)$ est triviale, on obtient une application $\mathcal{E}|_{f^{-1}(U)} \to \mathcal{O}_{f^{-1}(U)}$ qui scinde l'injection $\mathcal{O}_X \to \mathcal{E}$. Cela définit une seconde section $\sigma' : f^{-1}(U) \to P$ disjointe de $\sigma$. Puisque $\xi$ est non trivial, nous concluons que $\sigma'$ ne peut pas se prolonger à tout $X$ tout en restant disjointe de $\sigma$. Soit $X' \subset P$ l'image schématique de $\sigma'$ (Morphismes, définition \ref{morphisms-definition-scheme-theoretic-image}). Considérons le diagramme $$
\xymatrix{
& X' \ar[rd]_g \ar[r] & P \ar[d]_\pi \\
f^{-1}(U) \ar[ru]_{\sigma'} \ar[rr] & & X \ar@/_/[u]_\sigma
}
$$ Le morphisme $P \setminus \sigma(X) \to X$ est affine. Si $X' \cap \sigma(X) = \emptyset$, alors $X' \to X$ est à la fois affine et propre, donc fini (Morphismes, lemme \ref{morphisms-lemma-finite-proper}), puis un isomorphisme (car $X$ est normal ; voir Morphismes, lemme \ref{morphisms-lemma-finite-birational-over-normal}). Cela est impossible, comme indiqué ci-dessus.

\medskip\noindent Soit $X^\nu$ la normalisation de $X'$. Puisque $A$ est de Nagata, $X^\nu \to X'$ est fini (Morphismes, lemmes \ref{morphisms-lemma-nagata-normalization} et \ref{morphisms-lemma-ubiquity-nagata}). Soit $Z \subset X^\nu$ l'image réciproque du diviseur de Cartier effectif $\sigma(X) \subset P$. D'après ce qui précède, $Z$ n'est pas vide et est contenu dans la fibre fermée de $X^\nu \to S$. Puisque $P \to X$ est lisse, $\sigma(X)$ est un diviseur de Cartier effectif (Diviseurs, lemme \ref{divisors-lemma-section-smooth-regular-immersion}). Par conséquent, $Z \subset X^\nu$ est lui aussi un diviseur de Cartier effectif. Puisque le faisceau conormal de $\sigma(X)$ dans $P$ est $\mathcal{O}_X$, le faisceau conormal de $Z$ dans $X^\nu$ (qui est a priori inversible) est $\mathcal{O}_Z$ d'après Morphismes, lemme \ref{morphisms-lemma-conormal-functorial-flat}. Cela contredit le lemme \ref{lemma-nontrivial-normal-bundle} et achève la démonstration. \end{proof}

\begin{lemma} \label{lemma-R1-injective} Dans la situation \ref{situation-vanishing}, supposons $X$ normal et $A$ de Nagata. Alors $$
\Hom_{D(A)}(\kappa[-1], Rf_*\mathcal{O}_X)
$$ est nul. On utilise ici $D(A) = D_\QCoh(\mathcal{O}_S)$ pour considérer $Rf_*\mathcal{O}_X$ comme un objet de $D(A)$. \end{lemma}

\begin{proof} Par adjonction entre $Rf_*$ et $Lf^*$, une telle application équivaut à une application $\alpha : Lf^*\kappa[-1] \to \mathcal{O}_X$. Remarquons que $$
H^i(Lf^*\kappa[-1]) =
\left\{
\begin{matrix}
0 & \text{si} & i > 1 \\
\mathcal{O}_{X_s} & \text{si} & i = 1 \\
\text{un }\mathcal{O}_{X_s}\text{-module} & \text{si} & i \leq 0
\end{matrix}
\right.
$$
\begingroup\emergencystretch=4em
Puisque $\Hom(H^0(Lf^*\kappa[-1]), \mathcal{O}_X) = 0$, car $\mathcal{O}_X$ est sans torsion, la suite spectrale pour $\Ext$ (Cohomologie sur les sites, exemple \ref{sites-cohomology-example-hom-complex-into-sheaf}) implique que $\Hom_{D(\mathcal{O}_X)}(Lf^*\kappa[-1], \mathcal{O}_X)$ est égal à $\Ext^1_{\mathcal{O}_X}(\mathcal{O}_{X_s}, \mathcal{O}_X)$. Nous concluons que $\alpha : Lf^*\kappa[-1] \to \mathcal{O}_X$ est donné par une extension\par\endgroup
$$
0 \to \mathcal{O}_X \to \mathcal{E} \to \mathcal{O}_{X_s} \to 0
$$ D'après le lemme \ref{lemma-H1-injective}, l'image réciproque de cette extension par la surjection $\mathcal{O}_X \to \mathcal{O}_{X_s}$ est nulle (car cette image réciproque est manifestement scindée sur $f^{-1}(U)$). Ainsi $1 \in \mathcal{O}_{X_s}$ se relève en une section globale $s$ de $\mathcal{E}$. En multipliant $s$ par le faisceau d'idéaux $\mathcal{I}$ de $X_s$, on obtient une application de $\mathcal{O}_X$-modules $c_s : \mathcal{I} \to \mathcal{O}_X$. En appliquant $f_*$, on obtient une application $A$-linéaire $f_*c_s : \mathfrak m \to A$. Comme $A$ est un anneau local noethérien normal intègre, cette application est donnée par la multiplication par un élément $a \in A$. En remplaçant $s$ par $s -  a$, on trouve que $s$ est annulé par $\mathcal{I}$ et que l'extension est triviale, comme souhaité. \end{proof}

\begin{remark} \label{remark-dualizing-setup} Soit $X$ un schéma intègre noethérien normal de dimension $2$. Dans ce cas, les propriétés suivantes sont équivalentes : \begin{enumerate} \item $X$ possède un complexe dualisant $\omega_X^\bullet$ ; \item il existe un $\mathcal{O}_X$-module cohérent $\omega_X$ tel que $\omega_X[n]$ soit un complexe dualisant, où $n$ peut être un entier quelconque. \end{enumerate} Cela résulte du fait que $X$ est de Cohen--Macaulay (Propriétés, lemme \ref{properties-lemma-normal-dimension-2-Cohen-Macaulay}) et de Dualité pour les schémas, lemme \ref{duality-lemma-dualizing-module-CM-scheme}. Dans cette situation, nous dirons que $\omega_X$ est un {\it module dualisant}, conformément à Dualité pour les schémas, section \ref{duality-section-dualizing-module}. En particulier, lorsque $A$ est un anneau local noethérien normal intègre de dimension $2$, nous dirons que {\it $A$ possède un module dualisant $\omega_A$} si la propriété ci-dessus est satisfaite. Dans ce cas, si $X \to \Spec(A)$ est une modification normale, alors $X$ possède aussi un module dualisant ; voir Dualité pour les schémas, exemple \ref{duality-example-proper-over-local}. Dans cette situation, nous notons toujours $\omega_X$ le module dualisant normalisé par rapport à $\omega_A$, c'est-à-dire tel que $\omega_X[2]$ soit le complexe dualisant normalisé relativement à $\omega_A[2]$. Voir Dualité pour les schémas, section \ref{duality-section-glue}. \end{remark}

\noindent Le théorème d'annulation de Grauert--Riemenschneider énoncé dans la proposition suivante est une conséquence formelle du lemme \ref{lemma-R1-injective} et de la théorie générale de la dualité.

\begin{proposition}[Grauert--Riemenschneider] \label{proposition-Grauert-Riemenschneider} Dans la situation \ref{situation-vanishing}, supposons que \begin{enumerate} \item $X$ soit un schéma normal, \item $A$ soit de Nagata et possède un complexe dualisant $\omega_A^\bullet$. \end{enumerate} Soit $\omega_X$ le module dualisant de $X$ (remarque \ref{remark-dualizing-setup}). Alors $R^1f_*\omega_X = 0$. \end{proposition}

\begin{proof} \begingroup\emergencystretch=3em
Dans cette démonstration, nous utiliserons l'identification $D(A) = D_\QCoh(\mathcal{O}_S)$ pour identifier les $\mathcal{O}_S$-modules quasi-cohérents aux $A$-modules. De plus, nous pouvons supposer $\omega_A^\bullet$ normalisé ; voir Complexes dualisants, section \ref{dualizing-section-dualizing-local}. Comme $X$ est un schéma noethérien normal de dimension $2$, il est de Cohen--Macaulay (Propriétés, lemme \ref{properties-lemma-normal-dimension-2-Cohen-Macaulay}). Ainsi $\omega_X^\bullet = \omega_X[2]$ (Dualité pour les schémas, lemme \ref{duality-lemma-dualizing-module-CM-scheme}, et la normalisation dans Dualité pour les schémas, exemple \ref{duality-example-proper-over-local}). Si la proposition est fausse, il existe une application non nulle $R^1f_*\omega_X \to \kappa$. Autrement dit, on obtient une application non nulle $\alpha : Rf_*\omega_X^\bullet \to \kappa[1]$. En appliquant $R\Hom_A(-, \omega_A^\bullet)$, on obtient une application non nulle\par\endgroup
$$
\beta : \kappa[-1] \longrightarrow Rf_*\mathcal{O}_X
$$, ce qui est impossible d'après le lemme \ref{lemma-R1-injective}. Pour voir que $R\Hom_A(-, \omega_A^\bullet)$ a bien l'effet annoncé, remarquons d'abord que $$
R\Hom_A(\kappa[1], \omega_A^\bullet) =
R\Hom_A(\kappa, \omega_A^\bullet)[-1] =
\kappa[-1]
$$ puisque $\omega_A^\bullet$ est normalisé et que l'on a $$
R\Hom_A(Rf_*\omega_X^\bullet, \omega_A^\bullet) =
Rf_*R\SheafHom_{\mathcal{O}_X}(\omega_X^\bullet, \omega_X^\bullet) =
Rf_*\mathcal{O}_X
$$ La première égalité résulte de Dualité pour les schémas, exemple \ref{duality-example-iso-on-RSheafHom-noetherian}, et du fait que $\omega_X^\bullet = f^!\omega_A^\bullet$ par construction ; la seconde, du fait que $\omega_X^\bullet$ est un complexe dualisant sur $X$ (ce qui remonte à Dualité pour les schémas, lemme \ref{duality-lemma-shriek-dualizing}). \end{proof}





\section{Bornitude} \label{section-bounded}

\noindent Dans cette section, nous entamons l'étude qui conduira à une réduction au cas des singularités rationnelles pour les schémas de dimension $2$.

\begin{lemma} \label{lemma-exact-sequence} Soit $(A, \mathfrak m, \kappa)$ un anneau local noethérien normal intègre de dimension $2$. Considérons un diagramme commutatif $$
\xymatrix{
X' \ar[rd]_{f'} \ar[rr]_g & & X \ar[ld]^f \\
& \Spec(A)
}
$$ où $f$ et $f'$ sont des modifications comme dans la situation \ref{situation-vanishing} et où $X$ est normal. On a alors une suite exacte courte $$
0 \to H^1(X, \mathcal{O}_X) \to H^1(X', \mathcal{O}_{X'}) \to
H^0(X, R^1g_*\mathcal{O}_{X'}) \to 0
$$ De plus, $\dim(\text{Supp}(R^1g_*\mathcal{O}_{X'})) = 0$ et $R^1g_*\mathcal{O}_{X'}$ est engendré par ses sections globales. \end{lemma}

\begin{proof} Nous utiliserons sans autre mention les observations faites après la situation \ref{situation-vanishing}. Comme $X$ est normal et $g$ est dominant et birationnel, on a $g_*\mathcal{O}_{X'} = \mathcal{O}_X$ ; voir par exemple Compléments sur les morphismes, lemme \ref{more-morphisms-lemma-geometrically-connected-fibres-towards-normal}. Puisque les fibres de $g$ sont de dimension $\leq 1$, on a $R^pg_*\mathcal{O}_{X'} = 0$ pour $p > 1$ ; voir par exemple Cohomologie des schémas, lemme \ref{coherent-lemma-higher-direct-images-zero-above-dimension-fibre}. Le support de $R^1g_*\mathcal{O}_{X'}$ est contenu dans l'ensemble des points de $X$ où les fibres de $g'$ sont de dimension $\geq 1$. Il est donc contenu dans l'ensemble des images des composantes irréductibles $C' \subset X'_s$ qui s'envoient sur des points de $X_s$ ; c'est un ensemble fini de points fermés (rappelons que $X'_s \to X_s$ est un morphisme entre schémas propres de dimension $1$ sur $\kappa$). Alors $R^1g_*\mathcal{O}_{X'}$ est engendré par ses sections globales d'après Cohomologie des schémas, lemme \ref{coherent-lemma-coherent-support-dimension-0}. En utilisant le morphisme $f : X \to S$ et les références ci-dessus, on obtient $H^p(X, \mathcal{F}) = 0$ pour $p > 1$ et pour tout $\mathcal{O}_X$-module cohérent $\mathcal{F}$. La suite exacte courte du lemme résulte donc de la suite spectrale de Leray pour $g$ et $\mathcal{O}_{X'}$ ; voir Cohomologie, lemme \ref{cohomology-lemma-Leray}. \end{proof}

\begin{lemma} \label{lemma-bound-a-torsion} Soit $(A, \mathfrak m, \kappa)$ un anneau local normal intègre de Nagata et de dimension $2$. Soit $a \in A$ non nul. Il existe un entier $N$ tel que, pour toute modification $f : X \to \Spec(A)$ avec $X$ normal, le $A$-module $$
M_{X, a} = \Coker(A \longrightarrow H^0(Z, \mathcal{O}_Z))
$$, où $Z \subset X$ est défini par $a$, soit de longueur bornée par $N$. \end{lemma}

\begin{proof} La suite exacte courte $
0 \to \mathcal{O}_X \xrightarrow{a} \mathcal{O}_X \to \mathcal{O}_Z \to 0
$ montre que \begin{equation}
\label{equation-a-torsion}
M_{X, a} = H^1(X, \mathcal{O}_X)[a]
\end{equation} Ici, $N[a] = \{n \in N \mid an = 0\}$ pour un $A$-module $N$. Ainsi, si $a$ divise $b$, alors $M_{X, a} \subset M_{X, b}$. Supposons que, pour un certain $c \in A$, les modules $M_{X, c}$ soient de longueurs bornées. Alors, pour tout $X$, on a une suite exacte $$
0 \to M_{X, c} \to M_{X, c^2} \to M_{X, c}
$$ où la deuxième flèche est donnée par la multiplication par $c$. Il s'ensuit que $M_{X, c^2}$ est également de longueur bornée. Il suffit donc de trouver un $c \in A$ pour lequel le lemme est vrai et tel que $a$ divise $c^n$ pour un certain $n > 0$. D'après Compléments d'algèbre, lemme \ref{more-algebra-lemma-divides-radical}, nous pouvons supposer que $A/(a)$ est un anneau réduit.

\medskip\noindent Supposons $A/(a)$ réduit. Soit $A/(a) \subset B$ la normalisation de $A/(a)$ dans son anneau total des fractions. Comme $A$ est de Nagata, $\Coker(A \to B)$ est fini. Nous affirmons que la longueur de ce module fini constitue une borne. Pour le voir, considérons $f : X \to \Spec(A)$ comme dans le lemme et soit $Z' \subset Z$ l'adhérence schématique de $Z \cap f^{-1}(U)$. Alors $Z' \to \Spec(A/(a))$ est fini, par exemple d'après Variétés, lemme \ref{varieties-lemma-finite-in-codim-1}. Par conséquent, $Z' = \Spec(B')$ avec $A/(a) \subset B' \subset B$. D'autre part, nous affirmons que l'application $$
H^0(Z, \mathcal{O}_Z) \to H^0(Z', \mathcal{O}_{Z'})
$$ est injective. En effet, si $s \in H^0(Z, \mathcal{O}_Z)$ appartient au noyau, alors la restriction de $s$ à $f^{-1}(U) \cap Z$ est nulle. L'image de $s$ dans $H^1(X, \mathcal{O}_X)$ s'annule donc dans $H^1(f^{-1}(U), \mathcal{O}_X)$. D'après le lemme \ref{lemma-H1-injective}, $s$ provient d'un élément $\tilde s$ de $A$. Mais, par hypothèse, $\tilde s$ s'envoie sur zéro dans $B'$, ce qui entraîne $s = 0$. En réunissant ces faits, on voit que $M_{X, a}$ est un sous-quotient de $B'/A$ : tout élément de $B'$ ne se prolonge pas nécessairement en une section globale de $\mathcal{O}_Z$, mais, dans tous les cas, la longueur de $M_{X, a}$ est bornée par celle de $B/A$. \end{proof}

\noindent Dans certains cas, la résolution des singularités se ramène au cas des singularités rationnelles.

\begin{definition} \label{definition-reduce-to-rational} Soit $(A, \mathfrak m, \kappa)$ un anneau local normal intègre de Nagata et de dimension $2$. \begin{enumerate} \item Nous disons que $A$ {\it définit une singularité rationnelle} si, pour toute modification normale $X \to \Spec(A)$, on a $H^1(X, \mathcal{O}_X) = 0$. \item Nous disons que {\it la réduction aux singularités rationnelles est possible pour $A$} si les longueurs des $A$-modules $$
H^1(X, \mathcal{O}_X)
$$ sont bornées pour toutes les modifications $X \to \Spec(A)$ avec $X$ normal. \end{enumerate} \end{definition}

\noindent Le sens de la terminologie employée en (2) est expliqué par le lemme \ref{lemma-reduce-to-rational}. Le lemme suivant affirme, en termes approximatifs, que les anneaux locaux des modifications de $\Spec(A)$, lorsque $A$ définit une singularité rationnelle, définissent eux aussi des singularités rationnelles.

\begin{lemma} \label{lemma-rational-propagates} Soit $(A, \mathfrak m, \kappa)$ un anneau local normal intègre de Nagata et de dimension $2$ qui définit une singularité rationnelle. Soit $A \subset B$ une extension locale d'anneaux intègres ayant le même corps des fractions, essentiellement de type fini, telle que $\dim(B) = 2$ et que $B$ soit normal. Alors $B$ définit une singularité rationnelle. \end{lemma}

\begin{proof} Choisissons une $A$-algèbre de type fini $C$ telle que $B = C_\mathfrak q$ pour un idéal premier $\mathfrak q \subset C$. Après avoir remplacé $C$ par l'image de $C$ dans $B$, nous pouvons supposer que $C$ est un anneau intègre dont le corps des fractions est celui de $A$. On peut alors choisir une immersion fermée $\Spec(C) \to \mathbf{A}^n_A$ et prendre l'adhérence dans $\mathbf{P}^n_A$ pour conclure que $B$ est isomorphe à $\mathcal{O}_{X, x}$ pour un certain point fermé $x \in X$ d'une modification projective $X \to \Spec(A)$. (Morphismes, lemme \ref{morphisms-lemma-dimension-formula}, montre que $\kappa(x)$ est fini sur $\kappa$, puis Morphismes, lemme \ref{morphisms-lemma-algebraic-residue-field-extension-closed-point-fibre}, montre que $x$ est un point fermé.) Soit $\nu : X^\nu \to X$ la normalisation. Comme $A$ est de Nagata, le morphisme $\nu$ est fini (Morphismes, lemme \ref{morphisms-lemma-nagata-normalization}). Ainsi $X^\nu$ est projectif sur $A$ d'après Compléments sur les morphismes, lemme \ref{more-morphisms-lemma-category-projective}. Puisque $B = \mathcal{O}_{X, x}$ est normal, on a $\mathcal{O}_{X, x} = (\nu_*\mathcal{O}_{X^\nu})_x$. Il existe donc un unique point $x^\nu \in X^\nu$ au-dessus de $x$, et $\mathcal{O}_{X^\nu, x^\nu} = \mathcal{O}_{X, x}$. Nous pouvons ainsi supposer $X$ normal et projectif sur $A$. Soit $Y \to \Spec(\mathcal{O}_{X, x}) = \Spec(B)$ une modification avec $Y$ normal. Il faut montrer que $H^1(Y, \mathcal{O}_Y) = 0$. D'après Limites, lemme \ref{limits-lemma-modifications}, il existe un morphisme de schémas $g : X' \to X$ qui est un isomorphisme au-dessus de $X \setminus \{x\}$ et tel que $X' \times_X \Spec(\mathcal{O}_{X, x})$ soit isomorphe à $Y$. Alors $g$ est une modification, car il est propre d'après Limites, lemme \ref{limits-lemma-modifications-properties}. L'anneau local de $X'$ en un point de $x'$ est soit isomorphe à l'anneau local de $X$ en $g(x')$ si $g(x') \not = x$, soit, si $g(x') = x$, l'anneau local de $X'$ en $x'$ est isomorphe à celui de $Y$ au point correspondant. Ainsi $X'$ est normal puisque $X$ et $Y$ le sont. Par conséquent, $H^1(X', \mathcal{O}_{X'}) = 0$ d'après notre hypothèse sur $A$. Le lemme \ref{lemma-exact-sequence} donne $R^1g_*\mathcal{O}_{X'} = 0$. Cela signifie manifestement que $H^1(Y, \mathcal{O}_Y) = 0$, comme souhaité. \end{proof}

\begin{lemma} \label{lemma-reduce-to-rational} Soit $(A, \mathfrak m, \kappa)$ un anneau local normal intègre de Nagata et de dimension $2$. Si la réduction aux singularités rationnelles est possible pour $A$, alors il existe une suite finie d'éclatements normalisés $$
X = X_n \to X_{n - 1} \to \ldots \to X_1 \to X_0 = \Spec(A)
$$ en des points fermés telle que, pour tout point fermé $x \in X$, l'anneau local $\mathcal{O}_{X, x}$ définisse une singularité rationnelle. En particulier, $X \to \Spec(A)$ est une modification et $X$ est un schéma normal projectif sur $A$. \end{lemma}

\begin{proof} Choisissons une modification $X \to \Spec(A)$ avec $X$ normal qui maximise la longueur de $H^1(X, \mathcal{O}_X)$. D'après le lemme \ref{lemma-exact-sequence}, pour toute modification ultérieure $g : X' \to X$ avec $X'$ normal, on a $R^1g_*\mathcal{O}_{X'} = 0$ et $H^1(X, \mathcal{O}_X) = H^1(X', \mathcal{O}_{X'})$.

\medskip\noindent Soit $x \in X$ un point fermé. Nous allons montrer que $\mathcal{O}_{X, x}$ définit une singularité rationnelle. Soit $Y \to \Spec(\mathcal{O}_{X, x})$ une modification avec $Y$ normal. Il faut montrer que $H^1(Y, \mathcal{O}_Y) = 0$. D'après Limites, lemme \ref{limits-lemma-modifications}, il existe un morphisme de schémas $g : X' \to X$ qui est un isomorphisme au-dessus de $X \setminus \{x\}$ et tel que $X' \times_X \Spec(\mathcal{O}_{X, x})$ soit isomorphe à $Y$. Alors $g$ est une modification, car il est propre d'après Limites, lemme \ref{limits-lemma-modifications-properties}. L'anneau local de $X'$ en un point de $x'$ est soit isomorphe à l'anneau local de $X$ en $g(x')$ si $g(x') \not = x$, soit, si $g(x') = x$, l'anneau local de $X'$ en $x'$ est isomorphe à celui de $Y$ au point correspondant. Ainsi $X'$ est normal puisque $X$ et $Y$ le sont. Par maximalité, on a $R^1g_*\mathcal{O}_{X'} = 0$ (voir le premier paragraphe). Cela signifie manifestement que $H^1(Y, \mathcal{O}_Y) = 0$, comme souhaité.

\medskip\noindent En conclusion, nous avons trouvé une modification normale $X$ de $\Spec(A)$ telle que tous les anneaux locaux de $X$ aux points fermés définissent des singularités rationnelles. Choisissons alors une suite d'éclatements normalisés $X_n \to \ldots \to X_1 \to \Spec(A)$ telle que $X_n$ domine $X$ ; voir le lemme \ref{lemma-dominate-by-normalized-blowing-up}. Pour un point fermé $x' \in X_n$ s'envoyant sur $x \in X$, on peut appliquer le lemme \ref{lemma-rational-propagates} à l'homomorphisme d'anneaux $\mathcal{O}_{X, x} \to \mathcal{O}_{X_n, x'}$ pour voir que $\mathcal{O}_{X_n, x'}$ définit une singularité rationnelle. \end{proof}

\begin{lemma} \label{lemma-go-up-separable} Soit $A \to B$ un homomorphisme local fini injectif entre anneaux locaux normaux intègres de Nagata et de dimension $2$. Supposons que l'extension induite des corps des fractions soit séparable. Si la réduction aux singularités rationnelles est possible pour $A$, elle l'est aussi pour $B$. \end{lemma}

\begin{proof} Soit $n$ le degré de l'extension de corps des fractions $L/K$. Soit $\text{Trace}_{L/K} : L \to K$ la trace. Puisque l'extension est finie séparable, la forme bilinéaire $(h, g) \mapsto \text{Trace}_{L/K}(fg)$ est non dégénérée sur $L$, considéré comme $K$-espace vectoriel. Voir Corps, lemme \ref{fields-lemma-separable-trace-pairing}. Choisissons $b_1, \ldots, b_n \in B$ qui forment une base de $L$ sur $K$. D'après ce qui précède, $d = \det(\text{Trace}_{L/K}(b_ib_j)) \in A$ est non nul.

\medskip\noindent Soit $Y \to \Spec(B)$ une modification avec $Y$ normal. Il existe un éclatement $U$-admissible $X'$ de $\Spec(A)$ tel que le transformé strict $Y'$ de $Y$ soit fini sur $X'$ ; voir Compléments sur la platitude, lemme \ref{flat-lemma-finite-after-blowing-up}. Considérons le diagramme $$
\xymatrix{
Y' \ar[d] \ar[r] & Y \ar[r] & \Spec(B) \ar[d] \\
X' \ar[rr] & & \Spec(A)
}
$$ Après avoir remplacé $X'$ et $Y'$ par leurs normalisations, nous pouvons supposer que $X'$ et $Y'$ sont des modifications normales de $\Spec(A)$ et $\Spec(B)$. On se ramène ainsi au cas où il existe un diagramme commutatif $$
\xymatrix{
Y \ar[d]_\pi \ar[r]_-g & \Spec(B) \ar[d] \\
X \ar[r]^-f & \Spec(A)
}
$$ dans lequel $X$ et $Y$ sont des modifications normales de $\Spec(A)$ et $\Spec(B)$ et où $\pi$ est fini.

\medskip\noindent L'application trace de $L$ sur $K$ se prolonge en une application de $\mathcal{O}_X$-modules $\text{Trace} : \pi_*\mathcal{O}_Y \to \mathcal{O}_X$. Considérons l'application $$
\Phi : \pi_*\mathcal{O}_Y \longrightarrow \mathcal{O}_X^{\oplus n},\quad
s \longmapsto (\text{Trace}(b_1s), \ldots, \text{Trace}(b_ns))
$$ Cette application est injective (car elle l'est au point générique) et il existe une application $$
\mathcal{O}_X^{\oplus n} \longrightarrow \pi_*\mathcal{O}_Y,\quad
(s_1, \ldots, s_n) \longmapsto \sum b_i s_i
$$ dont le composé avec $\Phi$ a pour matrice $\text{Trace}(b_ib_j)$. Le conoyau de $\Phi$ est donc annulé par $d$. On obtient ainsi une suite exacte $$
H^0(X, \Coker(\Phi)) \to H^1(Y, \mathcal{O}_Y) \to
H^1(X, \mathcal{O}_X)^{\oplus n}
$$ Comme le terme de droite est borné par hypothèse, il suffit de montrer que la $d$-torsion de $H^1(Y, \mathcal{O}_Y)$ est bornée. C'est le contenu du lemme \ref{lemma-bound-a-torsion} et de (\ref{equation-a-torsion}). \end{proof}

\begin{lemma} \label{lemma-regular-rational} Soit $A$ un anneau local régulier de Nagata et de dimension $2$. Alors $A$ définit une singularité rationnelle. \end{lemma}

\begin{proof} (L'hypothèse selon laquelle $A$ est de Nagata n'est pas nécessaire pour cette démonstration, mais nous n'avons défini la notion de singularité rationnelle que pour les anneaux locaux normaux intègres de Nagata et de dimension $2$.) Soit $X \to \Spec(A)$ une modification avec $X$ normal. D'après le lemme \ref{lemma-dominate-by-blowing-up-in-points}, on peut dominer $X$ par un schéma $X_n$ qui est le dernier terme d'une suite $$
X_n \to X_{n - 1} \to \ldots \to X_1 \to X_0 = \Spec(A)
$$ d'éclatements en des points fermés. D'après le lemme \ref{lemma-blowup-regular}, les schémas $X_i$ sont réguliers, donc en particulier normaux (Algèbre, lemme \ref{algebra-lemma-regular-normal}). Le lemme \ref{lemma-exact-sequence} donne $H^1(X, \mathcal{O}_X) \subset H^1(X_n, \mathcal{O}_{X_n})$. Il suffit donc de démontrer $H^1(X_n, \mathcal{O}_{X_n}) = 0$. En utilisant de nouveau le lemme \ref{lemma-exact-sequence}, il suffit de montrer $R^1(X_i \to X_{i - 1})_*\mathcal{O}_{X_i} = 0$ pour $i = 1, \ldots, n$. Cela résulte du lemme \ref{lemma-cohomology-of-blowup}. \end{proof}

\begin{lemma} \label{lemma-bound-dualizing-implies-bound} Soit $A$ un anneau local normal intègre de Nagata et de dimension $2$ qui possède un complexe dualisant $\omega_A^\bullet$. S'il existe un élément non nul $d \in A$ tel que, pour toute modification normale $X \to \Spec(A)$, le conoyau de l'application trace $$
\Gamma(X, \omega_X) \to \omega_A
$$ soit annulé par $d$, alors la réduction aux singularités rationnelles est possible pour $A$. \end{lemma}

\begin{proof} \begingroup\emergencystretch=2em
Pour $X \to \Spec(A)$ comme dans l'énoncé, il faut borner $H^1(X, \mathcal{O}_X)$. Soit $\omega_X$ le module dualisant de $X$ comme dans l'énoncé de Grauert--Riemenschneider (proposition \ref{proposition-Grauert-Riemenschneider}). L'application trace est l'application $Rf_*\omega_X \to \omega_A$ décrite dans Dualité pour les schémas, section \ref{duality-section-trace}. D'après Grauert--Riemenschneider, on a $Rf_*\omega_X = f_*\omega_X$ ; l'application trace induit donc bien une application $\Gamma(X, \omega_X) \to \omega_A$. Par dualité, on a $Rf_*\omega_X = R\Hom_A(Rf_*\mathcal{O}_X, \omega_A)$ (on utilise ici que $\omega_X[2]$ est le complexe dualisant sur $X$ normalisé relativement à $\omega_A[2]$ ; voir Dualité pour les schémas, lemme \ref{duality-lemma-duality-bootstrap}, ou plus directement la section \ref{duality-section-duality}, voire, encore plus directement, l'exemple \ref{duality-example-iso-on-RSheafHom-noetherian}). Le triangle distingué\par\endgroup
$$
A \to Rf_*\mathcal{O}_X \to R^1f_*\mathcal{O}_X[-1] \to A[1]
$$ est transformé par $R\Hom_A(-, \omega_A)$ en la suite exacte courte $$
0 \to f_*\omega_X \to \omega_A \to
\Ext_A^2(R^1f_*\mathcal{O}_X, \omega_A) \to 0
$$ (et $\Ext_A^i(R^1f_*\mathcal{O}_X, \omega_A) = 0$ pour $i \not = 2$ ; cela résulte aussi de la discussion ci-dessous). Puisque $R^1f_*\mathcal{O}_X$ est supporté dans $\{\mathfrak m\}$, le théorème de dualité locale affirme que $$
\Ext_A^2(R^1f_*\mathcal{O}_X, \omega_A) =
\Ext_A^0(R^1f_*\mathcal{O}_X, \omega_A[2]) =
\Hom_A(R^1f_*\mathcal{O}_X, E)
$$ est le dual de Matlis de $R^1f_*\mathcal{O}_X$ (et que les autres groupes d'extensions sont nuls) ; voir Complexes dualisants, lemme \ref{dualizing-lemma-special-case-local-duality}. Par l'équivalence de catégories inhérente à la dualité de Matlis (Complexes dualisants, proposition \ref{dualizing-proposition-matlis}), si $R^1f_*\mathcal{O}_X$ n'est pas annulé par $d$, alors le $\Ext^2$ ci-dessus ne l'est pas non plus. On voit donc que $H^1(X, \mathcal{O}_X)$ est annulé par $d$. La bornitude requise résulte alors du lemme \ref{lemma-bound-a-torsion} et de (\ref{equation-a-torsion}). \end{proof}

\begin{lemma} \label{lemma-compare-differentials-dualizing} Soit $p$ un nombre premier. Soit $A$ un anneau local régulier de dimension $2$ et de caractéristique $p$. Soit $A_0 \subset A$ un sous-anneau tel que $\Omega_{A/A_0}$ soit libre de rang $r < \infty$. Posons $\omega_A = \Omega^r_{A/A_0}$. Si $X \to \Spec(A)$ est obtenu par une suite d'éclatements en des points fermés, alors il existe une application $$
\varphi_X : (\Omega^r_{X/\Spec(A_0)})^{**} \longrightarrow \omega_X
$$ qui prolonge l'identification donnée au point générique. \end{lemma}

\begin{proof} Remarquons que $A$ est de Gorenstein (Complexes dualisants, lemme \ref{dualizing-lemma-regular-gorenstein}) ; le module inversible $\omega_A$ sert donc bien de module dualisant. En outre, tout $X$ comme dans le lemme possède un module dualisant inversible $\omega_X$, puisque $X$ est régulier (donc de Gorenstein) et propre sur $A$ ; voir la remarque \ref{remark-dualizing-setup} et le lemme \ref{lemma-blowup-regular}. Supposons construite l'application $\varphi_X : (\Omega^r_{X/A_0})^{**} \to \omega_X$ et supposons que $b : X' \to X$ soit l'éclatement en un point fermé. Posons $\Omega^r_X = (\Omega^r_{X/A_0})^{**}$ et $\Omega^r_{X'} = (\Omega^r_{X'/A_0})^{**}$. Puisque $\omega_{X'} = b^!(\omega_X)$, une application $\Omega^r_{X'} \to \omega_{X'}$ équivaut à une application $Rb_*(\Omega^r_{X'}) \to \omega_X$. Voir la discussion de la remarque \ref{remark-dualizing-setup} et Dualité pour les schémas, section \ref{duality-section-duality}. Il suffit donc à son tour de construire une application $$
Rb_*(\Omega^r_{X'}) \longrightarrow \Omega^r_X
$$ Les faisceaux $\Omega^r_{X'}$ et $\Omega^r_X$ sont inversibles ; voir Diviseurs, lemme \ref{divisors-lemma-reflexive-over-regular-dim-2}. Considérons la suite exacte $$
b^*\Omega_{X/A_0} \to \Omega_{X'/A_0} \to \Omega_{X'/X} \to 0
$$ Un calcul local montre que $\Omega_{X'/X}$ est isomorphe à un module inversible sur le diviseur exceptionnel $E$ ; voir le lemme \ref{lemma-differentials-of-blowup}. Il en résulte que l'on a l'une des deux possibilités $$
\Omega^r_{X'} \cong (b^*\Omega^r_X)(E)
\quad\text{ou}\quad
\Omega^r_{X'} \cong b^*\Omega^r_X
$$ ; voir Diviseurs, lemme \ref{divisors-lemma-wedge-product-ses}. (La seconde possibilité ne se produit jamais en caractéristique zéro, mais peut se produire en caractéristique $p$.) Dans les deux cas, on a $R^1b_*(\Omega^r_{X'}) = 0$ et $b_*(\Omega^r_{X'}) = \Omega^r_X$ d'après le lemme \ref{lemma-cohomology-of-blowup}. \end{proof}

\begin{lemma} \label{lemma-go-up-degree-p} Soit $p$ un nombre premier. Soit $A$ un anneau local régulier complet de dimension $2$ et de caractéristique $p$. Soit $L/K$ une extension inséparable de degré $p$ du corps des fractions $K$ de $A$. Soit $B \subset L$ la clôture intégrale de $A$. Alors la réduction aux singularités rationnelles est possible pour $B$. \end{lemma}

\begin{proof} On a $A = k[[x, y]]$. Écrivons $L = K[x]/(x^p - f)$ pour un certain $f \in A$ et notons $g \in B$ la classe de congruence de $x$, c'est-à-dire l'élément tel que $g^p = f$. D'après Algèbre, lemme \ref{algebra-lemma-derivative-zero-pth-power}, $\text{d}f$ est non nul dans $\Omega_{K/\mathbf{F}_p}$. D'après Compléments d'algèbre, lemme \ref{more-algebra-lemma-power-series-ring-subfields}, il existe un sous-corps $k^p \subset k' \subset k$ avec $p^e = [k : k'] < \infty$ tel que $\text{d}f$ soit non nul dans $\Omega_{K/K_0}$, où $K_0$ est le corps des fractions de $A_0 = k'[[x^p, y^p]] \subset A$. Alors $$
\Omega_{A/A_0} =
A \otimes_k \Omega_{k/k'} \oplus A \text{d}x \oplus A \text{d}y
$$ est libre fini de rang $e + 2$. Posons $\omega_A = \Omega^{e + 2}_{A/A_0}$. Considérons l'application canonique $$
\text{Tr} :
\Omega^{e + 2}_{B/A_0}
\longrightarrow
\Omega^{e + 2}_{A/A_0} = \omega_A
$$ du lemme \ref{lemma-trace-extends}. Par dualité, elle détermine une application $$
c : \Omega^{e + 2}_{B/A_0} \to \omega_B = \Hom_A(B, \omega_A)
$$ Affirmation : le conoyau de $c$ est annulé par un élément non nul de $B$.

\medskip\noindent Puisque $\text{d}f$ est non nul dans $\Omega_{A/A_0}$, on peut trouver $\eta_1, \ldots, \eta_{e + 1} \in \Omega_{A/A_0}$ tels que $\theta = \eta_1 \wedge \ldots \wedge \eta_{e + 1} \wedge \text{d}f$ soit non nul dans $\omega_A = \Omega^{e + 2}_{A/A_0}$. Pour démontrer l'affirmation, nous allons construire des éléments $\omega_i$ de $\Omega^{e + 2}_{B/A_0}$, $i = 0, \ldots, p - 1$, qui sont envoyés sur $\varphi_i \in \omega_B = \Hom_A(B, \omega_A)$ avec $\varphi_i(g^j) = \delta_{ij}\theta$ pour $j = 0, \ldots, p - 1$. Puisque $\{1, g, \ldots, g^{p - 1}\}$ est une base de $L/K$, cela démontre l'affirmation. Posons $\eta = \eta_1 \wedge \ldots \wedge \eta_{e + 1}$, de sorte que $\theta = \eta \wedge \text{d}f$. Posons $\omega_i = \eta \wedge g^{p - 1 - i}\text{d}g$. Par construction, on a alors $$
\varphi_i(g^j) = \text{Tr}(g^j \eta \wedge g^{p - 1 - i}\text{d}g) =
\text{Tr}(\eta \wedge g^{p - 1 - i + j}\text{d}g) = \delta_{ij} \theta
$$ d'après la description explicite de l'application trace dans le lemme \ref{lemma-trace-higher}.

\medskip\noindent Soit $Y \to \Spec(B)$ une modification normale. Exactement comme dans la démonstration du lemme \ref{lemma-go-up-separable}, on peut se ramener au cas où $Y$ est fini sur une modification $X$ de $\Spec(A)$. D'après le lemme \ref{lemma-dominate-by-blowing-up-in-points}, on peut même supposer que $X \to \Spec(A)$ est obtenu par une suite d'éclatements en des points fermés. Considérons le diagramme : $$
\xymatrix{
Y \ar[d]_\pi \ar[r]_-g & \Spec(B) \ar[d] \\
X \ar[r]^-f & \Spec(A)
}
$$ On peut appliquer le lemme \ref{lemma-trace-extends} à $\pi$, ce qui fournit la première flèche de $$
\pi_*(\Omega^{e + 2}_{Y/A_0})
\xrightarrow{\text{Tr}}
(\Omega^{e + 2}_{X/A_0})^{**}
\xrightarrow{\varphi_X}
\omega_X
$$ ; la seconde provient du lemme \ref{lemma-compare-differentials-dualizing} (puisque $f$ est une suite d'éclatements en des points fermés). Par dualité pour le morphisme fini $\pi$, cela correspond à une application $$
c_Y : \Omega^{e + 2}_{Y/A_0} \longrightarrow \omega_Y
$$ qui prolonge l'application $c$ ci-dessus. Ainsi, l'image de $\Gamma(Y, \omega_Y) \to \omega_B$ contient celle de $c$. Notre affirmation montre que le conoyau est annulé par un élément non nul fixé de $B$. On conclut par le lemme \ref{lemma-bound-dualizing-implies-bound}. \end{proof}






\section{Singularités rationnelles} \label{section-rational-singularities}

\noindent Dans cette section, nous ramenons le cas des points singuliers rationnels à celui des points singuliers rationnels de Gorenstein. Voir \cite{Lipman-rational} et \cite{Mattuck}.

\begin{situation} \label{situation-rational} Ici, $(A, \mathfrak m, \kappa)$ est un anneau local normal intègre de Nagata et de dimension $2$ qui définit une singularité rationnelle. Soit $s$ le point fermé de $S = \Spec(A)$ et soit $U = S \setminus \{s\}$. Soit $f : X \to S$ une modification avec $X$ normal. Nous notons $C_1, \ldots, C_r$ les composantes irréductibles de la fibre spéciale $X_s$ de $f$. \end{situation}

\begin{lemma} \label{lemma-globally-generated} Dans la situation \ref{situation-rational}, soit $\mathcal{F}$ un $\mathcal{O}_X$-module quasi-cohérent. Alors \begin{enumerate} \item $H^p(X, \mathcal{F}) = 0$ pour $p \not \in \{0, 1\}$ ; \item $H^1(X, \mathcal{F}) = 0$ si $\mathcal{F}$ est engendré par ses sections globales. \end{enumerate} \end{lemma}

\begin{proof} La partie (1) résulte de Cohomologie des schémas, lemme \ref{coherent-lemma-higher-direct-images-zero-above-dimension-fibre}. Si $\mathcal{F}$ est engendré par ses sections globales, il existe une surjection $\bigoplus_{i \in I} \mathcal{O}_X \to \mathcal{F}$. D'après la partie (1) et la suite exacte longue de cohomologie, celle-ci induit une surjection sur $H^1$. Puisque $H^1(X, \mathcal{O}_X) = 0$, car $S$ définit une singularité rationnelle, et puisque $H^1(X, -)$ commute aux sommes directes (Cohomologie, lemme \ref{cohomology-lemma-quasi-separated-cohomology-colimit}), on conclut. \end{proof}

\begin{lemma} \label{lemma-sections-powers-I-rational} Dans la situation \ref{situation-rational}, supposons que $E = X_s$ soit un diviseur de Cartier effectif. Soit $\mathcal{I}$ le faisceau d'idéaux de $E$. Alors $H^0(X, \mathcal{I}^n) = \mathfrak m^n$ et $H^1(X, \mathcal{I}^n) = 0$. \end{lemma}

\begin{proof} On a $H^0(X, \mathcal{O}_X) = A$ ; voir la discussion qui suit la situation \ref{situation-vanishing}. Alors $\mathfrak m \subset H^0(X, \mathcal{I}) \subset H^0(X, \mathcal{O}_X)$. La seconde inclusion n'est pas une égalité, car $X_s \not = \emptyset$. Ainsi $H^0(X, \mathcal{I}) = \mathfrak m$. Comme $\mathcal{I}^n = \mathfrak m^n\mathcal{O}_X$, notre lemme \ref{lemma-globally-generated} montre que $H^1(X, \mathcal{I}^n) = 0$.

\medskip\noindent Choisissons des générateurs $x_1, \ldots, x_{\mu + 1}$ de $\mathfrak m$. Ils définissent des sections globales de $\mathcal{I}$ qui l'engendrent. On obtient donc une suite exacte courte $$
0 \to \mathcal{F} \to \mathcal{O}_X^{\oplus \mu + 1} \to \mathcal{I} \to 0
$$ Alors $\mathcal{F}$ est un $\mathcal{O}_X$-module localement libre fini de rang $\mu$ et $\mathcal{F} \otimes \mathcal{I}$ est engendré par ses sections globales d'après Constructions, lemme \ref{constructions-lemma-globally-generated-omega-twist-1}. Par conséquent, $\mathcal{F} \otimes \mathcal{I}^n$ est engendré par ses sections globales pour tout $n \geq 1$. Ainsi, pour $n \geq 2$, on peut considérer la suite exacte $$
0 \to \mathcal{F} \otimes \mathcal{I}^{n - 1} \to
(\mathcal{I}^{n - 1})^{\oplus \mu + 1} \to
\mathcal{I}^n \to 0
$$ En appliquant la suite exacte longue de cohomologie et en utilisant $H^1(X, \mathcal{F} \otimes \mathcal{I}^{n - 1}) = 0$ d'après le lemme \ref{lemma-globally-generated}, on obtient que tout élément de $H^0(X, \mathcal{I}^n)$ est de la forme $\sum x_i a_i$ pour certains $a_i \in H^0(X, \mathcal{I}^{n - 1})$. Cela montre par récurrence que $H^0(X, \mathcal{I}^n) = \mathfrak m^n$. \end{proof}

\begin{lemma} \label{lemma-blow-up-normal-rational} Dans la situation \ref{situation-rational}, l'éclatement de $\Spec(A)$ en $\mathfrak m$ est normal. \end{lemma}

\begin{proof} Soit $X' \to \Spec(A)$ l'éclatement ; autrement dit, $$
X' = \text{Proj}(A \oplus \mathfrak m \oplus \mathfrak m^2 \oplus \ldots).
$$ est le Proj de l'algèbre de Rees. Cela montre en particulier que $X'$ est intègre et que $X' \to \Spec(A)$ est une modification projective. Soit $X$ la normalisation de $X'$. Puisque $A$ est de Nagata, $\nu : X \to X'$ est fini (Morphismes, lemme \ref{morphisms-lemma-nagata-normalization}). Soit $E' \subset X'$ le diviseur exceptionnel et soit $E \subset X$ son image réciproque. Soient $\mathcal{I}' \subset \mathcal{O}_{X'}$ et $\mathcal{I} \subset \mathcal{O}_X$ leurs faisceaux d'idéaux. Rappelons que $\mathcal{I}' = \mathcal{O}_{X'}(1)$ (Diviseurs, lemme \ref{divisors-lemma-blowing-up-projective}). Remarquons que $\mathcal{I} = \nu^*\mathcal{I}'$ et que $E$ est un diviseur de Cartier effectif (Diviseurs, lemme \ref{divisors-lemma-pullback-effective-Cartier-defined}). Nous cherchons à montrer que $\nu$ est un isomorphisme. Comme $\nu$ est fini, il suffit de montrer que $\mathcal{O}_{X'} \to \nu_*\mathcal{O}_X$ est un isomorphisme. Sinon, il existe un $n \geq 0$ tel que $$
H^0(X', (\mathcal{I}')^n) \not =
H^0(X', (\nu_*\mathcal{O}_X) \otimes (\mathcal{I}')^n)
$$, par exemple parce que l'on peut retrouver les $\mathcal{O}_{X'}$-modules quasi-cohérents à partir de leurs modules gradués associés ; voir Propriétés, lemme \ref{properties-lemma-ample-quasi-coherent}. La formule de projection donne $$
H^0(X', (\nu_*\mathcal{O}_X) \otimes (\mathcal{I}')^n) =
H^0(X, \nu^*(\mathcal{I}')^n) =
H^0(X, \mathcal{I}^n) = \mathfrak m^n
$$, la dernière égalité résultant du lemme \ref{lemma-sections-powers-I-rational}. D'autre part, il existe manifestement une injection $\mathfrak m^n \to H^0(X', (\mathcal{I}')^n)$. Puisque $H^0(X', (\mathcal{I}')^n)$ est sans torsion, nous concluons qu'il y a égalité pour tout $n$, donc que $X = X'$. \end{proof}

\begin{lemma} \label{lemma-cohomology-blow-up-rational} Dans la situation \ref{situation-rational}, soit $X$ l'éclatement de $\Spec(A)$ en $\mathfrak m$. Soit $E \subset X$ le diviseur exceptionnel. Avec $\mathcal{O}_X(1) = \mathcal{I}$ comme d'habitude et $\mathcal{O}_E(1) = \mathcal{O}_X(1)|_E$, on a : \begin{enumerate} \item $E$ est une courbe propre de Cohen--Macaulay sur $\kappa$ ; \item $\mathcal{O}_E(1)$ est très ample ; \item $\deg(\mathcal{O}_E(1)) \geq 1$, l'égalité n'ayant lieu que si $A$ est un anneau local régulier ; \item $H^1(E, \mathcal{O}_E(n)) = 0$ pour $n \geq 0$ ; \item $H^0(E, \mathcal{O}_E(n)) = \mathfrak m^n/\mathfrak m^{n + 1}$ pour $n \geq 0$. \end{enumerate} \end{lemma}

\begin{proof} Puisque $\mathcal{O}_X(1)$ est très ample par construction, sa restriction à la fibre spéciale $E$ l'est aussi. D'après le lemme \ref{lemma-blow-up-normal-rational}, le schéma $X$ est normal. Alors $E$ est de Cohen--Macaulay d'après Diviseurs, lemme \ref{divisors-lemma-normal-effective-Cartier-divisor-S1}. Le lemme \ref{lemma-sections-powers-I-rational} s'applique, et les suites exactes $$
0 \to \mathcal{I}^{n + 1} \to \mathcal{I}^n \to i_*\mathcal{O}_E(n) \to 0
$$ ainsi que la suite exacte longue de cohomologie donnent (4) et (5). En particulier, on a $$
\deg(\mathcal{O}_E(1)) = \chi(E, \mathcal{O}_E(1)) - \chi(E, \mathcal{O}_E) =
\dim(\mathfrak m/\mathfrak m^2) - 1
$$ d'après Variétés, définition \ref{varieties-definition-degree-invertible-sheaf}. Cela donne également (3). \end{proof}

\begin{lemma} \label{lemma-dualizing-rational} Dans la situation \ref{situation-rational}, supposons que $A$ possède un complexe dualisant $\omega_A^\bullet$. Si $\omega_X$ est le module dualisant de $X$, alors l'application trace $H^0(X, \omega_X) \to \omega_A$ est un isomorphisme et, par conséquent, il existe une application canonique $f^*\omega_A \to \omega_X$. \end{lemma}

\begin{proof} \begingroup\emergencystretch=2em
D'après Grauert--Riemenschneider (proposition \ref{proposition-Grauert-Riemenschneider}), on a $Rf_*\omega_X = f_*\omega_X$. Par dualité, on dispose d'une suite exacte courte\par\endgroup
$$
0 \to f_*\omega_X \to \omega_A \to
\Ext^2_A(R^1f_*\mathcal{O}_X, \omega_A) \to 0
$$ (voir par exemple la démonstration du lemme \ref{lemma-bound-dualizing-implies-bound}) et, puisque $A$ définit une singularité rationnelle, on obtient $f_*\omega_X = \omega_A$. \end{proof}

\begin{lemma} \label{lemma-dualizing-blow-up-rational} Dans la situation \ref{situation-rational}, supposons que $A$ possède un complexe dualisant $\omega_A^\bullet$ et ne soit pas régulier. Soit $X$ l'éclatement de $\Spec(A)$ en $\mathfrak m$, de diviseur exceptionnel $E \subset X$. Soit $\omega_X$ le module dualisant de $X$. Alors : \begin{enumerate} \item $\omega_E = \omega_X|_E \otimes \mathcal{O}_E(-1)$ ; \item $H^1(X, \omega_X(n)) = 0$ pour $n \geq 0$ ; \item l'application $f^*\omega_A \to \omega_X$ du lemme \ref{lemma-dualizing-rational} est surjective. \end{enumerate} \end{lemma}

\begin{proof} Nous utiliserons sans autre mention les résultats du lemme \ref{lemma-cohomology-blow-up-rational}. Remarquons que $\omega_E = \omega_X|_E \otimes \mathcal{O}_E(-1)$ d'après Dualité pour les schémas, lemmes \ref{duality-lemma-sheaf-with-exact-support-effective-Cartier} et \ref{duality-lemma-twisted-inverse-image-closed}. Ainsi $\omega_X|_E = \omega_E(1)$. Considérons les suites exactes courtes $$
0 \to \omega_X(n + 1) \to \omega_X(n) \to i_*\omega_E(n + 1) \to 0
$$ D'après Courbes algébriques, lemme \ref{curves-lemma-vanishing-twist}, on a $H^1(E, \omega_E(n + 1)) = 0$ pour $n \geq 0$. Les applications $$
\ldots \to H^1(X, \omega_X(2)) \to H^1(X, \omega_X(1)) \to H^1(X, \omega_X)
$$ sont donc surjectives. Puisque $H^1(X, \omega_X(n))$ est nul pour $n \gg 0$ (Cohomologie des schémas, lemme \ref{coherent-lemma-kill-by-twisting}), on en déduit (2).

\medskip\noindent D'après Courbes algébriques, lemme \ref{curves-lemma-tensor-omega-with-globally-generated-invertible}, $\omega_X|_E = \omega_E \otimes \mathcal{O}_E(1)$ est engendré par ses sections globales. Puisque nous avons vu ci-dessus que $H^1(X, \omega_X(1)) = 0$, l'application $H^0(X, \omega_X) \to H^0(E, \omega_X|_E)$ est surjective. Nous concluons que $\omega_X$ est engendré par ses sections globales ; cela donne (3), car $\Gamma(X, \omega_X) = \omega_A$ est utilisé dans le lemme \ref{lemma-dualizing-rational} pour définir l'application. \end{proof}

\begin{lemma} \label{lemma-rational-to-gorenstein} Soit $(A, \mathfrak m, \kappa)$ un anneau local normal intègre de Nagata et de dimension $2$ qui définit une singularité rationnelle. Supposons que $A$ possède un complexe dualisant. Il existe alors une suite finie d'éclatements en des points fermés singuliers $$
X = X_n \to X_{n - 1} \to \ldots \to X_1 \to X_0 = \Spec(A)
$$ telle que $X_i$ soit normal pour tout $i$ et que le faisceau dualisant $\omega_X$ de $X$ soit un $\mathcal{O}_X$-module inversible. \end{lemma}

\begin{proof} Le module dualisant $\omega_A$ est un $A$-module fini dont le germe au point générique est inversible. En effet, $\omega_A \otimes_A K$ est un module dualisant pour le corps des fractions $K$ de $A$, donc est de rang $1$. Il existe ainsi un éclatement $b : Y \to \Spec(A)$ tel que le transformé strict de $\omega_A$ relativement à $b$ soit un $\mathcal{O}_Y$-module inversible ; voir Diviseurs, lemme \ref{divisors-lemma-blowup-fitting-ideal}. D'après le lemme \ref{lemma-dominate-by-normalized-blowing-up}, on peut choisir une suite d'éclatements normalisés $$
X_n \to X_{n - 1} \to \ldots \to X_1 \to \Spec(A)
$$ telle que $X_n$ domine $Y$. Le lemme \ref{lemma-blow-up-normal-rational} et un raisonnement par récurrence montrent que chaque $X_i \to X_{i - 1}$ est simplement un éclatement.

\medskip\noindent Nous affirmons que $\omega_{X_n}$ est inversible. Puisque $\omega_{X_n}$ est un $\mathcal{O}_{X_n}$-module cohérent, il suffit de vérifier que ses germes sont des modules inversibles. Si $x \in X_n$ est un point régulier, cela résulte du fait que les schémas réguliers sont de Gorenstein (Complexes dualisants, lemme \ref{dualizing-lemma-regular-gorenstein}). Si $x$ est un point singulier de $X_n$, alors chacune des images $x_i \in X_i$ de $x$ est un point singulier (car l'éclatement d'un point régulier est régulier d'après le lemme \ref{lemma-blowup-regular}). Considérons l'application canonique $f_n^*\omega_A \to \omega_{X_n}$ du lemme \ref{lemma-dualizing-rational}. Pour tout $i$, le morphisme $X_{i + 1} \to X_i$ est soit l'éclatement de $x_i$, soit un isomorphisme en $x_i$. Puisque $x_i$ est toujours un point singulier, le lemme \ref{lemma-dualizing-blow-up-rational} et une récurrence montrent que les applications $f_i^*\omega_A \to \omega_{X_i}$ sont toujours surjectives sur les germes en $x_i$. Par conséquent, $$
(f_n^*\omega_A)_x \longrightarrow \omega_{X_n, x}
$$ est surjective. D'autre part, par notre choix de $b$, le quotient de $f_n^*\omega_A$ par son sous-module de torsion est un module inversible $\mathcal{L}$. De plus, le module dualisant est sans torsion (Dualité pour les schémas, lemme \ref{duality-lemma-dualizing-module}). Il s'ensuit que $\mathcal{L}_x \cong  \omega_{X_n, x}$, ce qui achève la démonstration. \end{proof}






\section{Arcs formels} \label{section-arcs}

\noindent Soit $X$ un schéma localement noethérien. Dans cette section, nous appelons {\it arc formel} dans $X$ un morphisme $a : T \to X$ où $T$ est le spectre d'un anneau de valuation discrète complet $R$ dont le corps résiduel $\kappa$ est identifié au corps résiduel de l'image $p$ du point fermé de $\Spec(R)$. Nous dirons alors que l'arc formel $a$ est {\it centré en $p$}. Nous disons que l'arc formel $T \to X$ est {\it non singulier} si l'application induite $\mathfrak m_p/\mathfrak m_p^2 \to \mathfrak m_R/\mathfrak m_R^2$ est surjective.

\medskip\noindent Soient $a : T \to X$ et $T = \Spec(R)$ un arc formel non singulier centré en un point fermé $p$ de $X$. Supposons $X$ localement noethérien. Soit $b : X_1 \to X$ l'éclatement de $X$ en $x$. Comme $a$ est non singulier, il existe un élément $f \in \mathfrak m_p$ qui s'envoie sur une uniformisante de $R$. En particulier, le point générique de $T$ s'envoie sur un point de $X$ distinct de $p$. Autrement dit, si $K$ désigne le corps des fractions de $R$, la restriction de $a$ définit un morphisme $\Spec(K) \to X \setminus \{p\}$. Puisque le morphisme $b$ est propre et est un isomorphisme au-dessus de $X \setminus \{x\}$, le critère valuatif de propreté fournit un unique morphisme $a_1$ qui rend commutatif le diagramme suivant $$
\xymatrix{
T \ar[r]_{a_1} \ar[rd]_a & X_1 \ar[d]^{b} \\
& X
}
$$ Soit $p_1 \in X_1$ l'image du point fermé de $T$. Remarquons que $p_1$ est un point fermé, car c'est un point $\kappa = \kappa(p)$-rationnel de la fibre de $X_1 \to X$ au-dessus de $x$. Comme on dispose d'une factorisation $$
\mathcal{O}_{X, x} \to \mathcal{O}_{X_1, p_1} \to R
$$, on voit que $a_1$ est lui aussi un arc formel non singulier.

\medskip\noindent On peut répéter le procédé et obtenir une suite d'éclatements $$
\xymatrix{
T \ar[d]_a \ar[rd]_{a_1} \ar[rrd]_{a_2} \ar[rrrd]^{a_3} \\
(X, p) & (X_1, p_1) \ar[l] & (X_2, p_2) \ar[l] &
(X_3, p_3) \ar[l] & \ldots \ar[l]
}
$$ Les suites d'éclatements de ce type se caractérisent comme suit.

\begin{lemma} \label{lemma-sequence-blowups} Soit $X$ un schéma localement noethérien. Soit $$
(X, p) = (X_0, p_0) \leftarrow (X_1, p_1) \leftarrow (X_2, p_2) \leftarrow
(X_3, p_3) \leftarrow \ldots
$$ une suite d'éclatements telle que : \begin{enumerate} \item $p_i$ soit fermé, s'envoie sur $p_{i - 1}$ et vérifie $\kappa(p_i) = \kappa(p_{i - 1})$ ; \item il existe un élément $x_1 \in \mathfrak m_p$ dont l'image dans $\mathfrak m_{p_i}$, pour $i > 0$, définit le diviseur exceptionnel $E_i \subset X_i$. \end{enumerate} Alors la suite est obtenue, comme ci-dessus, à partir d'un arc non singulier $a : T \to X$. \end{lemma}

\begin{proof} Écrivons $\mathcal{O}_n = \mathcal{O}_{X_n, p_n}$ et $\mathcal{O} = \mathcal{O}_{X, p}$. Notons $\mathfrak m \subset \mathcal{O}$ et $\mathfrak m_n \subset \mathcal{O}_n$ les idéaux maximaux.

\medskip\noindent Nous affirmons que $x_1^t \not \in \mathfrak m_n^{t + 1}$. En effet, si ce n'était pas le cas, alors, dans l'anneau local $\mathcal{O}_{n + 1}$, l'élément $x_1^t$ appartiendrait à l'idéal de $(t + 1)E_{n + 1}$. Cela contredit l'hypothèse selon laquelle $x_1$ définit $E_{n + 1}$.

\medskip\noindent Pour tout $n$, choisissons des générateurs $y_{n, 1}, \ldots, y_{n, t_n}$ de $\mathfrak m_n$. Comme $\mathfrak m_n \mathcal{O}_{n + 1} = x_1\mathcal{O}_{n + 1}$ d'après l'hypothèse (2), on peut écrire $y_{n, i} = a_{n, i} x_1$ pour certains $a_{n, i} \in \mathcal{O}_{n + 1}$. Puisque l'application $\mathcal{O}_n \to \mathcal{O}_{n + 1}$ induit, d'après (1), un isomorphisme sur les corps résiduels, on peut choisir $c_{n, i} \in \mathcal{O}_n$ ayant la même classe résiduelle que $a_{n, i}$. On voit alors que $$
\mathfrak m_n = (x_1, z_{n, 1}, \ldots, z_{n, t_n}),
\quad z_{n, i} = y_{n, i} - c_{n, i} x_1
$$ et que les éléments $z_{n, i}$ s'envoient sur des éléments de $\mathfrak m_{n + 1}^2$ dans $\mathcal{O}_{n + 1}$.

\medskip\noindent Considérons $$
J_n = \Ker(\mathcal{O} \to \mathcal{O}_n/\mathfrak m_n^{n + 1})
$$ Nous affirmons que $\mathcal{O}/J_n$ est de longueur $n + 1$ et que $\mathcal{O}/(x_1) + J_n$ est égal au corps résiduel. Pour $n = 0$, c'est immédiat. Supposons l'assertion vraie pour $n$. Soit $f \in J_n$. Alors, dans $\mathcal{O}_n$, on a $$
f = a x_1^{n + 1} + x_1^n A_1(z_{n, i}) +
x_1^{n - 1} A_2(z_{n, i}) + \ldots + A_{n + 1}(z_{n, i})
$$ pour un certain $a \in \mathcal{O}_n$ et certains $A_i$ homogènes de degré $i$ à coefficients dans $\mathcal{O}_n$. Puisque $\mathcal{O} \to \mathcal{O}_n$ identifie les corps résiduels, on peut choisir $a \in \mathcal{O}$ (raisonner comme dans la construction de $z_{n, i}$ ci-dessus). En prenant l'image dans $\mathcal{O}_{n + 1}$, on voit que $f$ et $a x_1^{n + 1}$ ont la même image modulo $\mathfrak m_{n + 1}^{n + 2}$. Puisque $x_n^{n + 1} \not \in \mathfrak m_{n + 1}^{n + 2}$, il s'ensuit que $J_n/J_{n + 1}$ est de longueur $1$, ce qui démontre l'affirmation.

\medskip\noindent Considérons $R = \lim \mathcal{O}/J_n$. C'est un quotient du complété $\mathfrak m$-adique de $\mathcal{O}$ ; c'est donc un anneau local noethérien complet. D'autre part, il n'est pas de longueur finie et $x_1$ engendre son idéal maximal. Ainsi $R$ est un anneau de valuation discrète complet. L'application $\mathcal{O} \to R$ se relève en un homomorphisme local $\mathcal{O}_n \to R$ pour tout $n$. Il y a deux façons de le montrer : (1) pour tout $n$, on peut employer un procédé analogue pour construire $\mathcal{O}_n \to R_n$, puis montrer que $\mathcal{O} \to \mathcal{O}_n \to R_n$ se factorise par un isomorphisme $R \to R_n$ ; ou (2) on peut utiliser Diviseurs, lemme \ref{divisors-lemma-characterize-affine-blowup}, pour montrer que $\mathcal{O}_n$ est un localisé d'une algèbre d'éclatement affine itérée et construire explicitement une application $\mathcal{O}_n \to R$. Cela étant dit, il est clair que notre suite d'éclatements provient de l'arc non singulier $a : T = \Spec(R) \to X$. \end{proof}

\noindent Le lemme suivant est une sorte de lemme de désingularisation de N\'eron.

\begin{lemma} \label{lemma-sequence-blowups-along-arc-becomes-nonsingular} Soit $(A, \mathfrak m, \kappa)$ un anneau local noethérien intègre de dimension $2$. Soit $A \to R$ une surjection sur un anneau de valuation discrète complet. Elle définit un arc non singulier $a : T = \Spec(R) \to \Spec(A)$. Soit $$
\Spec(A) = X_0 \leftarrow X_1 \leftarrow X_2 \leftarrow X_3 \leftarrow \ldots
$$ la suite d'éclatements construite à partir de $a$. Si $A_\mathfrak p$ est un anneau local régulier, où $\mathfrak p = \Ker(A \to R)$, alors, pour un certain $i$, le schéma $X_i$ est régulier en $x_i$. \end{lemma}

\begin{proof} Soit $x_1 \in \mathfrak m$ un élément qui s'envoie sur une uniformisante de $R$. Remarquons que $\kappa(\mathfrak p) = K$ est le corps des fractions de $R$. Écrivons $\mathfrak p = (x_2, \ldots, x_r)$ avec $r$ minimal. Si $r = 2$, alors $\mathfrak m = (x_1, x_2)$ et $A$ est régulier, ce qui démontre le lemme. Supposons $r > 2$. Après renumérotation si nécessaire, nous pouvons supposer que $x_2$ s'envoie sur une uniformisante de $A_\mathfrak p$. Alors $\mathfrak p/\mathfrak p^2 + (x_2)$ est annulé par une puissance de $x_1$. Pour $i > 2$, il existe $n_i \geq 0$ et $a_i \in A$ tels que $$
x_1^{n_i} x_i - a_i x_2 = \sum\nolimits_{2 \leq j \leq k} a_{jk} x_jx_k
$$ pour certains $a_{jk} \in A$. Si $n_i = 0$ pour un certain $i$, on peut supprimer $x_i$ de la liste des générateurs de $\mathfrak p$ et conclure par récurrence sur $r$. Si, pour un certain $i$, l'élément $a_i$ est une unité, on peut supprimer $x_2$ de la liste des générateurs de $\mathfrak p$ et conclure de la même manière. Ainsi, soit $a_i \in \mathfrak p$, soit $a_i = u_i x_1^{m_1} \bmod \mathfrak p$ pour certains $m_1 > 0$ et une unité $u_i \in A$. On se trouve donc dans l'un des deux cas $$
x_1^{n_i} x_i = \sum\nolimits_{2 \leq j \leq k} a_{jk} x_jx_k
\quad\text{ou}\quad
x_1^{n_i} x_i - u_i x_1^{m_i} x_2 =
\sum\nolimits_{2 \leq j \leq k} a_{jk} x_jx_k
$$ Nous allons montrer qu'après éclatement les entiers $n_i$ et $m_i$ diminuent, ce qui achèvera la démonstration.

\medskip\noindent Examinons ce que deviennent ces équations dans l'algèbre d'éclatement affine $A' = A[\mathfrak m/x_1]$. Comme $\mathfrak m = (x_1, \ldots, x_r)$, on voit que $A'$ est engendré sur $R$ par $y_i = x_i/x_1$ pour $i \geq 2$. Manifestement, $A \to R$ se prolonge en $A' \to R$, de noyau $(y_2, \ldots, y_r)$. On voit alors que l'on se trouve dans l'un des deux cas $$
x_1^{n_i - 1} y_i = \sum\nolimits_{2 \leq j \leq k} a_{jk} y_jy_k
\quad\text{ou}\quad
x_1^{n_i - 1} y_i - u_i x_1^{m_1 - 1} y_2 =
\sum\nolimits_{2 \leq j \leq k} a_{jk} y_jy_k
$$, ce qui achève la démonstration. \end{proof}




\section{Changement de base au complété} \label{section-aux}

\noindent Le lemme élémentaire suivant se révélera utile dans la suite.

\begin{lemma} \label{lemma-iso-completions} Soit $(A, \mathfrak m, \kappa)$ un anneau local dont l'idéal maximal $\mathfrak m$ est de type fini. Soit $X$ un schéma sur $A$. Soit $Y = X \times_{\Spec(A)} \Spec(A^\wedge)$, où $A^\wedge$ est le complété $\mathfrak m$-adique de $A$. Pour un point $q \in Y$, d'image $p \in X$ situé au-dessus du point fermé de $\Spec(A)$, l'homomorphisme d'anneaux locaux $\mathcal{O}_{X, p} \to \mathcal{O}_{Y, q}$ induit un isomorphisme sur les complétés. \end{lemma}

\begin{proof} Nous pouvons supposer $X$ affine. Écrivons alors $X = \Spec(B)$. Soit $\mathfrak q \subset B' = B \otimes_A A^\wedge$ l'idéal premier correspondant à $q$ et soit $\mathfrak p \subset B$ l'idéal premier correspondant à $p$. D'après Algèbre, lemme \ref{algebra-lemma-hathat-finitely-generated}, on a $$
B'/(\mathfrak m^\wedge)^n B' =
A^\wedge/(\mathfrak m^\wedge)^n \otimes_A B =
A/\mathfrak m^n \otimes_A B = B/\mathfrak m^n B
$$ pour tout $n$. Puisque $\mathfrak m B \subset \mathfrak p$ et $\mathfrak m^\wedge B' \subset \mathfrak q$, on voit que $B/\mathfrak p^n$ et $B'/\mathfrak q^n$ sont tous deux des quotients de l'anneau affiché ci-dessus par la $n$-ième puissance du même idéal premier. Le lemme en résulte. \end{proof}

\begin{lemma} \label{lemma-port-regularity-to-completion} Soit $(A, \mathfrak m, \kappa)$ un anneau local noethérien. Soit $X \to \Spec(A)$ un morphisme localement de type fini. Posons $Y = X \times_{\Spec(A)} \Spec(A^\wedge)$. Soit $y \in Y$, d'image $x \in X$. Alors : \begin{enumerate} \item si $\mathcal{O}_{Y, y}$ est régulier, alors $\mathcal{O}_{X, x}$ est régulier ; \item si $y$ appartient à la fibre fermée, alors $\mathcal{O}_{Y, y}$ est régulier $\Leftrightarrow \mathcal{O}_{X, x}$ est régulier ; \item si $X$ est propre sur $A$, alors $X$ est régulier si et seulement si $Y$ est régulier. \end{enumerate} \end{lemma}

\begin{proof} Puisque $A \to A^\wedge$ est fidèlement plat (Algèbre, lemme \ref{algebra-lemma-completion-faithfully-flat}), $Y \to X$ est plat. La partie (1) résulte donc d'Algèbre, lemme \ref{algebra-lemma-descent-regular}. Le lemme \ref{lemma-iso-completions} montre que le morphisme $Y \to X$ induit un isomorphisme sur les anneaux locaux complétés aux points des fibres spéciales. La partie (2) résulte alors de Compléments d'algèbre, lemme \ref{more-algebra-lemma-completion-regular}. Si $X$ est propre sur $A$, alors $Y$ est propre sur $A^\wedge$ (Morphismes, lemme \ref{morphisms-lemma-base-change-proper}), et tout point fermé de $X$ et de $Y$ appartient à la fibre fermée. Par conséquent, $Y$ est un schéma régulier si et seulement si $X$ l'est, d'après Propriétés, lemme \ref{properties-lemma-characterize-regular}. \end{proof}

\begin{lemma} \label{lemma-descend-admissible-blowup} Soit $(A, \mathfrak m)$ un anneau local noethérien de complété $A^\wedge$. Soient $U \subset \Spec(A)$ et $U^\wedge \subset \Spec(A^\wedge)$ les spectres épointés. Si $Y \to \Spec(A^\wedge)$ est un éclatement $U^\wedge$-admissible, alors il existe un éclatement $U$-admissible $X \to \Spec(A)$ tel que $Y = X \times_{\Spec(A)} \Spec(A^\wedge)$. \end{lemma}

\begin{proof} Par définition, il existe un idéal $J \subset A^\wedge$ tel que $V(J) = \{\mathfrak m A^\wedge\}$ et tel que $Y$ soit l'éclatement de $S^\wedge$ le long du sous-schéma fermé défini par $J$ ; voir Diviseurs, définition \ref{divisors-definition-admissible-blowup}. Comme $A^\wedge$ est noethérien, cela entraîne $\mathfrak m^n A^\wedge \subset J$ pour un certain $n$. Puisque $A^\wedge/\mathfrak m^n A^\wedge = A/\mathfrak m^n$, on trouve un idéal $\mathfrak m^n \subset I \subset A$ tel que $J = I A^\wedge$. Soit $X \to S$ l'éclatement en $I$. Comme $A \to A^\wedge$ est plat, on conclut que le changement de base de $X$ est $Y$ d'après Diviseurs, lemme \ref{divisors-lemma-flat-base-change-blowing-up}. \end{proof}

\begin{lemma} \label{lemma-blowup-still-good} Soit $(A, \mathfrak m, \kappa)$ un anneau local normal intègre de Nagata et de dimension $2$. Supposons que $A$ définisse une singularité rationnelle et que le complété $A^\wedge$ de $A$ soit normal. Alors : \begin{enumerate} \item $A^\wedge$ définit une singularité rationnelle ; \item si $X \to \Spec(A)$ est l'éclatement en $\mathfrak m$, alors, pour tout point fermé $x \in X$, le complété $\mathcal{O}_{X, x}$ est normal. \end{enumerate} \end{lemma}

\begin{proof} \begingroup\emergencystretch=1em
Soit $Y \to \Spec(A^\wedge)$ une modification avec $Y$ normal. Il faut montrer que $H^1(Y, \mathcal{O}_Y) = 0$. D'après Variétés, lemme \ref{varieties-lemma-modification-normal-iso-over-codimension-1}, $Y \to \Spec(A^\wedge)$ est un isomorphisme au-dessus du spectre épointé $U^\wedge = \Spec(A^\wedge) \setminus \{\mathfrak m^\wedge\}$. D'après le lemme \ref{lemma-dominate-by-scheme-modification}, il existe un éclatement $U^\wedge$-admissible $Y' \to \Spec(A^\wedge)$ qui domine $Y$. Le lemme \ref{lemma-descend-admissible-blowup} fournit un éclatement $U$-admissible $X \to \Spec(A)$ dont le changement de base à $A^\wedge$ domine $Y$. Puisque $A$ est de Nagata, on peut remplacer $X$ par sa normalisation ; après quoi $X \to \Spec(A)$ est une modification normale (mais peut ne plus être un éclatement $U$-admissible). Alors $H^1(X, \mathcal{O}_X) = 0$ puisque $A$ définit une singularité rationnelle. Il s'ensuit que $H^1(X \times_{\Spec(A)} \Spec(A^\wedge),
\mathcal{O}_{X \times_{\Spec(A)} \Spec(A^\wedge)}) = 0$ par changement de base plat (Cohomologie des schémas, lemme \ref{coherent-lemma-flat-base-change-cohomology}, et platitude de $A \to A^\wedge$ d'après Algèbre, lemme \ref{algebra-lemma-completion-flat}). On obtient $H^1(Y, \mathcal{O}_Y) = 0$ d'après le lemme \ref{lemma-exact-sequence}.\par\endgroup

\medskip\noindent Enfin, soit $X \to \Spec(A)$ l'éclatement de $\Spec(A)$ en $\mathfrak m$. Alors $Y = X \times_{\Spec(A)} \Spec(A^\wedge)$ est l'éclatement de $\Spec(A^\wedge)$ en $\mathfrak m^\wedge$. D'après le lemme \ref{lemma-blow-up-normal-rational}, $Y$ et $X$ sont tous deux normaux. D'autre part, $A^\wedge$ est excellent (Compléments d'algèbre, proposition \ref{more-algebra-proposition-ubiquity-excellent}) ; tout ouvert affine de $Y$ est donc le spectre d'un anneau normal excellent intègre (Compléments d'algèbre, lemme \ref{more-algebra-lemma-finite-type-over-excellent}). Ainsi, pour $y \in Y$, l'homomorphisme d'anneaux $\mathcal{O}_{Y, y} \to \mathcal{O}_{Y, y}^\wedge$ est régulier et Compléments d'algèbre, lemme \ref{more-algebra-lemma-normal-goes-up}, montre que $\mathcal{O}_{Y, y}^\wedge$ est normal. Si $x \in X$ est un point fermé de la fibre spéciale, il existe un unique point fermé $y \in Y$ au-dessus de $x$. Puisque $\mathcal{O}_{X, x} \to \mathcal{O}_{Y, y}$ induit un isomorphisme sur les complétés (lemme \ref{lemma-iso-completions}), on conclut. \end{proof}

\begin{lemma} \label{lemma-formally-unramified} Soit $(A, \mathfrak m)$ un anneau local noethérien. Soit $X$ un schéma sur $A$. Supposons que : \begin{enumerate} \item $A$ soit analytiquement non ramifié (Algèbre, définition \ref{algebra-definition-analytically-unramified}) ; \item $X$ soit localement de type fini sur $A$ ; \item $X \to \Spec(A)$ soit \'etale aux points génériques des composantes irréductibles de $X$. \end{enumerate} Alors la normalisation de $X$ est finie sur $X$. \end{lemma}

\begin{proof} Puisque $A$ est analytiquement non ramifié, il est réduit d'après Algèbre, lemme \ref{algebra-lemma-analytically-unramified-easy}. Comme la normalisation de $X$ ne dépend que du réduit de $X$, nous pouvons remplacer $X$ par son réduit $X_{red}$ ; remarquons que $X_{red} \to X$ est un isomorphisme sur l'ouvert $U$ où $X \to \Spec(A)$ est \'etale, car $U$ est réduit (Descente, lemme \ref{descent-lemma-reduced-local-smooth}), de sorte que la condition (3) reste vraie après ce remplacement. De plus, nous pouvons supposer, et nous le faisons, que $X = \Spec(B)$ est affine.

\medskip\noindent L'application $$
K = \prod\nolimits_{\mathfrak p \subset A\text{ minimal}} \kappa(\mathfrak p)
\longrightarrow
K^\wedge = \prod\nolimits_{\mathfrak p^\wedge \subset A^\wedge\text{ minimal}}
\kappa(\mathfrak p^\wedge)
$$ est injective car $A \to A^\wedge$ est fidèlement plat (Algèbre, lemme \ref{algebra-lemma-completion-faithfully-flat}) et induit donc une application surjective entre les ensembles d'idéaux premiers minimaux (par descente pour les homomorphismes plats, voir Algèbre, section \ref{algebra-section-going-up}). Les deux membres sont des produits finis de corps puisque nos anneaux sont noethériens. Posons $L = \prod_{\mathfrak q \subset B\text{ minimal}} \kappa(\mathfrak q)$. Notre hypothèse (3) entraîne que $L = B \otimes_A K$ et que $K \to L$ est un homomorphisme d'anneaux fini \'etale (cela tient au fait que $A \to B$ est génériquement fini ; on peut par exemple utiliser Algèbre, lemme \ref{algebra-lemma-generically-finite}, ou les résultats plus détaillés de Morphismes, section \ref{morphisms-section-generically-finite}). Puisque $B$ est réduit, on voit que $B \subset L$. Cela implique que $$
C = B \otimes_A A^\wedge \subset
L \otimes_A A^\wedge = L \otimes_K K^\wedge = M
$$. Alors $M$ est l'anneau total des fractions de $C$ et est un produit fini de corps, car c'est une algèbre séparable finie sur $K^\wedge$. Il s'ensuit que $C$ est réduit et que sa normalisation $C'$ est la clôture intégrale de $C$ dans $M$. La normalisation $B'$ de $B$ est la clôture intégrale de $B$ dans $L$. Par platitude de $A \to A^\wedge$, on obtient une application injective $B' \otimes_A A^\wedge \to M$ dont l'image est contenue dans $C'$. On a ainsi $$
B' \otimes_A A^\wedge \longrightarrow C'
$$. Comme $A^\wedge$ est de Nagata (d'après Algèbre, lemme \ref{algebra-lemma-Noetherian-complete-local-Nagata}), on voit que $C'$ est fini sur $C = B \otimes_A A^\wedge$ (voir Algèbre, lemmes \ref{algebra-lemma-Noetherian-complete-local-Nagata} et \ref{algebra-lemma-nagata-in-reduced-finite-type-finite-integral-closure}). Puisque $C$ est noethérien, on conclut que $B' \otimes_A A^\wedge$ est fini sur $C = B \otimes_A A^\wedge$. Par descente fidèlement plate (Algèbre, lemme \ref{algebra-lemma-descend-properties-modules}), on voit donc que $B'$ est fini sur $B$, ce qu'il fallait démontrer. \end{proof}

\begin{lemma} \label{lemma-normalization-completion} Soit $(A, \mathfrak m, \kappa)$ un anneau local noethérien. Soit $X \to \Spec(A)$ un morphisme localement de type fini. Posons $Y = X \times_{\Spec(A)} \Spec(A^\wedge)$. Si le complémentaire de la fibre spéciale dans $Y$ est normal, alors la normalisation $X^\nu \to X$ est finie et le changement de base de $X^\nu$ à $\Spec(A^\wedge)$ redonne la normalisation de $Y$. \end{lemma}

\begin{proof} On se ramène immédiatement au cas où $X = \Spec(B)$ est affine, avec $B$ une $A$-algèbre de type fini. Posons $C = B \otimes_A A^\wedge$, de sorte que $Y = \Spec(C)$. Puisque $A \to A^\wedge$ est fidèlement plat, pour tout idéal premier $\mathfrak q \subset B$ il existe un idéal premier $\mathfrak r \subset C$ au-dessus de $\mathfrak q$. Alors $B_\mathfrak q \to C_\mathfrak r$ est fidèlement plat. Par conséquent, si $\mathfrak q$ n'est pas situé au-dessus de $\mathfrak m$, alors $C_\mathfrak r$ est normal par hypothèse sur $Y$, et l'on conclut que $B_\mathfrak q$ est normal d'après Algèbre, lemme \ref{algebra-lemma-descent-normal}. On voit ainsi que $X$ est normal en dehors de la fibre spéciale.

\medskip\noindent Rappelons que l'anneau local noethérien complet $A^\wedge$ est de Nagata (Algèbre, lemme \ref{algebra-lemma-Noetherian-complete-local-Nagata}). La normalisation $Y^\nu \to Y$ est donc finie (Morphismes, lemme \ref{morphisms-lemma-nagata-normalization}) et est un isomorphisme en dehors de la fibre spéciale. Écrivons $Y^\nu = \Spec(C')$. Alors $C \to C'$ est fini et est un isomorphisme en dehors de $V(\mathfrak m C)$. Puisque $B \to C$ est plat et induit un isomorphisme $B/\mathfrak m B \to C/\mathfrak m C$, il existe un homomorphisme fini d'anneaux $B \to B'$ dont le changement de base à $C$ redonne $C \to C'$. Voir Compléments d'algèbre, lemme \ref{more-algebra-lemma-application-formal-glueing} et remarque \ref{more-algebra-remark-formal-glueing-algebras}. On obtient ainsi un morphisme fini $X' \to X$ qui est un isomorphisme en dehors de la fibre spéciale et dont le changement de base redonne $Y^\nu \to Y$. D'après la discussion du premier paragraphe, $X'$ est normal aux points qui ne sont pas sur la fibre spéciale. À un point $x \in X'$ de la fibre spéciale correspondent un point $y \in Y^\nu$ et une application plate $\mathcal{O}_{X', x} \to \mathcal{O}_{Y^\nu, y}$. Puisque $\mathcal{O}_{Y^\nu, y}$ est normal, $\mathcal{O}_{X', x}$ l'est aussi ; voir Algèbre, lemme \ref{algebra-lemma-descent-normal}. Ainsi $X'$ est normal, et il s'ensuit que c'est la normalisation de $X$. \end{proof}

\begin{lemma} \label{lemma-normalized-blowup-completion} Soit $(A, \mathfrak m, \kappa)$ un anneau local noethérien intègre dont le complété $A^\wedge$ est normal. Alors, pour toute suite $$
Y_n \to Y_{n - 1} \to \ldots \to Y_1 \to \Spec(A^\wedge)
$$ d'éclatements normalisés, il existe une suite d'éclatements normalisés (propres) $$
X_n \to X_{n - 1} \to \ldots \to X_1 \to \Spec(A)
$$ dont le changement de base à $A^\wedge$ redonne la suite donnée. \end{lemma}

\begin{proof} Étant donnée la suite $Y_n \to \ldots \to Y_1 \to Y_0 = \Spec(A^\wedge)$, construisons par récurrence $X_n \to \ldots \to X_1 \to X_0 = \Spec(A)$. Le cas initial est $i = 0$. Étant donné $X_i$ dont le changement de base est $Y_i$, soit $Y'_i \to Y_i$ l'éclatement au point fermé $y_i \in Y_i$ tel que $Y_{i + 1}$ soit la normalisation de $Y_i$. Puisque les fibres fermées de $Y_i$ et de $X_i$ sont isomorphes, le point $y_i$ correspond à un point fermé $x_i$ de la fibre spéciale de $X_i$. Soit $X'_i \to X_i$ l'éclatement de $X_i$ en $x_i$. Alors le changement de base de $X'_i$ à $\Spec(A^\wedge)$ est isomorphe à $Y'_i$. D'après le lemme \ref{lemma-normalization-completion}, la normalisation $X_{i + 1} \to X'_i$ est finie et son changement de base à $\Spec(A^\wedge)$ est isomorphe à $Y_{i + 1}$. \end{proof}














\section{Points doubles rationnels} \label{section-rational-double-points}

\noindent Dans la section \ref{section-rational-singularities}, nous avons montré que la résolution des singularités rationnelles de dimension $2$ se ramène au cas de Gorenstein. Une singularité rationnelle de surface qui est de Gorenstein est un point double rationnel. Nous allons les résoudre par des calculs explicites.

\medskip\noindent D'après la discussion d'Exemples, section \ref{examples-section-bad}, il existe un anneau local noethérien normal intègre $A$ dont le complété est isomorphe à $\mathbf{C}[[x, y, z]]/(z^2)$. Dans ce cas, on pourrait dire que $A$ définit un point double rationnel ; en revanche, $\Spec(A)$ n'admet pas de résolution des singularités. Ce type de comportement ne peut pas se produire si $A$ est un anneau de Nagata ; voir Algèbre, lemme \ref{algebra-lemma-local-nagata-domain-analytically-unramified}.

\medskip\noindent La situation est toutefois pire encore : il existe un anneau local normal intègre de Nagata $A$ dont le complété est $\mathbf{C}[[x, y, z]]/(yz)$, et un autre dont le complété est $\mathbf{C}[[x, y, z]]/(y^2 - z^3)$. Il s'agit de l'exemple 2.5 de \cite{Nishimura-few}. C'est pourquoi, dans cette section, nous devons supposer que le complété de notre anneau est normal.

\begin{situation} \label{situation-rational-double-point} Dans ce qui suit, $(A, \mathfrak m, \kappa)$ est un anneau local normal intègre de Nagata, de dimension $2$, qui définit une singularité rationnelle, dont le complété est normal et qui est de Gorenstein. Nous supposons que $A$ n'est pas régulier. \end{situation}

\noindent Les arguments de cette section montreront que, dans cette situation, l'éclatement répété des points singuliers résout $\Spec(A)$. Au cours de la démonstration, nous aurons besoin du lemme suivant.

\begin{lemma} \label{lemma-issquare} Soit $\kappa$ un corps. Soit $I \subset \kappa[x, y]$ un idéal. Supposons que $$
a + b x + c y + d x^2 + exy + f y^2 \in I^2
$$ pour des $a, b, c, d, e, f \in k$ non tous nuls. Si la colongueur de $I$ dans $\kappa[x, y]$ est $> 1$, alors $a + b x + c y + d x^2 + exy + f y^2 = j(g + hx + iy)^2$ pour certains $j, g, h, i \in \kappa$. \end{lemma}

\begin{proof} Considérons les dérivées partielles $b + 2dx + ey$ et $c + ex + 2fy$. D'après les règles de Leibniz, elles appartiennent à $I$. Si l'une d'elles est non nulle, alors, après un changement linéaire de coordonnées, c'est-à-dire de la forme $x \mapsto \alpha + \beta x + \gamma y$ et $y \mapsto \delta + \epsilon x + \zeta y$, nous pouvons supposer que $x \in I$. On voit alors que $I = (x)$ ou que $I = (x, F)$, où $F$ est un polynôme unitaire de degré $\geq 2$ en $y$. Dans le premier cas, l'assertion est claire. Dans le second, remarquons que tout élément de $I^2$ peut s'écrire sous la forme $$
A(x, y) x^2 + B(y) x F + C(y) F^2
$$ pour certains $A(x, y) \in \kappa[x, y]$ et $B, C \in \kappa[y]$. Ainsi $$
a + b x + c y + d x^2 + exy + f y^2 = A(x, y) x^2 + B(y) x F + C(y) F^2
$$ et, par des considérations de degré, on voit que $B = C = 0$ et que $A$ est une constante.

\medskip\noindent Pour achever la démonstration, il reste à traiter le cas où les deux dérivées partielles sont nulles. Cela ne peut se produire qu'en caractéristique $2$, et l'on obtient alors $$
a + d x^2 + f y^2 \in I^2
$$. Nous pouvons supposer $f$ non nul (sinon, intervertissons les rôles de $x$ et de $y$). Après division par $f$, nous nous ramenons au cas où la caractéristique de $\kappa$ est $2$ et où $$
a + d x^2 + y^2 \in I^2
$$. Si $a$ et $d$ sont des carrés dans $\kappa$, c'est terminé. Sinon, il existe une dérivation $\theta : \kappa \to \kappa$ telle que $\theta(a) \not = 0$ ou $\theta(d) \not = 0$ ; voir Algèbre, lemme \ref{algebra-lemma-derivative-zero-pth-power}. Nous pouvons la prolonger en une dérivation de $\kappa[x, y]$ en posant $\theta(x) = \theta(y) = 0$. On trouve alors que $$
\theta(a) + \theta(d) x^2 \in I
$$. Le cas $\theta(d) = 0$ est absurde. Nous pouvons donc supposer que $\alpha + x^2 \in I$ pour un certain $\alpha \in \kappa$. En combinant cela avec ce qui précède, on obtient $a + \alpha d + y^2 \in I$. Par conséquent, $$
J = (\alpha + x^2, a + \alpha d + y^2) \subset I
$$, de codimension au plus $2$. Remarquons que $J/J^2$ est libre sur $\kappa[x, y]/J$, de base $\alpha + x^2$ et $a + \alpha d + y^2$. Ainsi, $a + d x^2 + y^2 =
1 \cdot (a + \alpha d + y^2) + d \cdot (\alpha + x^2) \in I^2$ implique que l'inclusion $J \subset I$ est stricte. On trouve donc un élément non nul de la forme $g + hx + iy + jxy$ dans $I$. Si $j = 0$, alors $I$ contient une forme linéaire et l'on peut conclure comme dans le premier paragraphe. Ainsi $j \not = 0$ et $\dim_\kappa(I/J) = 1$ (sinon, on pourrait trouver dans $I$ un élément comme ci-dessus avec $j = 0$). On conclut que $I$ est de la forme $(\alpha + x^2, \beta + y^2, g + hx + iy + jxy)$, avec $j \not = 0$, et de colongueur $3$. Dans ce cas, $a + dx^2 + y^2 \in I^2$ est impossible. On peut le montrer par un calcul direct, mais nous préférons raisonner comme suit. Pour établir cette assertion, nous pouvons supposer que $\kappa$ est algébriquement clos. Nous pouvons alors effectuer le changement de coordonnées $x \mapsto \sqrt{\alpha} + x$ et $y \mapsto \sqrt{\beta} + y$ et supposer que $I = (x^2, y^2, g' + h'x + i'y + jxy)$, avec le même $j$. Alors $g' = h' = i' = 0$, faute de quoi la colongueur de $I$ ne serait pas $3$. On obtient ainsi $I = (x^2, y^2, xy)$, et le résultat est clair. \end{proof}

\noindent Soit $(A, \mathfrak m, \kappa)$ comme dans la situation \ref{situation-rational-double-point}. Soit $X \to \Spec(A)$ l'éclatement en $\mathfrak m$ de $\Spec(A)$. D'après le lemme \ref{lemma-blow-up-normal-rational}, $X$ est normal. Toutes les singularités de $X$ sont rationnelles d'après le lemme \ref{lemma-rational-propagates}. Puisque $\omega_A = A$, le lemme \ref{lemma-dualizing-blow-up-rational} montre que $\omega_X \cong \mathcal{O}_X$ (voir la discussion de la remarque \ref{remark-dualizing-setup} pour les conventions). Toutes les singularités de $X$ sont donc de Gorenstein. En outre, les anneaux locaux de $X$ aux points fermés ont des complétés normaux d'après le lemme \ref{lemma-blowup-still-good}. Autrement dit, en éclatant $\Spec(A)$, on obtient une surface normale $X$ dont les points singuliers sont comme dans la situation \ref{situation-rational-double-point}. Nous l'utiliserons ci-dessous sans plus de commentaires. (Remarque : nous verrons au cours de la discussion ci-dessous que ces points singuliers sont en nombre fini.)

\medskip\noindent Soit $E \subset X$ le diviseur exceptionnel. On a $\omega_E = \mathcal{O}_E(-1)$ d'après le lemme \ref{lemma-dualizing-blow-up-rational}. Le lemme \ref{lemma-cohomology-blow-up-rational} donne $\kappa = H^0(E, \mathcal{O}_E)$. Ainsi, $E$ est une courbe de Gorenstein et, par Riemann--Roch tel qu'il est présenté dans Courbes algébriques, section \ref{curves-section-Riemann-Roch}, on a $$
\chi(E, \mathcal{O}_E) = 1 - g = -(1/2) \deg(\omega_E) =
(1/2)\deg(\mathcal{O}_E(1))
$$, où $g = \dim_\kappa H^1(E, \mathcal{O}_E) \geq 0$. Puisque $\deg(\mathcal{O}_E(1))$ est positif d'après Variétés, lemme \ref{varieties-lemma-ampleness-in-terms-of-degrees-components}, on trouve que $g = 0$ et $\deg(\mathcal{O}_E(1)) = 2$. Il s'ensuit que $$
\dim_\kappa (\mathfrak m^n/\mathfrak m^{n + 1}) = 2n + 1
$$ d'après le lemme \ref{lemma-cohomology-blow-up-rational} et Riemann--Roch sur $E$.

\medskip\noindent \begingroup\emergencystretch=2em
Choisissons $x_1, x_2, x_3 \in \mathfrak m$ qui s'envoient sur une base de $\mathfrak m/\mathfrak m^2$. Comme $\dim_\kappa(\mathfrak m^2/\mathfrak m^3) = 5$, les images de $x_i x_j$, $i \geq j$ dans cet espace vectoriel sur $\kappa$ satisfont une relation. Autrement dit, on peut trouver $a_{ij} \in A$, $i \geq j$, non tous contenus dans $\mathfrak m$, tels que\par\endgroup
$$
a_{11} x_1^2 + a_{12} x_1x_2 + a_{13}x_1x_3 + a_{22} x_2^2 +
a_{23} x_2x_3 + a_{33} x_3^2 =
\sum a_{ijk} x_ix_jx_k
$$ pour un certain $a_{ijk} \in A$, où $i \leq j \leq k$. Notons $a \mapsto \overline{a}$ l'application $A \to \kappa$. La forme quadratique $q = \sum \overline{a}_{ij} t_i t_j \in \kappa[t_1, t_2, t_3]$ est bien définie, compte tenu de nos choix, à multiplication près par un élément de $\kappa^*$. Si, au cours de nos arguments, nous constatons que $\overline{a}_{ij} = 0$ dans $\kappa$, nous pouvons absorber le terme $a_{ij} x_i x_j$ dans le membre de droite et supposer que $a_{ij} = 0$ ; cette opération modifie les $a_{ijk}$, mais pas les autres $a_{i'j'}$.

\medskip\noindent L'éclatement est recouvert par $3$ cartes affines correspondant aux « variables » $x_1, x_2, x_3$. Par symétrie, il suffit d'étudier l'une de ces cartes. Pour cela, soit $$
A' = A[\mathfrak m/x_1]
$$ l'algèbre affine d'éclatement (comme dans Algèbre, section \ref{algebra-section-blow-up}). Puisque $x_1, x_2, x_3$ engendrent $\mathfrak m$, on voit que $A'$ est engendrée par $y_2 = x_2/x_1$ et $y_3 = x_3/x_1$ sur $A$. Nous utiliserons parfois $y_1 = 1$ pour simplifier les formules. De plus, en considérant la relation ci-dessus, on trouve que $$
a_{11} + a_{12} y_2 + a_{13} y_3 + a_{22} y_2^2 +
a_{23} y_2y_3 + a_{33} y_3^2 =
x_1 (\sum a_{ijk} y_iy_jy_k)
$$ dans $A'$. Rappelons que $x_1 \in A'$ définit le diviseur exceptionnel $E$ sur cet ouvert affine de $X$, lequel est donc donné schématiquement par $$
\kappa[y_2, y_3]/
(\overline{a}_{11} + \overline{a}_{12} y_2 + \overline{a}_{13} y_3 +
\overline{a}_{22} y_2^2 + \overline{a}_{23} y_2y_3 + \overline{a}_{33} y_3^2)
$$. Autrement dit, $E \subset \mathbf{P}^2_\kappa = \text{Proj}(\kappa[t_1, t_2, t_3])$ est le sous-schéma des zéros de la forme quadratique $q$ introduite ci-dessus.

\medskip\noindent La forme quadratique $q$ est un invariant important de la singularité définie par $A$. Disons que nous sommes dans le {\bf cas II} si $q$ est le carré d'une forme linéaire multiplié par un élément de $\kappa^*$, et dans le {\bf cas I} sinon. Remarquons que nous sommes dans le cas II exactement lorsque, après modification de notre choix de $x_1, x_2, x_3$, on a $$
x_3^2 = \sum a_{ijk}x_ix_jx_k
$$ dans l'anneau local $A$.

\medskip\noindent Soit $\mathfrak m' \subset A'$ un idéal maximal situé au-dessus de $\mathfrak m$, de corps résiduel $\kappa'$. Autrement dit, $\mathfrak m'$ correspond à un point fermé $p \in E$ du diviseur exceptionnel. Rappelons que la surjection $$
\kappa[y_2, y_3] \to \kappa'
$$ a un noyau engendré par deux éléments $f_2, f_3 \in \kappa[y_2, y_3]$ (voir par exemple Algèbre, exemple \ref{algebra-example-spec-kxy}, ou la démonstration d'Algèbre, lemme \ref{algebra-lemma-dim-affine-space}). Soient $z_2, z_3 \in A'$ des éléments qui s'envoient respectivement sur $f_2, f_3$ dans $\kappa[y_2, y_3]$. On voit alors que $\mathfrak m' = (x_1, z_2, z_3)$, car $x_2$ et $x_3$ deviennent divisibles par $x_1$ dans $A'$.

\medskip\noindent {\bf Assertion.} Si $X$ est singulier en $p$, alors $\kappa' = \kappa$ ou bien nous sommes dans le cas II. En effet, si $A'_{\mathfrak m'}$ est singulier, alors $\dim_{\kappa'} \mathfrak m'/(\mathfrak m')^2 = 3$, ce qui implique que $\dim_{\kappa'} \overline{\mathfrak m}'/(\overline{\mathfrak m}')^2 = 2$, où $\overline{m}'$ est l'idéal maximal de $\mathcal{O}_{E, p} = \mathcal{O}_{X, p}/x_1\mathcal{O}_{X, p}$. Cela entraîne que $$
q(1, y_2, y_3) =
\overline{a}_{11} + \overline{a}_{12} y_2 + \overline{a}_{13} y_3 +
\overline{a}_{22} y_2^2 + \overline{a}_{23} y_2y_3 + \overline{a}_{33} y_3^2
\in (f_2, f_3)^2
$$, sans quoi il existerait une relation entre les classes de $z_2$ et de $z_3$ dans $\overline{\mathfrak m}'/(\overline{\mathfrak m}')^2$. L'assertion résulte maintenant du lemme \ref{lemma-issquare}.

\medskip\noindent Résolution dans le cas I. D'après l'assertion, tout point singulier de $X$ est $\kappa$-rationnel. Choisissons un tel point singulier $p$. Nous pouvons choisir nos $x_1, x_2, x_3 \in \mathfrak m$ de sorte que $p$ appartienne à la carte décrite ci-dessus et ait pour coordonnées $y_2 = y_3 = 0$. Comme c'est un point singulier, un raisonnement analogue à celui de la démonstration de l'assertion donne $q(1, y_2, y_3) \in (y_2, y_3)^2$. Nous pouvons donc choisir $a_{11} = a_{12} = a_{13} = 0$ et $q(t_1, t_2, t_3) = q(t_2, t_3)$. Il s'ensuit que $$
E = V(q) \subset \mathbf{P}^1_\kappa
$$ est soit la réunion de deux droites distinctes qui se rencontrent en $p$, soit une courbe de degré $2$ possédant un unique point $\kappa$-rationnel (nous omettons un petit détail ; utiliser que $q$ n'est pas, à un scalaire près, le carré d'une forme linéaire). Dans les deux cas, on conclut que $X$ possède un unique point singulier $p$, qui est $\kappa$-rationnel. Nous avons besoin d'un peu plus d'information dans ce cas. Tout d'abord, en examinant les termes de degré supérieur dans l'expression ci-dessus, on trouve que $\overline{a}_{111} = 0$, car $p$ est singulier. On peut alors écrire $a_{111} = b_{111} x_1 \bmod (x_2, x_3)$ pour un certain $b_{111} \in A$. La forme quadratique en $p$ associée aux générateurs $x_1, y_2, y_3$ de $\mathfrak m'$ est alors $$
q' =
\overline{b}_{111} t_1^2 +
\overline{a}_{112} t_1 t_2 + \overline{a}_{113} t_1 t_3 +
\overline{a}_{22} t_2^2 +
\overline{a}_{23} t_2 t_3 +
\overline{a}_{33} t_3^2
$$. On voit que $E' = V(q')$ rencontre la droite $t_1 = 0$ soit en deux points, soit en un point de degré $2$. On conclut que $p$ relève du cas I.

\medskip\noindent Supposons que l'éclatement $X' \to X$ de $X$ en $p$ ait de nouveau un point singulier $p'$. On voit alors que $p'$ est un point $\kappa$-rationnel et l'on peut éclater pour obtenir $X'' \to X'$. Si ce procédé ne s'arrête pas, on obtient une suite d'éclatements $$
\Spec(A) \leftarrow X \leftarrow X' \leftarrow X'' \leftarrow \ldots
$$. Nous voulons montrer que le lemme \ref{lemma-sequence-blowups} s'applique à cette situation. Pour cela, nous devons préciser le choix de l'élément $x_1$ de $\mathfrak m$. Supposons que $A$ relève du cas I et que $X$ ait un point singulier. Nous dirons que $x_1 \in \mathfrak m$ est une {\it bonne coordonnée} si, pour tout choix de $x_2, x_3$ (ou, de manière équivalente, pour un certain choix), la forme quadratique $q(t_1, t_2, t_3)$ possède la propriété que $q(0, t_2, t_3)$ n'est pas un carré multiplié par un scalaire. Nous avons vu ci-dessus qu'il existe une bonne coordonnée. Si $x_1$ est une bonne coordonnée, alors le point singulier $p \in E$ de $X$ n'appartient pas à l'hypersurface $t_1 = 0$, car celle-ci soit ne possède pas de point rationnel, soit, si elle en possède un, n'est pas singulière en ce point sur $X$. Remarquons que cela équivaut à dire que l'image de $x_1$ dans $\mathcal{O}_{X, p}$ définit le diviseur exceptionnel $E$. Les calculs ci-dessus montrent maintenant que, si $x_1$ est une bonne coordonnée pour $A$, alors $x_1 \in \mathfrak m'\mathcal{O}_{X, p}$ est une bonne coordonnée pour $p$. Cela utilise bien sûr le fait que la notion de bonne coordonnée ne dépend pas du choix de $x_2$, $x_3$ employé pour effectuer le calcul. Ainsi, $x_1$ s'envoie sur une bonne coordonnée en $p'$, en $p''$, etc. Le lemme \ref{lemma-sequence-blowups} s'applique donc, et notre suite d'éclatements provient d'un arc non singulier $A \to R$. Alors l'application $A^\wedge \to R$ est surjective. Puisque le complété de $A$ est normal, le lemme \ref{lemma-sequence-blowups-along-arc-becomes-nonsingular} montre qu'après un nombre fini d'éclatements $$
\Spec(A^\wedge) \leftarrow X^\wedge \leftarrow (X')^\wedge \leftarrow \ldots
$$, le schéma obtenu $(X^{(n)})^\wedge$ est régulier. Comme $(X^{(n)})^\wedge \to X^{(n)}$ induit des isomorphismes sur les anneaux locaux complétés (lemme \ref{lemma-iso-completions}), on conclut qu'il en va de même pour $X^{(n)}$.

\medskip\noindent Résolution dans le cas II. On a ici $$
x_3^2 = \sum a_{ijk}x_ix_jx_k
$$ dans $A$ pour un certain choix de générateurs $x_1, x_2, x_3$ de $\mathfrak m$. Alors $q = t_3^2$ et $E = 2C$, où $C$ est une droite. Rappelons que, dans $A'$, on obtient $$
y_3^2 = x_1(\sum a_{ijk} y_iy_jy_k)
$$. Puisque nous savons que $X$ est normal, l'anneau $\mathcal{O}_{X, \xi}$ au point générique $\xi$ de $C$ est un anneau de valuation discrète. L'élément $y_3 \in A'$ s'envoie sur une uniformisante de $\mathcal{O}_{X, \xi}$. Comme $x_1$ définit schématiquement $E$, qui est $C$ avec multiplicité $2$, on voit que $x_1$ est égal à une unité multipliée par $y_3^2$ dans $\mathcal{O}_{X, \xi}$. En considérant l'égalité ci-dessus, on conclut que $$
h(y_2) = \overline{a}_{111} + \overline{a}_{112} y_2 +
\overline{a}_{122} y_2^2 +
\overline{a}_{222} y_2^3
$$ doit être non nul dans le corps résiduel de $\xi$. Supposons maintenant que $p \in C$ définisse un point singulier. Alors $y_3$ s'annule en $p$ et $p$ doit correspondre à un zéro de $h$, d'après le raisonnement employé dans la démonstration de l'assertion ci-dessus. Si $h$ n'a pas un zéro double en $p$, alors la forme quadratique $q'$ en $p$ n'est pas un carré, et l'on conclut que $p$ relève du cas I, traité ci-dessus\footnote{L'idéal maximal en $p$ dans $A'$ est engendré par $y_3, x_1$ et par un troisième élément $g$ dont l'image dans $\kappa[y_2]$ est le diviseur premier de $h$ correspondant à $p$. Si ce diviseur premier ne divise pas $h$ deux fois, on voit alors que la forme quadratique en $p$ est de la forme $$
y_3^2 - x_1((\text{un élément})x_1 + (\text{un élément})y_3 + (\text{une unité})g)
$$, ce qui ne peut jamais être un carré dans $\kappa[y_3, x_1, g]$.}. Comme le degré de $h$ est $3$, il y a au plus un point singulier $p \in C$ relevant du cas II ; de plus, ce point est $\kappa$-rationnel. Après modification de notre choix de $x_1, x_2, x_3$, nous pouvons supposer qu'il s'agit du point $y_2 = y_3 = 0$. Alors $h = \overline{a}_{122} y_2^2 + \overline{a}_{222} y_2^3$. En outre, il faut encore que $\overline{a}_{113} = 0$ pour que la forme quadratique $q'$ ait la forme requise. Ainsi, l'anneau local $\mathcal{O}_{X, p}$ définit une singularité comme au paragraphe suivant.

\medskip\noindent Le dernier cas à traiter est celui où l'on peut choisir nos générateurs $x_1, x_2, x_3$ de $\mathfrak m$ de sorte que $$
x_3^2 + x_1(a x_2^2 + b x_2x_3 + c x_3^2) \in \mathfrak m^4
$$ pour un certain $a, b, c \in A$. C'est une sous-classe du cas II. Si $\overline{a} = 0$, alors on peut écrire $a = a_1 x_1 + a_2 x_2 + a_3 x_3$ et, après éclatement, on obtient $$
y_3^2 + x_1(a_1 x_1 y_2^2 + a_2 x_1 y_2^3 + a_3 x_1 y_2^2 y_3 +
b y_2 y_3 + c y_3^2) = x_1^2 (\sum a_{ijkl}y_iy_jy_ky_l)
$$. Cela signifie que $X$ n'est pas normal\footnote{En effet, l'équation montre qu'il est singulier le long du lieu de dimension $1$ donné par $x_1 = y_3 = 0$, ce qui est impossible pour une surface normale.}, contradiction. D'après le résultat du paragraphe précédent, si l'éclatement $X$ possède un point singulier $p$ qui relève du cas II, alors ce point est unique et $\kappa$-rationnel. En calculant les algèbres affines d'éclatement $A[\frac{\mathfrak m}{x_2}]$ et $A[\frac{\mathfrak m}{x_3}]$, le lecteur voit aisément que $p$ ne peut pas appartenir aux ouverts correspondants de $X$. Ainsi, $p$ appartient au spectre de $A[\frac{\mathfrak m}{x_1}]$. En effectuant l'éclatement comme précédemment, on voit que $p$ doit être le point de coordonnées $y_2 = y_3 = 0$ et que la nouvelle équation est de la forme $$
y_3^2 + x_1(a y_2^2 + b y_2 y_3 + c y_3^2) \in (\mathfrak m')^4
$$, qui a la même forme qu'auparavant et possède la propriété que $x_1$ définit le diviseur exceptionnel. Ainsi, si le procédé ne s'arrête pas, on obtient une suite infinie d'éclatements, et sur chacun d'eux $x_1$ définit le diviseur exceptionnel dans l'anneau local du point singulier. On peut donc achever la démonstration à l'aide des lemmes \ref{lemma-sequence-blowups} et \ref{lemma-sequence-blowups-along-arc-becomes-nonsingular} et du même raisonnement que précédemment.

\begin{lemma} \label{lemma-resolve-rational-double-points} Soit $(A, \mathfrak m, \kappa)$ un anneau local normal intègre de Nagata, de dimension $2$, qui définit une singularité rationnelle, dont le complété est normal et qui est de Gorenstein. Il existe alors une suite finie d'éclatements en des points fermés singuliers $$
X_n \to X_{n - 1} \to \ldots \to X_1 \to X_0 = \Spec(A)
$$ telle que $X_n$ soit régulier et que chacun des schémas intermédiaires $X_i$ soit normal et ne possède qu'un nombre fini de points singuliers du même type. \end{lemma}

\begin{proof} C'est exactement ce qui a été démontré dans la discussion ci-dessus. \end{proof}




\section{Propriétés impliquées} \label{section-existence-gives}

\noindent Dans cette section, nous montrons que, pour un schéma intègre noethérien, l'existence d'une altération régulière a de nombreuses conséquences. Les lecteurs qui ne s'intéressent pas aux « mauvais » anneaux noethériens peuvent omettre cette section.

\begin{lemma} \label{lemma-regular-alteration-implies} Soit $Y$ un schéma intègre noethérien. Supposons qu'il existe une altération $f : X \to Y$ avec $X$ régulier. Alors la normalisation $Y^\nu \to Y$ est finie et $Y$ possède un ouvert dense régulier. \end{lemma}

\begin{proof} Il suffit de le démontrer lorsque $Y = \Spec(A)$, où $A$ est un anneau noethérien intègre. Soit $B$ la clôture intégrale de $A$ dans son corps des fractions. Posons $C = \Gamma(X, \mathcal{O}_X)$. D'après Cohomologie des schémas, lemme \ref{coherent-lemma-proper-over-affine-cohomology-finite}, on voit que $C$ est un $A$-module fini. Comme $X$ est normal (Propriétés, lemme \ref{properties-lemma-regular-normal}), on voit que $C$ est un anneau normal intègre (Propriétés, lemme \ref{properties-lemma-normal-integral-sections}). Ainsi $B \subset C$, et l'on conclut que $B$ est fini sur $A$, puisque $A$ est noethérien.

\medskip\noindent Il existe un ouvert non vide $V \subset Y$ tel que $f^{-1}V \to V$ soit fini ; voir Morphismes, définition \ref{morphisms-definition-alteration}. Après avoir rétréci $V$, nous pouvons supposer que $f^{-1}V \to V$ est plat (Morphismes, proposition \ref{morphisms-proposition-generic-flatness}). Ainsi $f^{-1}V \to V$ est fidèlement plat. Alors $V$ est régulier d'après Algèbre, lemme \ref{algebra-lemma-descent-regular}. \end{proof}

\begin{lemma} \label{lemma-algebra-helper} Soit $(A, \mathfrak m)$ un anneau local noethérien. Soient $B \subset C$ des $A$-algèbres finies. Supposons que (a) $B$ soit un anneau normal et que (b) le complété $\mathfrak m$-adique $C^\wedge$ soit un anneau normal. Alors $B^\wedge$ est un anneau normal. \end{lemma}

\begin{proof} Considérons le diagramme commutatif $$
\xymatrix{
B \ar[r] \ar[d] & C \ar[d] \\
B^\wedge \ar[r] & C^\wedge
}
$$. Rappelons que la complétion $\mathfrak m$-adique est exacte sur la catégorie des $A$-modules finis, car elle est donnée par le produit tensoriel avec la $A$-algèbre plate $A^\wedge$ (Algèbre, lemme \ref{algebra-lemma-completion-flat}). Nous allons utiliser le critère de Serre (Algèbre, lemme \ref{algebra-lemma-criterion-normal}) pour démontrer que l'anneau noethérien $B^\wedge$ est normal. Soit $\mathfrak q \subset B^\wedge$ un idéal premier situé au-dessus de $\mathfrak p \subset B$. Si $\dim(B_\mathfrak p) \geq 2$, alors $\text{depth}(B_\mathfrak p) \geq 2$ et, puisque $B_\mathfrak p \to B^\wedge_\mathfrak q$ est plat, on trouve que $\text{depth}(B^\wedge_\mathfrak q) \geq 2$ (Algèbre, lemme \ref{algebra-lemma-apply-grothendieck}). Si $\dim(B_\mathfrak p) \leq 1$, alors $B_\mathfrak p$ est soit un anneau de valuation discrète, soit un corps. Dans ce cas, $C_\mathfrak p$ est fidèlement plat sur $B_\mathfrak p$ (car il est fini et sans torsion). Par conséquent, $B^\wedge_\mathfrak p \to C^\wedge_\mathfrak p$ est fidèlement plat, et il en va de même après localisation en $\mathfrak q$. Comme $C^\wedge$, et donc toute localisation de cet anneau, est $(S_2)$, on conclut que $B^\wedge_\mathfrak p$ est $(S_2)$ d'après Algèbre, lemme \ref{algebra-lemma-descent-Sk}. Au total, on trouve que $(S_2)$ est vérifiée pour $B^\wedge$. Pour démontrer que $B^\wedge$ est $(R_1)$, il suffit de considérer les idéaux premiers $\mathfrak q \subset B^\wedge$ tels que $\dim(B^\wedge_\mathfrak q) \leq 1$. Puisque $\dim(B^\wedge_\mathfrak q) = \dim(B_\mathfrak p) +
\dim(B^\wedge_\mathfrak q/\mathfrak p B^\wedge_\mathfrak q)$ d'après Algèbre, lemme \ref{algebra-lemma-dimension-base-fibre-total}, on trouve que $\dim(B_\mathfrak p) \leq 1$ et l'on voit, comme précédemment, que $B^\wedge_\mathfrak q \to C^\wedge_\mathfrak q$ est fidèlement plat. On conclut à l'aide d'Algèbre, lemme \ref{algebra-lemma-descent-Rk}. \end{proof}

\begin{lemma} \label{lemma-regular-alteration-implies-local} Soit $(A, \mathfrak m, \kappa)$ un anneau local noethérien intègre. Supposons qu'il existe une altération $f : X \to \Spec(A)$ avec $X$ régulier. Alors : \begin{enumerate} \item il existe un élément non nul $f \in A$ tel que $A_f$ soit régulier ; \item la clôture intégrale $B$ de $A$ dans son corps des fractions est finie sur $A$ ; \item le complété $\mathfrak m$-adique de $B$ est un anneau normal, c'est-à-dire que les complétés de $B$ en ses idéaux maximaux sont des anneaux normaux intègres ; \item la fibre formelle générique de $A$ est régulière. \end{enumerate} \end{lemma}

\begin{proof} Les parties (1) et (2) résultent du lemme \ref{lemma-regular-alteration-implies}. Nous devons reprendre une partie de sa démonstration afin de fixer les notations nécessaires à celle de (3). Posons $C = \Gamma(X, \mathcal{O}_X)$. D'après Cohomologie des schémas, lemme \ref{coherent-lemma-proper-over-affine-cohomology-finite}, on voit que $C$ est un $A$-module fini. Comme $X$ est normal (Propriétés, lemme \ref{properties-lemma-regular-normal}), on voit que $C$ est un anneau normal intègre (Propriétés, lemme \ref{properties-lemma-normal-integral-sections}). Ainsi $B \subset C$, et l'on conclut que $B$ est fini sur $A$, puisque $A$ est noethérien. D'après le lemme \ref{lemma-algebra-helper}, pour prouver (3), il suffit de montrer que le complété $\mathfrak m$-adique $C^\wedge$ est normal.

\medskip\noindent D'après Algèbre, lemme \ref{algebra-lemma-completion-finite-extension}, le complété $C^\wedge$ est le produit des complétés de $C$ aux idéaux premiers de $C$ situés au-dessus de $\mathfrak m$. Ceux-ci sont en nombre fini et sont les idéaux maximaux $\mathfrak m_1, \ldots, \mathfrak m_r$ de $C$. (Le résultat correspondant pour $B$ explique la dernière assertion du lemme.) En remplaçant donc $A$ par $C_{\mathfrak m_i}$ et $X$ par $X_i = X \times_{\Spec(C)} \Spec(C_{\mathfrak m_i})$, on se ramène au cas traité au paragraphe suivant. (Remarquons que $\Gamma(X_i, \mathcal{O}) = C_{\mathfrak m_i}$ d'après Cohomologie des schémas, lemme \ref{coherent-lemma-flat-base-change-cohomology}.)

\medskip\noindent Dans ce paragraphe, $A$ est un anneau local noethérien normal intègre et $f : X \to \Spec(A)$ est une altération régulière avec $\Gamma(X, \mathcal{O}_X) = A$. Il faut montrer que le complété $A^\wedge$ de $A$ est un anneau normal intègre. D'après le lemme \ref{lemma-port-regularity-to-completion}, $Y = X \times_{\Spec(A)} \Spec(A^\wedge)$ est régulier. Puisque $\Gamma(Y, \mathcal{O}_Y) = A^\wedge$ d'après Cohomologie des schémas, lemme \ref{coherent-lemma-flat-base-change-cohomology}, on conclut comme précédemment que $A^\wedge$ est normal. En effet, $Y$ est normal d'après Propriétés, lemme \ref{properties-lemma-regular-normal}. Il est connexe puisque $\Gamma(Y, \mathcal{O}_Y) = A^\wedge$ est local. Ainsi $Y$ est normal et intègre (car un schéma noethérien connexe et normal est intègre). Par conséquent, $\Gamma(Y, \mathcal{O}_Y) = A^\wedge$ est un anneau normal intègre d'après Propriétés, lemme \ref{properties-lemma-normal-integral-sections}. Cela démontre (3).

\medskip\noindent Démonstration de (4). Notons $\eta \in \Spec(A)$ le point générique et indiquons par un indice $\eta$ le changement de base à $\eta$. Puisque $f$ est une altération, le schéma $X_\eta$ est fini et fidèlement plat sur $\eta$. Comme $Y = X \times_{\Spec(A)} \Spec(A^\wedge)$ est régulier d'après le lemme \ref{lemma-port-regularity-to-completion}, on voit que $Y_\eta$ est régulier (en tant que limite d'ouverts de $Y$). Alors $Y_\eta \to \Spec(A^\wedge \otimes_A \kappa(\eta))$ est fini et fidèlement plat sur la fibre formelle générique. On conclut par Algèbre, lemme \ref{algebra-lemma-descent-regular}. \end{proof}











\section{Résolution} \label{section-resolution}



\noindent Voici une définition.

\begin{definition} \label{definition-resolution} Soit $Y$ un schéma intègre noethérien. Une {\it résolution des singularités} de $Y$ est une modification $f : X \to Y$ telle que $X$ soit régulier. \end{definition}

\noindent Dans le cas des surfaces, nous souhaitons parfois disposer d'un peu plus d'information.

\begin{definition} \label{definition-resolution-surface} Soit $Y$ un schéma intègre noethérien de dimension $2$. Nous disons que $Y$ admet une {\it résolution des singularités par éclatements normalisés} s'il existe une suite $$
Y_n \to Y_{n - 1} \to \ldots \to Y_1 \to Y_0 \to Y
$$ telle que : \begin{enumerate} \item $Y_i$ soit propre sur $Y$ pour $i = 0, \ldots, n$ ; \item $Y_0 \to Y$ soit la normalisation ; \item $Y_i \to Y_{i - 1}$ soit un éclatement normalisé pour $i = 1, \ldots, n$ ; \item $Y_n$ soit régulier. \end{enumerate} \end{definition}

\noindent Remarquons que la condition (1) implique que la normalisation $Y_0$ de $Y$ est finie sur $Y$ et que les normalisations utilisées dans les éclatements normalisés sont elles aussi finies.

\begin{lemma} \label{lemma-existence-implies-existence-by-normalized-blowing-ups} Soit $(A, \mathfrak m, \kappa)$ un anneau local noethérien. Supposons que $A$ soit normal et de dimension $2$. Si $\Spec(A)$ admet une résolution des singularités, alors $\Spec(A)$ admet une résolution par éclatements normalisés. \end{lemma}

\begin{proof} D'après le lemme \ref{lemma-regular-alteration-implies-local}, le complété $A^\wedge$ de $A$ est normal. Le lemme \ref{lemma-port-regularity-to-completion} montre que $\Spec(A^\wedge)$ admet une résolution. D'après le lemme \ref{lemma-normalized-blowup-completion}, toute suite $Y_n \to Y_{n - 1} \to \ldots \to \Spec(A^\wedge)$ d'éclatements normalisés provient d'une suite d'éclatements normalisés $X_n \to \ldots \to \Spec(A)$. De plus, si $Y_n$ est régulier, alors $X_n$ est régulier d'après le lemme \ref{lemma-port-regularity-to-completion}. Il suffit donc de démontrer le lemme lorsque $A$ est complet.

\medskip\noindent Supposons en outre que $A$ soit complet. Nous utiliserons le fait que $A$ est de Nagata (Algèbre, proposition \ref{algebra-proposition-ubiquity-nagata}), excellent (Compléments d'algèbre, proposition \ref{more-algebra-proposition-ubiquity-excellent}) et possède un complexe dualisant (Complexes dualisants, lemme \ref{dualizing-lemma-ubiquity-dualizing}). Il en va de même pour tout anneau essentiellement de type fini sur $A$. Si $B$ est un anneau local normal intègre excellent, alors son complété $B^\wedge$ est normal (car $B \to B^\wedge$ est régulier et Compléments d'algèbre, lemme \ref{more-algebra-lemma-normal-goes-up}, s'applique). Nous l'utiliserons sans plus de commentaires dans la suite de la démonstration.

\medskip\noindent Soit $X \to \Spec(A)$ une résolution des singularités. Choisissons une suite d'éclatements normalisés $$
Y_n \to Y_{n - 1} \to \ldots \to Y_1 \to \Spec(A)
$$ qui domine $X$ (lemme \ref{lemma-dominate-by-normalized-blowing-up}). Le morphisme $Y_n \to X$ est un isomorphisme en dehors d'un nombre fini de points de $X$. Nous pouvons donc appliquer le lemme \ref{lemma-dominate-by-blowing-up-in-points} pour trouver une suite d'éclatements $$
X_m \to X_{m - 1} \to \ldots \to X
$$ en des points fermés telle que $X_m$ domine $Y_n$. On a le diagramme $$
\xymatrix{
& Y_n \ar[rd] \ar[rr] & & \Spec(A) \\
X_m \ar[rr] \ar[ru] & & X \ar[ru]
}
$$. Pour démontrer le lemme, il suffit de montrer qu'un nombre fini d'éclatements normalisés de $Y_n$ produit un schéma régulier. Le diagramme ci-dessus montre que $Y_n$ admet une résolution, à savoir $X_m$. Comme $Y_n$ est une surface normale, cela implique que $Y_n$ possède au plus un nombre fini de singularités $y_1, \ldots, y_t$ (car $X_m \to Y_n$ est un isomorphisme en dehors des fibres de dimension $1$ ; voir Variétés, lemme \ref{varieties-lemma-modification-normal-iso-over-codimension-1}).

\medskip\noindent Soit $x_a \in X$ l'image de $y_a$. Alors $\mathcal{O}_{X, x_a}$ est régulier et définit donc une singularité rationnelle (lemme \ref{lemma-regular-rational}). Appliquons le lemme \ref{lemma-rational-propagates} à $\mathcal{O}_{X, x_a} \to \mathcal{O}_{Y_n, y_a}$ : il en résulte que $\mathcal{O}_{Y_n, y_a}$ définit une singularité rationnelle. D'après le lemme \ref{lemma-rational-to-gorenstein}, il existe une suite finie d'éclatements en des points fermés singuliers $$
Y_{a, n_a} \to Y_{a, n_a - 1} \to \ldots \to \Spec(\mathcal{O}_{Y_n, y_a})
$$ telle que $Y_{a, n_a}$ soit de Gorenstein, c'est-à-dire possède un module dualisant inversible. D'après le lemme \ref{lemma-equivalence-sequence-blowups} (qui est essentiellement trivial), appliqué avec $n' = \sum n_a$, ces suites correspondent à une suite d'éclatements $$
Y_{n + n'} \to Y_{n + n' - 1} \to \ldots \to Y_n
$$ telle que $Y_{n + n'}$ soit normal et que les anneaux locaux de $Y_{n + n'}$ soient de Gorenstein. À l'aide des références données ci-dessus, nous pouvons dominer $Y_{n + n'}$ par une suite d'éclatements $X_{m + m'} \to \ldots \to X_m$ dominant $Y_{n + n'}$, comme dans le diagramme suivant $$
\xymatrix{
 & Y_{n + n'} \ar[rr] & & Y_n \ar[rd] \ar[rr] & & \Spec(A) \\
X_{m + m'} \ar[ru] \ar[rr] & & X_m \ar[rr] \ar[ru] & & X \ar[ru]
}
$$. Ainsi, $Y_{n + n'}$ possède de nouveau un nombre fini de points singuliers $y'_1, \ldots, y'_s$, mais cette fois les singularités sont des points doubles rationnels ; plus précisément, les anneaux locaux $\mathcal{O}_{Y_{n + n'}, y'_b}$ sont comme dans le lemme \ref{lemma-resolve-rational-double-points}. En raisonnant exactement comme ci-dessus, on conclut que le lemme est vrai. \end{proof}

\begin{lemma} \label{lemma-resolve-complete} Soit $(A, \mathfrak m, \kappa)$ un anneau local noethérien complet. Supposons que $A$ soit un anneau normal intègre de dimension $2$. Alors $\Spec(A)$ admet une résolution des singularités. \end{lemma}

\begin{proof} Un anneau local noethérien complet est J-2 (Compléments d'algèbre, proposition \ref{more-algebra-proposition-ubiquity-J-2}), de Nagata (Algèbre, proposition \ref{algebra-proposition-ubiquity-nagata}), excellent (Compléments d'algèbre, proposition \ref{more-algebra-proposition-ubiquity-excellent}) et possède un complexe dualisant (Complexes dualisants, lemme \ref{dualizing-lemma-ubiquity-dualizing}). Il en va de même pour tout anneau essentiellement de type fini sur $A$. Si $B$ est un anneau local normal intègre excellent, alors son complété $B^\wedge$ est normal (car $B \to B^\wedge$ est régulier et Compléments d'algèbre, lemme \ref{more-algebra-lemma-normal-goes-up}, s'applique). Autrement dit, les anneaux locaux rencontrés dans la suite de la démonstration possèdent les propriétés d'excellence dont nous avons besoin.

\medskip\noindent Choisissons $A_0 \subset A$ avec $A_0$ un anneau local régulier complet et $A_0 \to A$ fini ; voir Algèbre, lemme \ref{algebra-lemma-complete-local-Noetherian-domain-finite-over-regular}. Cela induit une extension finie de corps des fractions $K/K_0$. Nous allons raisonner par récurrence sur $[K : K_0]$. Le cas initial est celui où le degré vaut $1$ ; alors $A_0 = A$ et le résultat est vrai.

\medskip\noindent Supposons qu'il existe un corps intermédiaire $K_0 \subset L \subset K$, $K_0 \not = L \not = K$. Soit $B \subset A$ la clôture intégrale de $A_0$ dans $L$. Par récurrence, choisissons une résolution des singularités $Y \to \Spec(B)$. Soit $X$ la normalisation de $Y \times_{\Spec(B)} \Spec(A)$. On a le diagramme $$
\xymatrix{
X \ar[r] \ar[d] & \Spec(A) \ar[d] \\
Y \ar[r] & \Spec(B)
}
$$. Puisque $A$ est J-2, le lieu régulier de $X$ est ouvert. Comme $X$ est une surface normale, on conclut que $X$ possède au plus un nombre fini de points singuliers $x_1, \ldots, x_n$ ; ce sont des points fermés tels que $\dim(\mathcal{O}_{X, x_i}) = 2$. Pour tout $i$, soit $y_i \in Y$ son image. Comme $\mathcal{O}_{Y, y_i}^\wedge \to \mathcal{O}_{X, x_i}^\wedge$ est fini, de degré inférieur au précédent, l'hypothèse de récurrence montre que $\mathcal{O}_{X, x_i}^\wedge$ admet une résolution des singularités. D'après le lemme \ref{lemma-existence-implies-existence-by-normalized-blowing-ups}, il existe une suite $$
Z^\wedge_{i, n_i} \to \ldots \to Z^\wedge_{i, 1} \to
\Spec(\mathcal{O}_{X, x_i}^\wedge)
$$ d'éclatements normalisés telle que $Z^\wedge_{i, n_i}$ soit régulier. Le lemme \ref{lemma-normalized-blowup-completion} fournit une suite correspondante d'éclatements normalisés $$
Z_{i, n_i} \to \ldots \to Z_{i, 1} \to \Spec(\mathcal{O}_{X, x_i})
$$. Alors $Z_{i, n_i}$ est un schéma régulier d'après le lemme \ref{lemma-port-regularity-to-completion}. Le lemme \ref{lemma-equivalence-sequence-normalized-blowups} permet d'insérer ces éclatements normalisés dans une suite correspondante $$
Z_n \to Z_{n - 1} \to \ldots \to Z_1 \to X
$$, et bien sûr $Z_n$ est lui aussi régulier (il suffit de considérer les anneaux locaux). Cela démontre l'étape de récurrence.

\medskip\noindent Supposons qu'il n'existe aucun corps intermédiaire $K_0 \subset L \subset K$ avec $K_0 \not = L \not = K$. Alors, soit $K/K_0$ est séparable, soit la caractéristique de $K$ est $p$ et $[K : K_0] = p$. Le lemme \ref{lemma-go-up-separable} ou le lemme \ref{lemma-go-up-degree-p} implique alors que la réduction aux singularités rationnelles est possible. D'après le lemme \ref{lemma-reduce-to-rational}, il existe donc une modification normale $X \to \Spec(A)$ telle que, pour tout point singulier $x$ de $X$, l'anneau local $\mathcal{O}_{X, x}$ définisse une singularité rationnelle. Puisque $A$ est J-2, on trouve que $X$ ne possède qu'un nombre fini de points singuliers $x_1, \ldots, x_n$. D'après le lemme \ref{lemma-rational-to-gorenstein}, il existe une suite finie d'éclatements en des points fermés singuliers $$
X_{i, n_i} \to X_{i, n_i - 1} \to \ldots \to \Spec(\mathcal{O}_{X, x_i})
$$ telle que $X_{i, n_i}$ soit de Gorenstein, c'est-à-dire possède un module dualisant inversible. D'après le lemme \ref{lemma-equivalence-sequence-blowups} (qui est essentiellement trivial), appliqué avec $n = \sum n_a$, ces suites correspondent à une suite d'éclatements $$
X_n \to X_{n - 1} \to \ldots \to X
$$ telle que $X_n$ soit normal et que les anneaux locaux de $X_n$ soient de Gorenstein. De nouveau, $X_n$ possède un nombre fini de points singuliers $x'_1, \ldots, x'_s$, mais cette fois les singularités sont des points doubles rationnels ; plus précisément, les anneaux locaux $\mathcal{O}_{X_n, x'_i}$ sont comme dans le lemme \ref{lemma-resolve-rational-double-points}. En raisonnant exactement comme ci-dessus, on conclut que le lemme est vrai. \end{proof}

\noindent Nous arrivons enfin au théorème principal de ce chapitre.

\begin{theorem}[Lipman]
\label{theorem-resolve}
\begin{reference}
\cite[Théorème, page 151]{Lipman}
\end{reference}
Soit $Y$ un schéma intègre noethérien de dimension deux. Les assertions suivantes sont équivalentes :
\begin{enumerate}
\item il existe une altération $X \to Y$ avec $X$ régulier ;
\item il existe une résolution des singularités de $Y$ ;
\item $Y$ admet une résolution des singularités par éclatements normalisés ;
\item la normalisation $Y^\nu \to Y$ est finie, $Y^\nu$ ne possède qu'un nombre fini de points singuliers $y_1, \ldots, y_m$ et, pour tout $y_i$, le complété de $\mathcal{O}_{Y^\nu, y_i}$ est normal.
\end{enumerate}
\end{theorem}

\begin{proof} Les implications (3) $\Rightarrow$ (2) $\Rightarrow$ (1) sont immédiates.

\medskip\noindent Soit $X \to Y$ une altération avec $X$ régulier. Alors $Y^\nu \to Y$ est fini d'après le lemme \ref{lemma-regular-alteration-implies}. Considérons la factorisation $f : X \to Y^\nu$ de Morphismes, lemme \ref{morphisms-lemma-normalization-normal}. Le morphisme $f$ est fini au-dessus d'un ouvert $V \subset Y^\nu$ contenant tout point de codimension $\leq 1$ de $Y^\nu$, d'après Variétés, lemme \ref{varieties-lemma-finite-in-codim-1}. Alors $f$ est plat sur $V$ d'après Algèbre, lemme \ref{algebra-lemma-CM-over-regular-flat}, et parce qu'un anneau local normal de dimension $\leq 2$ est de Cohen--Macaulay par le critère de Serre (Algèbre, lemme \ref{algebra-lemma-criterion-normal}). Alors $V$ est régulier d'après Algèbre, lemme \ref{algebra-lemma-descent-regular}. Comme $Y^\nu$ est noethérien, on conclut que $Y^\nu \setminus V = \{y_1, \ldots, y_m\}$ est fini. D'après le lemme \ref{lemma-regular-alteration-implies-local}, le complété de $\mathcal{O}_{Y^\nu, y_i}$ est normal. On voit ainsi que (1) $\Rightarrow$ (4).

\medskip\noindent Supposons (4). Il faut démontrer (3). Nous pouvons immédiatement remplacer $Y$ par sa normalisation. Soient $y_1, \ldots, y_m \in Y$ les points singuliers. En appliquant les lemmes \ref{lemma-resolve-complete} et \ref{lemma-existence-implies-existence-by-normalized-blowing-ups}, on trouve qu'il existe une suite finie d'éclatements normalisés $$
Y_{i, n_i} \to Y_{i, n_i - 1} \to \ldots \to \Spec(\mathcal{O}^\wedge_{Y, y_i})
$$ telle que $Y_{i, n_i}$ soit régulier. D'après le lemme \ref{lemma-normalized-blowup-completion}, il existe une suite correspondante d'éclatements normalisés $$
X_{i, n_i} \to \ldots \to X_{i, 1} \to \Spec(\mathcal{O}_{Y, y_i})
$$. Alors $X_{i, n_i}$ est un schéma régulier d'après le lemme \ref{lemma-port-regularity-to-completion}. Le lemme \ref{lemma-equivalence-sequence-normalized-blowups} permet d'insérer ces éclatements normalisés dans une suite correspondante $$
X_n \to X_{n - 1} \to \ldots \to X_1 \to Y
$$, et bien sûr $X_n$ est lui aussi régulier (il suffit de considérer les anneaux locaux). Cela achève la démonstration. \end{proof}






\section{Résolution plongée} \label{section-embedded}

\noindent Étant donnée une courbe sur une surface, il existe un éclatement qui transforme la courbe en un diviseur à croisements normaux stricts. Dans cette section, nous utiliserons le fait qu'un schéma localement noethérien de dimension un est normal si et seulement s'il est régulier (Algèbre, lemme \ref{algebra-lemma-characterize-dvr}). Nous utiliserons également le fait que tout point d'un schéma localement noethérien se spécialise en un point fermé (Propriétés, lemme \ref{properties-lemma-locally-Noetherian-closed-point}).

\begin{lemma} \label{lemma-resolve-curve} Soit $Y$ un schéma intègre noethérien de dimension un. Les assertions suivantes sont équivalentes : \begin{enumerate} \item il existe une altération $X \to Y$ avec $X$ régulier ; \item il existe une résolution des singularités de $Y$ ; \item il existe une suite finie $Y_n \to Y_{n - 1} \to \ldots \to Y_1 \to Y$ d'éclatements en des points fermés telle que $Y_n$ soit régulier ; \item la normalisation $Y^\nu \to Y$ est finie. \end{enumerate} \end{lemma}

\begin{proof} Les implications (3) $\Rightarrow$ (2) $\Rightarrow$ (1) sont immédiates. L'implication (1) $\Rightarrow$ (4) résulte du lemme \ref{lemma-regular-alteration-implies}. Remarquons qu'un schéma normal de dimension un est régulier ; l'implication (4) $\Rightarrow$ (2) est donc elle aussi claire. Il reste ainsi à montrer que les conditions équivalentes (1), (2) et (4) impliquent (3).

\medskip\noindent Soit $f : X \to Y$ une résolution des singularités. Comme $Y$ est de dimension un, on voit que $f$ est fini d'après Variétés, lemme \ref{varieties-lemma-finite-in-codim-1}. Nous allons construire des factorisations $$
X \to \ldots \to Y_2 \to Y_1 \to Y
$$ dans lesquelles $Y_i \to Y_{i - 1}$ est l'éclatement d'un point fermé et n'est pas un isomorphisme tant que $Y_{i - 1}$ n'est pas régulier. Chacun de ces morphismes sera fini (pour la même raison que ci-dessus), et l'on obtiendra un système correspondant $$
f_*\mathcal{O}_X \supset \ldots \supset
f_{2, *}\mathcal{O}_{Y_2} \supset
f_{1, *}\mathcal{O}_{Y_1} \supset \mathcal{O}_Y
$$, où $f_i : Y_i \to Y$ est le morphisme structural. Puisque $Y$ est noethérien, cette suite croissante de sous-modules cohérents doit se stabiliser (Cohomologie des schémas, lemme \ref{coherent-lemma-acc-coherent}), ce qui montre que, pour un certain $n$, le schéma $Y_n$ est régulier, comme souhaité. Pour construire $Y_i$ à partir de $Y_{i - 1}$, choisissons un point fermé singulier $y_{i - 1} \in Y_{i - 1}$ et soit $Y_i \to Y_{i - 1}$ l'éclatement correspondant. Comme $X$ est régulier de dimension $1$ (ses anneaux locaux aux points fermés sont donc des anneaux de valuation discrète, en particulier des anneaux principaux), le faisceau d'idéaux $\mathfrak m_{y_{i - 1}} \cdot \mathcal{O}_X$ est inversible. La propriété universelle de l'éclatement (Diviseurs, lemme \ref{divisors-lemma-universal-property-blowing-up}) fournit alors une factorisation $X \to Y_i$. Enfin, $Y_i \to Y_{i - 1}$ n'est pas un isomorphisme, puisque $\mathfrak m_{y_{i - 1}}$ n'est pas un idéal inversible. \end{proof}

\begin{lemma} \label{lemma-blowup-curve} Soit $X$ un schéma noethérien. Soit $Y \subset X$ un sous-schéma fermé intègre de dimension $1$ qui satisfait les conditions équivalentes du lemme \ref{lemma-resolve-curve}. Il existe alors une suite finie $$
X_n \to X_{n - 1} \to \ldots \to X_1 \to X
$$ d'éclatements en des points fermés telle que le transformé strict de $Y$ dans $X_n$ soit une courbe régulière. \end{lemma}

\begin{proof} Soit $Y_n \to Y_{n - 1} \to \ldots \to Y_1 \to Y$ la suite d'éclatements fournie par le lemme \ref{lemma-resolve-curve}. Soit $X_n \to X_{n - 1} \to \ldots \to X_1 \to X$ la suite correspondante d'éclatements de $X$. Cela convient, car le transformé strict est l'éclatement d'après Diviseurs, lemme \ref{divisors-lemma-strict-transform}. \end{proof}

\noindent Soit $X$ un schéma localement noethérien. Soient $Y, Z \subset X$ des sous-schémas fermés. Soit $p \in Y \cap Z$ un point fermé. Supposons que $Y$ soit intègre de dimension $1$ et que le point générique de $Y$ ne soit pas contenu dans $Z$. Dans cette situation, on peut considérer l'invariant \begin{equation}
\label{equation-multiplicity}
m_p(Y \cap Z) =
\text{length}_{\mathcal{O}_{X, p}}(\mathcal{O}_{Y \cap Z, p})
\end{equation}. C'est un entier $\geq 1$. En effet, si $I, J \subset \mathcal{O}_{X, p}$ sont les idéaux correspondant à $Y, Z$, on voit que $\mathcal{O}_{Y \cap Z, p} = \mathcal{O}_{X, p}/I + J$ a pour support $\{\mathfrak m_p\}$, car nous avons supposé que $Y \cap Z$ ne contient pas l'unique point de $Y$ qui se spécialise en $p$. La longueur est donc finie d'après Algèbre, lemme \ref{algebra-lemma-support-point}.

\begin{lemma} \label{lemma-blowup-nonsingular-curves-meeting-at-point} Dans la situation ci-dessus, soit $X' \to X$ l'éclatement de $X$ en $p$. Soient $Y', Z' \subset X'$ les transformés stricts de $Y, Z$. Si $\mathcal{O}_{Y, p}$ est régulier, alors : \begin{enumerate} \item $Y' \to Y$ est un isomorphisme ; \item $Y'$ rencontre la fibre exceptionnelle $E \subset X'$ en un point $q$ et $m_q(Y \cap E) = 1$ ; \item si $q \in Z'$ l'est aussi, alors $m_q(Y \cap Z') < m_p(Y \cap Z)$. \end{enumerate} \end{lemma}

\begin{proof} Puisque $\mathcal{O}_{X, p} \to \mathcal{O}_{Y, p}$ est surjectif et que $\mathcal{O}_{Y, p}$ est un anneau de valuation discrète, nous pouvons choisir un élément $x_1 \in \mathfrak m_p$ qui s'envoie sur une uniformisante dans $\mathcal{O}_{Y, p}$. Choisissons un ouvert affine $U = \Spec(A)$ contenant $p$ tel que $x_1 \in A$. Soit $\mathfrak m \subset A$ l'idéal maximal correspondant à $p$. Soient $I, J \subset A$ les idéaux qui définissent $Y, Z$ dans $\Spec(A)$. Après avoir rétréci $U$, nous pouvons supposer que $\mathfrak m = I + (x_1)$, autrement dit que $V(x_1) \cap U \cap Y = \{p\}$ schématiquement. On conclut que $p$ est un diviseur de Cartier effectif sur $Y$ et, puisque $Y'$ est l'éclatement de $Y$ en $p$ (Diviseurs, lemme \ref{divisors-lemma-strict-transform}), on voit que $Y' \to Y$ est un isomorphisme d'après Diviseurs, lemme \ref{divisors-lemma-blow-up-effective-Cartier-divisor}. La relation $\mathfrak m = I + (x_1)$ implique que $\mathfrak m^n \subset I + (x_1^n)$ ; on peut donc définir une application $$
\psi : A[\textstyle{\frac{\mathfrak m}{x_1}}] \longrightarrow A/I
$$ en envoyant $y/x_1^n \in A[\frac{\mathfrak m}{x_1}]$ sur la classe de $a$ dans $A/I$, où $a$ est choisi de sorte que $y \equiv ax_1^n \bmod I$. Alors $\psi$ correspond au morphisme de $Y \cap U$ dans $X'$ au-dessus de $U$ donné par $Y' \cong Y$. Puisque l'image de $x_1$ dans $A[\frac{\mathfrak m}{x_1}]$ définit le diviseur exceptionnel, on conclut que $m_q(Y', E) = 1$. Enfin, comme $J \subset \mathfrak m$, l'idéal $J' \subset A[\frac{\mathfrak m}{x_1}]$ contient certainement les éléments $f/x_1$ pour $f \in J$. Ainsi, si l'on choisit $f \in J$ dont l'image $\overline{f}$ dans $A/I$ a pour valuation minimale $m_p(Y \cap Z)$, on voit que $\psi(f/x_1) = \overline{f}/x_1$ dans $A/I$ a une valuation inférieure d'une unité, ce qui démontre la dernière partie du lemme. \end{proof}

\begin{lemma} \label{lemma-blowup-curves} Soit $X$ un schéma noethérien. Soient $Y_i \subset X$, $i = 1, \ldots, n$ des sous-schémas fermés intègres, chacun de dimension $1$ et satisfaisant les conditions équivalentes du lemme \ref{lemma-resolve-curve}. Il existe alors une suite finie $$
X_n \to X_{n - 1} \to \ldots \to X_1 \to X
$$ d'éclatements en des points fermés telle que les transformés stricts $Y'_i \subset X_n$ de $Y_i$ dans $X_n$ soient des courbes régulières deux à deux disjointes. \end{lemma}

\begin{proof} Il résulte du lemme \ref{lemma-blowup-curve} que nous pouvons supposer que $Y_i$ est une courbe régulière pour $i = 1, \ldots, n$. Pour tous $i \not = j$ et $p \in Y_i \cap Y_j$, nous disposons de l'invariant $m_p(Y_i \cap Y_j)$ (\ref{equation-multiplicity}). Si le maximum de ces nombres est $> 1$, on peut le diminuer (lemme \ref{lemma-blowup-nonsingular-curves-meeting-at-point}) en éclatant tous les points $p$ où ce maximum est atteint. Si le maximum vaut $1$, on peut alors séparer les courbes au moyen du même lemme, en éclatant tous ces points $p$. \end{proof}

\noindent Lorsque notre courbe est contenue dans une surface régulière, nous souhaitons souvent la transformer en un diviseur à croisements normaux.

\begin{lemma} \label{lemma-turn-into-effective-Cartier} Soit $X$ un schéma régulier de dimension $2$. Soit $Z \subset X$ un sous-schéma fermé propre. Il existe une suite $$
X_n \to \ldots \to X_1 \to X
$$ d'éclatements en des points fermés telle que l'image inverse $Z_n$ de $Z$ dans $X_n$ soit un diviseur de Cartier effectif. \end{lemma}

\begin{proof} Soit $D \subset Z$ le plus grand diviseur de Cartier effectif contenu dans $Z$. Alors $\mathcal{I}_Z \subset \mathcal{I}_D$, et le quotient est supporté en des points fermés d'après Diviseurs, lemme \ref{divisors-lemma-codim-1-part}. Nous pouvons donc écrire $\mathcal{I}_Z = \mathcal{I}_{Z'} \mathcal{I}_D$, où $Z' \subset X$ est un sous-schéma fermé qui, ensemblistement, consiste en un nombre fini de points fermés. En appliquant le lemme \ref{lemma-make-ideal-principal}, on trouve une suite d'éclatements comme dans l'énoncé du lemme telle que $\mathcal{I}_{Z'}\mathcal{O}_{X_n}$ soit inversible. Cela démontre le lemme. \end{proof}

\begin{lemma} \label{lemma-embedded-resolution} Soit $X$ un schéma régulier de dimension $2$. Soit $Z \subset X$ un sous-schéma fermé propre tel que toute composante irréductible $Y \subset Z$ de dimension $1$ satisfasse les conditions équivalentes du lemme \ref{lemma-resolve-curve}. Il existe alors une suite $$
X_n \to \ldots \to X_1 \to X
$$ d'éclatements en des points fermés telle que l'image inverse $Z_n$ de $Z$ dans $X_n$ soit un diviseur de Cartier effectif supporté sur un diviseur à croisements normaux stricts. \end{lemma}

\begin{proof} Soit $X' \to X$ l'éclatement d'un point fermé $p$. Alors l'image inverse $Z' \subset X'$ de $Z$ est supportée sur le transformé strict de $Z$ et sur le diviseur exceptionnel. Le diviseur exceptionnel est une courbe régulière (lemme \ref{lemma-blowup}), et le transformé strict $Y'$ de toute composante irréductible $Y$ est soit égal à $Y$, soit l'éclatement de $Y$ en $p$. Ainsi, ce procédé ne produit pas de nouvelles composantes singulières de dimension $1$. Il résulte donc des lemmes \ref{lemma-turn-into-effective-Cartier} et \ref{lemma-blowup-curves} que nous pouvons supposer que $Z$ est un diviseur de Cartier effectif et que toutes les composantes irréductibles $Y$ de $Z$ sont régulières. (Bien sûr, nous ne pouvons pas supposer que les composantes irréductibles sont deux à deux disjointes, car chaque fois que nous éclatons un point de $Z$, nous ajoutons une nouvelle composante irréductible à $Z$, à savoir le diviseur exceptionnel.)

\medskip\noindent Supposons que $Z$ soit un diviseur de Cartier effectif dont les composantes irréductibles $Y_i$ sont régulières. Pour tous $i \not = j$ et $p \in Y_i \cap Y_j$, nous disposons de l'invariant $m_p(Y_i \cap Y_j)$ (\ref{equation-multiplicity}). Si le maximum de ces nombres est $> 1$, nous pouvons le diminuer (lemme \ref{lemma-blowup-nonsingular-curves-meeting-at-point}) en éclatant tous les points $p$ où ce maximum est atteint (remarquons que les « nouveaux » invariants $m_{q_i}(Y'_i \cap E)$ valent toujours $1$). Si le maximum vaut $1$ et si, par exemple, $p \in Y_1 \cap \ldots \cap Y_r$ pour un certain $r > 2$ mais pour aucun des autres, alors, après avoir éclaté $p$, on voit que $Y'_1, \ldots, Y'_r$ ne se rencontrent pas aux points situés au-dessus de $p$, et que $m_{q_i}(Y'_i, E) = 1$, où $Y'_i \cap E = \{q_i\}$. En continuant donc à éclater les points où plus de $3$ composantes de $Z$ se rencontrent, on atteint une situation dans laquelle, pour tout point fermé $p \in X$, on a l'un des trois cas suivants : (a) aucune courbe $Y_i$ ne passe par $p$ ; (b) exactement une courbe $Y_i$ passe par $p$ et $\mathcal{O}_{Y_i, p}$ est régulier ; (c) exactement deux courbes $Y_i$, $Y_j$ passent par $p$, les anneaux locaux $\mathcal{O}_{Y_i, p}$, $\mathcal{O}_{Y_j, p}$ sont réguliers et $m_p(Y_i \cap Y_j) = 1$. Cela signifie que $\sum Y_i$ est un diviseur à croisements normaux stricts sur la surface régulière $X$ ; voir Morphismes \'etales, lemme \ref{etale-lemma-strict-normal-crossings}. \end{proof}






\section{Contraction des courbes exceptionnelles} \label{section-minus-one}

\noindent Soit $X$ un schéma noethérien. Soit $E \subset X$ un sous-schéma fermé possédant les propriétés suivantes : \begin{enumerate} \item $E$ est un diviseur de Cartier effectif sur $X$ ; \item il existe un corps $k$ et un isomorphisme de schémas $\mathbf{P}^1_k \to E$ ; \item l'image inverse du faisceau normal $\mathcal{N}_{E/X}$ par cet isomorphisme est isomorphe à $\mathcal{O}_{\mathbf{P}^1}(-1)$. \end{enumerate} Un tel sous-schéma fermé est appelé une {\it courbe exceptionnelle de première espèce}.

\medskip\noindent Soit $X'$ un schéma noethérien et soit $x \in X'$ un point fermé tel que $\mathcal{O}_{X', x}$ soit régulier de dimension $2$. Soit $b : X \to X'$ l'éclatement de $X'$ en $x$. Dans ce cas, la fibre exceptionnelle $E \subset X$ est une courbe exceptionnelle de première espèce. Cela résulte du lemme \ref{lemma-blowup}.

\medskip\noindent Question : toute courbe exceptionnelle de première espèce est-elle obtenue comme fibre d'un éclatement comme ci-dessus ? Autrement dit, existe-t-il toujours un morphisme propre de schémas $X \to X'$ tel que $E$ soit envoyée sur un point fermé $x \in X'$, que $\mathcal{O}_{X', x}$ soit régulier de dimension $2$ et que $X$ soit l'éclatement de $X'$ en $x$ ? Si tel est le cas, nous disons qu'{\it il existe une contraction de $E$}.

\begin{lemma} \label{lemma-factor-through-contraction} Soit $X$ un schéma noethérien. Soit $E \subset X$ une courbe exceptionnelle de première espèce. S'il existe une contraction $X \to X'$ de $E$, celle-ci possède la propriété universelle suivante : pour tout morphisme $\varphi : X \to Y$ tel que $\varphi(E)$ soit un point, il existe une unique factorisation $X \to X' \to Y$ de $\varphi$. \end{lemma}

\begin{proof} Soit $b : X \to X'$ une contraction de $E$. En tant qu'espace topologique, $X'$ est le quotient de $X$ par la relation qui identifie tous les points de $E$ à un seul point. En effet, $b$ est propre (Diviseurs, lemme \ref{divisors-lemma-blowing-up-projective} et Morphismes, lemme \ref{morphisms-lemma-locally-projective-proper}) et surjectif ; il définit donc une application submersive d'espaces topologiques (Topologie, lemme \ref{topology-lemma-closed-morphism-quotient-topology}). D'autre part, l'application canonique $\mathcal{O}_{X'} \to b_*\mathcal{O}_X$ est un isomorphisme. Cela est clair sur le complémentaire du point $x \in X'$, image de $E$ ; au niveau des germes en $x$, l'application est un isomorphisme d'après la partie (4) du lemme \ref{lemma-cohomology-of-blowup}. Ainsi, le couple $(X', \mathcal{O}_{X'})$ s'obtient à partir de $X$ en prenant le quotient comme espace topologique et en le munissant de $b_*\mathcal{O}_X$ comme faisceau structural.

\medskip\noindent Étant donné $\varphi$, soit $\varphi' : X' \to Y$ l'unique application d'espaces topologiques telle que $\varphi = \varphi' \circ b$. L'application $$
\varphi^\sharp : \varphi^{-1}\mathcal{O}_Y =
b^{-1}((\varphi')^{-1}\mathcal{O}_Y) \to \mathcal{O}_X
$$ est alors adjointe à une application $$
(\varphi')^\sharp :
(\varphi')^{-1}\mathcal{O}_Y \to b_*\mathcal{O}_X = \mathcal{O}_{X'}
$$. Le couple $(\varphi', (\varphi')^\sharp)$ est un morphisme d'espaces annelés de $X'$ vers $Y$ qui fournit la factorisation voulue. Comme $\varphi$ est un morphisme d'espaces localement annelés, il en est de même de $\varphi'$. En effet, la seule chose à vérifier est que l'homomorphisme $\mathcal{O}_{Y, y} \to \mathcal{O}_{X', x}$ est local, où $y \in Y$ est l'image de $E$ par $\varphi$. Cela est vrai, car l'image inverse d'un élément $f \in \mathfrak m_y$ est une fonction sur $X$ qui s'annule en tout point de $E$ ; par conséquent, l'image inverse de $f$ sur $X'$ est une fonction définie sur un voisinage de $x$ dans $X'$ et possédant la même propriété. Il est alors clair que cette fonction doit s'annuler en $x$, comme souhaité. \end{proof}

\begin{lemma} \label{lemma-contraction-unique} Soit $X$ un schéma noethérien. Soit $E \subset X$ une courbe exceptionnelle de première espèce. S'il existe une contraction de $E$, alors elle est unique à isomorphisme unique près. \end{lemma}

\begin{proof} Cela résulte immédiatement de la propriété universelle du lemme \ref{lemma-factor-through-contraction}. \end{proof}

\begin{lemma} \label{lemma-exceptional-first-kind-local} Soit $X$ un schéma noethérien. Soit $E \subset X$ une courbe exceptionnelle de première espèce. Posons $E_n = nE$ et notons $\mathcal{O}_n$ son faisceau structural. Alors $$
A = \lim H^0(E_n, \mathcal{O}_n)
$$ est un anneau local noethérien complet et régulier de dimension $2$, et $\Ker(A \to H^0(E_n, \mathcal{O}_n))$ est la $n$-ième puissance de son idéal maximal. \end{lemma}

\begin{proof} Rappelons qu'il existe un isomorphisme $\mathbf{P}^1_k \to E$ tel que l'image inverse du faisceau normal de $E$ dans $X$ soit isomorphe à $\mathcal{O}(-1)$. Alors $H^0(E, \mathcal{O}_E) = k$. Nous noterons $\mathcal{O}_n(iE)$ la restriction du faisceau inversible $\mathcal{O}_X(iE)$ à $E_n$, pour tous $n \geq 1$ et $i \in \mathbf{Z}$. Rappelons que $\mathcal{O}_X(-nE)$ est le faisceau d'idéaux de $E_n$. Ainsi, pour $d \geq 0$, nous obtenons une suite exacte courte $$
0 \to \mathcal{O}_E(-(d + n)E) \to
\mathcal{O}_{n + 1}(-dE) \to
\mathcal{O}_n(-dE) \to 0
$$. Comme $\mathcal{O}_E(-(d + n)E) = \mathcal{O}_{\mathbf{P}^1_k}(d + n)$, le premier groupe de cohomologie s'annule pour tous $d \geq 0$ et $n \geq 1$. On en déduit que les applications de transition du système $H^0(E_n, \mathcal{O}_n(-dE))$ sont surjectives. Pour $d = 0$, nous obtenons un système projectif de surjections d'anneaux tel que le noyau de chaque application de transition soit un idéal nilpotent. Ainsi, $A = \lim H^0(E_n, \mathcal{O}_n)$ est un anneau local de corps résiduel $k$ et d'idéal maximal $$
\lim \Ker(H^0(E_n, \mathcal{O}_n) \to H^0(E, \mathcal{O}_E)) =
\lim H^0(E_n, \mathcal{O}_n(-E))
$$.
\begingroup\emergencystretch=3em
Choisissons $x, y$ dans ce noyau, dont les images forment une $k$-base de $H^0(E, \mathcal{O}_E(-E)) = H^0(\mathbf{P}^1_k, \mathcal{O}(1))$. Alors $x^d, x^{d - 1}y, \ldots, y^d$ sont des éléments de $\lim H^0(E_n, \mathcal{O}_n(-dE))$ dont les images forment une base de $H^0(E, \mathcal{O}_E(-dE)) = H^0(\mathbf{P}^1_k, \mathcal{O}(d))$. On voit ainsi que $A$ est séparé et complet pour la topologie linéaire définie par les noyaux\par\endgroup
$$
I_n = \Ker(A \longrightarrow H^0(E_n, \mathcal{O}_n))
$$. Nous avons $x, y \in I_1$, $I_d I_{d'} \subset I_{d + d'}$, et $I_d/I_{d + 1}$ est un $k$-module libre de base $x^d, x^{d - 1}y, \ldots, y^d$. Nous allons montrer que $I_d = (x, y)^d$. En effet, si $z_e \in I_e$ avec $e \geq d$, alors nous pouvons écrire $$
z_e = a_{e, 0} x^d + a_{e, 1} x^{d - 1}y + \ldots + a_{e, d}y^d + z_{e + 1}
$$, où $a_{e, j} \in (x, y)^{e - d}$ et $z_{e + 1} \in I_{e + 1}$ d'après notre description de $I_d/I_{d + 1}$. Ainsi, en partant d'un certain $z = z_d \in I_d$, nous pouvons procéder par récurrence et obtenir $$
z = \sum\nolimits_{e \geq d} \sum\nolimits_j a_{e, j} x^{d - j} y^j
$$ avec des $a_{e, j} \in (x, y)^{e - d}$. Alors $a_j = \sum_{e \geq d} a_{e, j}$ existe (par complétude et parce que $a_{e, j} \in I_{e - d}$), et nous avons $z = \sum a_{e, j} x^{d - j} y^j$. Par conséquent, $I_d = (x, y)^d$. Ainsi, $A$ est complet pour la topologie $(x, y)$-adique. Alors $A$ est noethérien d'après Algèbre, lemme \ref{algebra-lemma-completion-Noetherian}. Il est clair que sa dimension vaut $2$ d'après la description de $(x, y)^d/(x, y)^{d + 1}$ et Algèbre, proposition \ref{algebra-proposition-dimension}. Comme son idéal maximal est engendré par deux éléments, il est régulier. \end{proof}

\begin{lemma} \label{lemma-contraction} Soit $X$ un schéma noethérien. Soit $E \subset X$ une courbe exceptionnelle de première espèce. S'il existe un morphisme $f : X \to Y$ tel que \begin{enumerate} \item $Y$ est noethérien, \item $f$ est propre, \item $f$ envoie $E$ sur un point $y$ de $Y$, \item $f$ est quasi-fini en tout point hors de $E$, \end{enumerate} alors il existe une contraction de $E$, et celle-ci est la factorisation de Stein de $f$. \end{lemma}

\begin{proof} Nous appliquons Compléments sur les morphismes, théorème \ref{more-morphisms-theorem-stein-factorization-Noetherian}, pour obtenir une factorisation de Stein $X \to X' \to Y$. Alors $X \to X'$ satisfait toutes les hypothèses du lemme (certains détails sont omis). Après avoir remplacé $Y$ par $X'$, nous pouvons donc supposer en outre que $f_*\mathcal{O}_X = \mathcal{O}_Y$ et que les fibres de $f$ sont géométriquement connexes.

\medskip\noindent Supposons que $f_*\mathcal{O}_X = \mathcal{O}_Y$ et que les fibres de $f$ soient géométriquement connexes. Remarquons que $y \in Y$ est un point fermé, car $f$ est fermé et $E$ est fermé. La restriction $f^{-1}(Y \setminus \{y\}) \to Y \setminus \{y\}$ de $f$ est un morphisme fini (Compléments sur les morphismes, lemme \ref{more-morphisms-lemma-characterize-finite}). Cette restriction est donc un isomorphisme puisque $f_*\mathcal{O}_X = \mathcal{O}_Y$, les morphismes finis étant affines. Pour montrer que $\mathcal{O}_{Y, y}$ est régulier de dimension $2$, considérons l'isomorphisme $$
\mathcal{O}_{Y, y}^\wedge \longrightarrow
\lim H^0(X \times_Y \Spec(\mathcal{O}_{Y, y}/\mathfrak m_y^n), \mathcal{O})
$$ fourni par Cohomologie des schémas, lemme \ref{coherent-lemma-formal-functions-stalk}. Soit $E_n = nE$ comme dans le lemme \ref{lemma-exceptional-first-kind-local}. Observons que $$
E_n \subset X \times_Y \Spec(\mathcal{O}_{Y, y}/\mathfrak m_y^n)
$$ car $E \subset X_y = X \times_Y \Spec(\kappa(y))$. D'autre part, puisque $E = f^{-1}(\{y\})$ ensemblistement (car les fibres de $f$ sont géométriquement connexes), la fibre schématique $X_y$ est schématiquement contenue dans $E_n$ pour un certain $n > 0$. Plus précisément, appliquons le lemme \ref{coherent-lemma-power-ideal-kills-sheaf} de Cohomologie des schémas au $\mathcal{O}_X$-module cohérent $\mathcal{F} = \mathcal{O}_{X_y}$ et au faisceau d'idéaux $\mathcal{I}$ de $E$, puis utilisons le fait que $\mathcal{I}^n$ est le faisceau d'idéaux de $E_n$. Cela montre que $$
X \times_Y \Spec(\mathcal{O}_{Y, y}/\mathfrak m_y^m) \subset E_{nm}
$$ La limite projective affichée ci-dessus est donc égale à $\lim H^0(E_n, \mathcal{O}_n)$, qui est un anneau local régulier de dimension deux d'après le lemme \ref{lemma-exceptional-first-kind-local}. Ainsi, $\mathcal{O}_{Y, y}$ est un anneau local régulier de dimension deux, car son complété l'est (Compléments d'algèbre, lemme \ref{more-algebra-lemma-completion-regular} et \ref{more-algebra-lemma-completion-dimension}).

\medskip\noindent Il reste à démontrer que $f : X \to Y$ est l'éclatement $b : Y' \to Y$ de $Y$ en $y$. Nous invitons le lecteur à trouver sa propre démonstration. Tout d'abord, le lemme \ref{lemma-exceptional-first-kind-local} entraîne également l'égalité schématique $X_y = E$. Comme le faisceau d'idéaux de $E$ est inversible, cela montre que $f^{-1}\mathfrak m_y \cdot \mathcal{O}_X$ est inversible. La propriété universelle de l'éclatement fournit donc une factorisation $$
X \to Y' \to Y
$$ du morphisme $f$ ; voir Diviseurs, lemme \ref{divisors-lemma-universal-property-blowing-up}. Rappelons que la fibre exceptionnelle $E' \subset Y'$ est une courbe exceptionnelle de première espèce d'après le lemme \ref{lemma-blowup}. Soit $g : E \to E'$ le morphisme induit. Comme, tant pour $E'$ que pour $E$, le faisceau conormal est engendré par les images inverses de $a$ et $b$, l'application canonique $g^*\mathcal{C}_{E'/Y'} \to \mathcal{C}_{E/X}$ (Morphismes, lemme \ref{morphisms-lemma-conormal-functorial}) est surjective. Les deux faisceaux étant inversibles, cette application est un isomorphisme. Comme $\mathcal{C}_{E/X}$ est de degré positif, $g$ ne peut être un morphisme constant. Ainsi, $g$ est à fibres finies, et $g$ est donc un morphisme fini (même référence que ci-dessus). Or $Y'$ est régulier, donc normal, en tout point de $E'$, tandis que $X \to Y'$ est birationnel et est un isomorphisme hors de $E'$ ; on conclut alors que $X \to Y'$ est un isomorphisme d'après Variétés, lemme \ref{varieties-lemma-modification-normal-iso-over-codimension-1}. \end{proof}

\begin{lemma} \label{lemma-pic-blowup} Soit $b : X \to X'$ la contraction d'une courbe exceptionnelle de première espèce $E \subset X$. Il existe alors une suite exacte courte $$
0 \to \Pic(X') \to \Pic(X) \to \mathbf{Z} \to 0
$$ où la première application est l'image inverse par $b$ et la seconde associe à $\mathcal{L}$ le degré de $\mathcal{L}$ sur la courbe exceptionnelle $E$. Cette suite est scindée par l'application $n \mapsto \mathcal{O}_X(-nE)$. \end{lemma}

\begin{proof} Puisque $E = \mathbf{P}^1_k$, le groupe de Picard de $E$ est $\mathbf{Z}$ ; voir Diviseurs, lemme \ref{divisors-lemma-Pic-projective-space-UFD}. Nous pouvons donc considérer la dernière application comme $\mathcal{L} \mapsto \mathcal{L}|_E$. Par définition des courbes exceptionnelles de première espèce, le degré de la restriction de $\mathcal{O}_X(E)$ à $E$ est $-1$. En combinant ces remarques, il suffit de montrer que $\Pic(X') \to \Pic(X)$ est injective et que son image est formée des faisceaux inversibles dont la restriction est isomorphe au faisceau trivial $\mathcal{O}_E$ sur $E$.

\medskip\noindent Étant donné un $\mathcal{O}_{X'}$-module inversible $\mathcal{L}'$, nous affirmons que l'application $\mathcal{L}' \to b_*b^*\mathcal{L}'$ est un isomorphisme. Cela est clair partout, sauf éventuellement au point image $x \in X'$ de $E$. Pour le vérifier sur les germes en $x$, nous pouvons remplacer $X'$ par un voisinage ouvert de $x$ et supposer que $\mathcal{L}'$ est $\mathcal{O}_{X'}$. Il faut alors montrer que l'application $\mathcal{O}_{X'} \to b_*\mathcal{O}_X$ est un isomorphisme. Cela résulte de la partie (4) du lemme \ref{lemma-cohomology-of-blowup}.

\medskip\noindent Soit $\mathcal{L}$ un $\mathcal{O}_X$-module inversible tel que $\mathcal{L}|_E = \mathcal{O}_E$. Nous affirmons alors que (1) $b_*\mathcal{L}$ est inversible et que (2) $b^*b_*\mathcal{L} \to \mathcal{L}$ est un isomorphisme. Les assertions (1) et (2) sont claires au-dessus de $X' \setminus \{x\}$. Il suffit donc de les démontrer après le changement de base à $\Spec(\mathcal{O}_{X', x})$. La formation de $b_*$ commute aux changements de base plats (Cohomologie des schémas, lemme \ref{coherent-lemma-flat-base-change-cohomology}), et il en va de même pour $b^*$ et pour la formation de l'application d'adjonction. Or, si $X'$ est le spectre d'un anneau local régulier, alors $\mathcal{L}$ est trivial d'après la description du groupe de Picard donnée au lemme \ref{lemma-blowup-pic}. L'affirmation est donc démontrée.

\medskip\noindent En combinant les affirmations démontrées dans les deux paragraphes précédents, on voit que l'application $\mathcal{L} \mapsto b_*\mathcal{L}$ est inverse de l'application $$
\Pic(X') \longrightarrow \Ker(\Pic(X) \to \Pic(E))
$$, ce qui démontre le lemme. \end{proof}

\begin{remark} \label{remark-pic-blowup} Soit $b : X \to X'$ la contraction d'une courbe exceptionnelle de première espèce $E \subset X$. Le lemme \ref{lemma-pic-blowup} fournit une identification $$
\Pic(X) = \Pic(X') \oplus \mathbf{Z}
$$ dans laquelle $\mathcal{L}$ correspond au couple $(\mathcal{L}', n)$ si et seulement si $\mathcal{L} = (b^*\mathcal{L}')(-nE)$, c'est-à-dire $\mathcal{L}(nE) = b^*\mathcal{L}'$. En fait, la preuve du lemme \ref{lemma-pic-blowup} montre que $\mathcal{L}' = b_*\mathcal{L}(nE)$. Bien entendu, l'application $\mathcal{L} \mapsto \mathcal{L}'$ est un homomorphisme de groupes. \end{remark}

\begin{lemma} \label{lemma-lift-sections-and-h1} Soit $X$ un schéma noethérien. Soit $E \subset X$ une courbe exceptionnelle de première espèce. Soit $\mathcal{L}$ un $\mathcal{O}_X$-module inversible. Soit $n$ l'entier tel que $\mathcal{L}|_E$ soit de degré $n$ lorsqu'on le considère comme module inversible sur $\mathbf{P}^1$. Alors \begin{enumerate} \item Si $H^1(X, \mathcal{L}) = 0$ et $n \geq 0$, alors $H^1(X, \mathcal{L}(iE)) = 0$ pour $0 \leq i \leq n + 1$. \item Si $n \leq 0$, alors $H^1(X, \mathcal{L}) \subset H^1(X, \mathcal{L}(E))$. \end{enumerate} \end{lemma}

\begin{proof} Observons que $\mathcal{L}|_E = \mathcal{O}(n)$ d'après Diviseurs, lemme \ref{divisors-lemma-Pic-projective-space-UFD}. Procédons par récurrence en utilisant la suite exacte longue de cohomologie associée à la suite exacte courte $$
0 \to \mathcal{L} \to \mathcal{L}(E) \to \mathcal{L}(E)|_E \to 0,
$$, ainsi que le fait que $H^1(\mathbf{P}^1, \mathcal{O}(d)) = 0$ pour $d \geq -1$ et que $H^0(\mathbf{P}^1, \mathcal{O}(d)) = 0$ pour $d \leq -1$. Nous omettons quelques détails. \end{proof}

\begin{lemma} \label{lemma-contract-ample} Soit $S = \Spec(R)$ un schéma noethérien affine. Soit $X \to S$ un morphisme propre. Soit $\mathcal{L}$ un faisceau inversible ample sur $X$. Soit $E \subset X$ une courbe exceptionnelle de première espèce. Alors \begin{enumerate} \item il existe une contraction $b : X \to X'$ de $E$, \item $X'$ est propre sur $S$, et \item le $\mathcal{O}_{X'}$-module inversible $\mathcal{L}'$ est ample, où $\mathcal{L}'$ est défini comme dans la remarque \ref{remark-pic-blowup}. \end{enumerate} \end{lemma}

\begin{proof} Soit $n$ le degré de $\mathcal{L}|_E$ comme dans le lemme \ref{lemma-lift-sections-and-h1}. Remarquons que $n > 0$, car $\mathcal{L}$ est ample sur $E$ (Variétés, lemme \ref{varieties-lemma-ample-curve}, et Propriétés, lemme \ref{properties-lemma-ample-on-closed}). Après avoir remplacé $\mathcal{L}$ par une puissance, nous pouvons supposer que $H^i(X, \mathcal{L}^{\otimes e}) = 0$ pour tous $i > 0$ et $e > 0$, voir Cohomologie des schémas, lemme \ref{coherent-lemma-vanshing-gives-ample}. Enfin, après avoir remplacé $\mathcal{L}$ par une autre puissance, nous pouvons supposer qu'il existe des sections globales $t_0, \ldots, t_n$ de $\mathcal{L}$ qui définissent une immersion fermée $\psi : X \to \mathbf{P}^n_S$, voir Morphismes, lemme \ref{morphisms-lemma-finite-type-over-affine-ample-very-ample}.

\medskip\noindent Posons $\mathcal{M} = \mathcal{L}(nE)$. Alors $\mathcal{M}|_E \cong \mathcal{O}_E$. Comme nous avons la suite exacte courte $$
0 \to \mathcal{M}(-E) \to \mathcal{M} \to \mathcal{O}_E \to 0
$$ et que $H^1(X, \mathcal{M}(-E))$ est nul (d'après le lemme \ref{lemma-lift-sections-and-h1} et le fait que $n > 0$), nous pouvons choisir une section $s_{n + 1}$ de $\mathcal{M}$ qui engendre $\mathcal{M}|_E$. Enfin, notons $s_0, \ldots, s_n$ les sections de $\mathcal{M}$ obtenues à partir des sections $t_0, \ldots, t_n$ de $\mathcal{L}$ choisies ci-dessus au moyen de $\mathcal{L} \subset \mathcal{L}(nE) = \mathcal{M}$. Prises ensemble, les sections $s_0, \ldots, s_n, s_{n + 1}$ engendrent $\mathcal{M}$ en tout point de $X$ et définissent donc un morphisme $$
\varphi : X \longrightarrow \mathbf{P}^{n + 1}_S
$$ sur $S$, voir Constructions, lemme \ref{constructions-lemma-projective-space}.

\medskip\noindent Nous allons vérifier ci-dessous les conditions du lemme \ref{lemma-contraction}. Une fois cela fait, nous voyons que la factorisation de Stein $X \to X' \to \mathbf{P}^{n + 1}_S$ de $\varphi$ est la contraction cherchée, ce qui démontre (1). De plus, le morphisme $X' \to \mathbf{P}^{n + 1}_S$ est fini ; par conséquent, $X'$ est propre sur $S$ (Morphismes, lemmes \ref{morphisms-lemma-finite-proper} et \ref{morphisms-lemma-composition-proper}). Cela démontre (2). Remarquons que $X'$ possède un faisceau inversible ample. En effet, l'image inverse $\mathcal{M}'$ de $\mathcal{O}_{\mathbf{P}^{n + 1}_S}(1)$ est ample d'après Morphismes, lemme \ref{morphisms-lemma-pullback-ample-tensor-relatively-ample}. Remarquons que $\mathcal{M}'$ a pour image inverse $\mathcal{M}$ sur $X$ (d'après Constructions, lemme \ref{constructions-lemma-projective-space}). Enfin, $\mathcal{M} = \mathcal{L}(nE)$. Puisque, dans les arguments ci-dessus, nous avons remplacé le faisceau $\mathcal{L}$ initial par une puissance positive, nous concluons que le $\mathcal{O}_{X'}$-module inversible $\mathcal{L}'$ mentionné en (3) du lemme est ample sur $X'$ d'après Propriétés, lemme \ref{properties-lemma-ample-power-ample}.

\medskip\noindent Observations immédiates : $\mathbf{P}^{n + 1}_S$ est noethérien et $\varphi$ est propre. Nous omettons les détails.

\medskip\noindent Remarquons ensuite que tout point de $U = X \setminus E$ est envoyé dans le sous-schéma ouvert $W$ de $\mathbf{P}^{n + 1}_S$ où l'une des $n + 1$ premières coordonnées homogènes est non nulle. En revanche, tout point de $E$ est envoyé sur un point dont les $n + 1$ premières coordonnées homogènes sont toutes nulles, donc en particulier dans le complémentaire de $W$. De plus, il est clair qu'il existe une factorisation $$
U = \varphi^{-1}(W) \xrightarrow{\varphi|_U}
W \xrightarrow{\text{pr}} \mathbf{P}^n_S
$$ de $\psi|_U$, où $\text{pr}$ est la projection définie par les $n + 1$ premières coordonnées et $\psi : X \to \mathbf{P}^n_S$ est l'immersion choisie ci-dessus. Il s'ensuit que $\varphi|_U : U \to W$ est quasi-fini.

\medskip\noindent Enfin, considérons l'application $\varphi|_E : E \to \mathbf{P}^{n + 1}_S$. Pour tout point $x \in E$, l'image $\varphi(x)$ a ses $n + 1$ premières coordonnées nulles ; autrement dit, le morphisme $\varphi|_E$ se factorise par le sous-schéma fermé $\mathbf{P}^0_S \cong S$. Le morphisme $E \to S = \Spec(R)$ se factorise sous la forme $E \to \Spec(H^0(E, \mathcal{O}_E)) \to \Spec(R)$ d'après Schémas, lemme \ref{schemes-lemma-morphism-into-affine}. Comme, par hypothèse, $H^0(E, \mathcal{O}_E)$ est un corps, nous concluons que $E$ s'envoie sur un point de $S \subset \mathbf{P}^{n + 1}_S$, ce qui achève la démonstration. \end{proof}

\begin{lemma} \label{lemma-contract-when-quasi-projective} Soit $S$ un schéma noethérien. Soit $f : X \to S$ un morphisme de type fini. Soit $E \subset X$ une courbe exceptionnelle de première espèce contenue dans une fibre de $f$. \begin{enumerate} \item Si $X$ est projectif sur $S$, alors il existe une contraction $X \to X'$ de $E$ et $X'$ est projectif sur $S$. \item Si $X$ est quasi-projectif sur $S$, alors il existe une contraction $X \to X'$ de $E$ et $X'$ est quasi-projectif sur $S$. \end{enumerate} \end{lemma}

\begin{proof} Les deux cas résultent du lemme \ref{lemma-contract-ample} et des résultats classiques sur les modules inversibles amples et les morphismes (quasi-)projectifs.

\medskip\noindent Démonstration de (1). La projectivité de $f$ signifie que $f$ est propre et qu'il existe un module inversible $f$-ample $\mathcal{L}$, voir Morphismes, lemme \ref{morphisms-lemma-projective-is-quasi-projective-proper} et définition \ref{morphisms-definition-quasi-projective}. Soit $U \subset S$ un ouvert affine contenant l'image de $E$. D'après le lemme \ref{lemma-contract-ample}, il existe une contraction $c : f^{-1}(U) \to V'$ de $E$ et un module inversible ample $\mathcal{N}'$ sur $V'$ dont l'image inverse sur $f^{-1}(U)$ est égale à $\mathcal{L}(nE)|_{f^{-1}(U)}$. Soit $v \in V'$ le point fermé tel que $c$ soit l'éclatement de $v$. Nous pouvons alors recoller $V'$ et $X \setminus E$ le long de $f^{-1}(U) \setminus E = V' \setminus \{v\}$ pour obtenir un schéma $X'$ sur $S$. Les morphismes $c$ et $\text{id}_{X \setminus E}$ se recollent en un morphisme $b : X \to X'$ qui est la contraction de $E$. L'image inverse de $U$ dans $X'$ est propre sur $U$. D'autre part, la restriction de $X' \to S$ au complémentaire de l'image de $v$ dans $S$ est isomorphe à la restriction de $X \to S$ à cet ouvert. Ainsi, $X' \to S$ est propre, puisque la propreté est locale sur la base d'après Morphismes, lemme \ref{morphisms-lemma-proper-local-on-the-base}. Enfin, les restrictions de $\mathcal{N}'$ et de $\mathcal{L}|_{X \setminus E}$ à $f^{-1}(U) \setminus E = V' \setminus \{v\}$ sont isomorphes ; elles se recollent donc en un module inversible $\mathcal{L}'$ sur $X'$. La restriction de $\mathcal{L}'$ à l'image inverse de $U$ dans $X'$ est ample, car il en est ainsi de $\mathcal{N}'$. Sur les ouverts affines de $S$ qui évitent l'image de $v$, il en va de même puisque c'est vrai pour $\mathcal{L}$. Ainsi, $\mathcal{L}'$ est relativement ample pour $(X' \to S)$ d'après Morphismes, lemme \ref{morphisms-lemma-characterize-relatively-ample}, et (1) est démontré.

\medskip\noindent Démonstration de (2). Nous pouvons écrire $X$ comme un sous-schéma ouvert d'un schéma $\overline{X}$ projectif sur $S$ d'après Morphismes, lemme \ref{morphisms-lemma-quasi-projective-open-projective}. D'après (1), il existe une contraction $b : \overline{X} \to \overline{X}'$, et $\overline{X}'$ est projectif sur $S$. Nous notons alors $X' \subset \overline{X}$ l'image de $X \to \overline{X}'$ ; c'est un ouvert puisque $b$ est un isomorphisme hors de $E$. Alors $X \to X'$ est la contraction cherchée. Remarquons que $X'$ est quasi-projectif sur $S$, car il possède un module inversible ample relativement à $S$, par la construction effectuée dans la démonstration de la partie (1). \end{proof}

\begin{lemma} \label{lemma-regular-dim-2-quasi-projective} Soit $S$ un schéma noethérien. Soit $f : X \to S$ un morphisme séparé de type fini, avec $X$ régulier de dimension $2$. Alors $X$ est quasi-projectif sur $S$. \end{lemma}

\begin{proof} Par le lemme de Chow (Cohomologie des schémas, lemme \ref{coherent-lemma-chow-Noetherian}), il existe un morphisme propre $\pi : X' \to X$ qui est un isomorphisme au-dessus d'un ouvert dense $U \subset X$ tel que $X' \to S$ soit H-quasi-projectif. D'après le lemme \ref{lemma-extend-rational-map-blowing-up}, il existe une suite d'éclatements en des points fermés $$
X_n \to \ldots \to X_1 \to X_0 = X
$$ et un $S$-morphisme $X_n \to X'$ prolongeant l'application rationnelle $U \to X'$. Remarquons que $X_n \to X$ est projectif d'après Diviseurs, lemme \ref{divisors-lemma-blowing-up-projective}, et Morphismes, lemme \ref{morphisms-lemma-composition-projective}. Il s'ensuit que $X_n \to X'$ est projectif d'après Morphismes, lemme \ref{morphisms-lemma-projective-permanence}. Par conséquent, $X_n \to S$ est quasi-projectif d'après Morphismes, lemme \ref{morphisms-lemma-composition-quasi-projective} (et le fait qu'un morphisme projectif est quasi-projectif ; voir Morphismes, lemme \ref{morphisms-lemma-projective-quasi-projective}). D'après le lemme \ref{lemma-contract-when-quasi-projective} (et l'unicité des contractions, lemme \ref{lemma-contraction-unique}), nous concluons que $X_{n - 1}, \ldots, X_0 = X$ sont quasi-projectifs sur $S$, comme souhaité. \end{proof}

\begin{lemma} \label{lemma-regular-dim-2-projective} Soit $S$ un schéma noethérien. Soit $f : X \to S$ un morphisme propre tel que $X$ soit régulier de dimension $2$. Alors $X$ est projectif sur $S$. \end{lemma}

\begin{proof} Cela résulte du lemme \ref{lemma-regular-dim-2-quasi-projective} et de Morphismes, lemme \ref{morphisms-lemma-projective-is-quasi-projective-proper}. \end{proof}






\section{Factorisation des morphismes birationnels} \label{section-factorizing}

\noindent Les morphismes propres birationnels entre surfaces non singulières sont des composés de transformations quadratiques.

\begin{lemma} \label{lemma-proper-birational-regular-surfaces} Soit $f : X \to Y$ un morphisme propre birationnel entre schémas noethériens intègres réguliers de dimension $2$. Alors $f$ est un composé d'éclatements en des points fermés. \end{lemma}

\begin{proof} Soit $V \subset Y$ le plus grand ouvert au-dessus duquel $f$ est un isomorphisme. Alors $V$ contient tous les points de codimension $1$ de $V$ (Variétés, lemme \ref{varieties-lemma-modification-normal-iso-over-codimension-1}). Soit $y \in Y$ un point fermé qui n'appartient pas à $V$. Nous voulons montrer que $f$ se factorise par l'éclatement $b : Y' \to Y$ de $Y$ en $y$. En effet, si tel est le cas, au moins une composante (et, en fait, exactement une) de la fibre $f^{-1}(y)$ s'enverra isomorphiquement sur la courbe exceptionnelle de $Y'$, et le nombre de courbes contenues dans les fibres de $X \to Y'$ sera strictement inférieur au nombre de courbes contenues dans les fibres de $X \to Y$ ; nous concluons donc par récurrence. Certains détails sont omis.

\medskip\noindent D'après le lemme \ref{lemma-extend-rational-map-blowing-up}, il existe une suite d'éclatements $$
X' = X_n \to X_{n - 1} \to \ldots \to X_1 \to X_0 = X
$$ en des points fermés situés au-dessus de la fibre $f^{-1}(y)$ et un morphisme $X' \to Y'$ tels que $$
\xymatrix{
X' \ar[d]_{f'} \ar[r] & X \ar[d]^f \\
Y' \ar[r] & Y
}
$$ soit commutatif. Nous voulons montrer que le morphisme $X' \to Y'$ se factorise par $X$ ; nous pourrons alors raisonner par récurrence sur $n$ pour nous ramener au cas où $X' \to X$ est l'éclatement de $X$ en un point fermé $x \in X$ qui s'envoie sur $y$.

\medskip\noindent Soit $E \subset X'$ la fibre exceptionnelle de l'éclatement $X' \to X$. Si $E$ s'envoie sur un point de $Y'$, alors le lemme \ref{lemma-factor-through-contraction} fournit la factorisation souhaitée. Nous allons montrer que, dans le cas contraire, nous obtenons une contradiction. En effet, si $f'(E)$ n'est pas un point, alors $E' = f'(E)$ est nécessairement la courbe exceptionnelle de $Y'$. Considérons le diagramme $$
\xymatrix{
E \ar[r] \ar[d]_g & X' \ar[d]_{f'} \ar[r] & X \ar[d]^f \\
E' \ar[r] & Y' \ar[r] & Y
}
$$ En raisonnant comme précédemment, $f'$ est un isomorphisme dans un voisinage ouvert du point générique de $E'$. Ainsi, $g : E \to E'$ est un morphisme fini birationnel. L'inverse de $g$, qui est une application rationnelle, est alors partout défini d'après Morphismes, lemme \ref{morphisms-lemma-extend-across}, et $g$ est un isomorphisme. Considérons l'application $$
g^*\mathcal{C}_{E'/Y'} \longrightarrow \mathcal{C}_{E/X'}
$$ de Morphismes, lemme \ref{morphisms-lemma-conormal-functorial}. Comme la source et le but sont des modules inversibles de degré $1$ sur $E = E' = \mathbf{P}^1_\kappa$ et que l'application est non nulle (puisque $f'$ est un isomorphisme au point générique de $E$), nous concluons que c'est un isomorphisme. D'après Morphismes, lemme \ref{morphisms-lemma-two-immersions}, nous concluons que $\Omega_{X'/Y'}|_E = 0$. Cela signifie que $f'$ est non ramifié en tout point de $E$ (Morphismes, lemme \ref{morphisms-lemma-unramified-at-point}). Ainsi, $f'$ est quasi-fini en tout point de $E$ (Morphismes, lemme \ref{morphisms-lemma-unramified-quasi-finite}). Par conséquent, le plus grand ouvert $V' \subset Y'$ au-dessus duquel $f'$ est un isomorphisme contient $E'$ d'après Variétés, lemme \ref{varieties-lemma-modification-normal-iso-over-codimension-1}. Il s'ensuit que l'image inverse de $y$ dans $X'$ est $E'$. L'image inverse de $y$ dans $X$ est donc $x$. Par conséquent, $x \in X$ appartient au plus grand ouvert au-dessus duquel $f$ est un isomorphisme, d'après Variétés, lemme \ref{varieties-lemma-modification-normal-iso-over-codimension-1}. Cela contredit l'hypothèse selon laquelle $y$ n'appartient pas à cet ouvert. \end{proof}

\begin{lemma} \label{lemma-birational-regular-surfaces} Soit $S$ un schéma noethérien. Soient $X$ et $Y$ des schémas propres intègres sur $S$, réguliers de dimension $2$. Alors $X$ et $Y$ sont $S$-birationnels si et seulement s'il existe un diagramme de $S$-morphismes $$
X = X_0 \leftarrow X_1 \leftarrow \ldots \leftarrow X_n = Y_m
\to \ldots \to Y_1 \to Y_0 = Y
$$ dans lequel chaque morphisme est l'éclatement d'un point fermé. \end{lemma}

\begin{proof} Soit $U \subset X$ un ouvert et soit $f : U \to Y$ l'application $S$-rationnelle donnée (qui est inversible comme application $S$-rationnelle). D'après le lemme \ref{lemma-extend-rational-map-blowing-up}, nous pouvons factoriser $f$ sous la forme $X_n \to \ldots \to X_1 \to X_0 = X$ et $f_n : X_n \to Y$. Comme $X_n$ est propre sur $S$ et $Y$ séparé sur $S$, le morphisme $f_n$ est propre. Il est clair que $f_n$ est birationnel. Le morphisme $f_n$ est donc un composé de contractions d'après le lemme \ref{lemma-proper-birational-regular-surfaces}. Nous omettons la démonstration de la réciproque. \end{proof}






\begin{multicols}{2}[\section{Autres chapitres}]
\noindent
Préliminaires
\begin{enumerate}
\item \hyperref[introduction-section-phantom]{Introduction}
\item \hyperref[conventions-section-phantom]{Conventions}
\item \hyperref[sets-section-phantom]{Théorie des ensembles}
\item \hyperref[categories-section-phantom]{Catégories}
\item \hyperref[topology-section-phantom]{Topologie}
\item \hyperref[sheaves-section-phantom]{Faisceaux sur les espaces}
\item \hyperref[sites-section-phantom]{Sites et faisceaux}
\item \hyperref[stacks-section-phantom]{Champs}
\item \hyperref[fields-section-phantom]{Corps}
\item \hyperref[algebra-section-phantom]{Algèbre commutative}
\item \hyperref[brauer-section-phantom]{Groupes de Brauer}
\item \hyperref[homology-section-phantom]{Algèbre homologique}
\item \hyperref[derived-section-phantom]{Catégories dérivées}
\item \hyperref[simplicial-section-phantom]{Méthodes simpliciales}
\item \hyperref[more-algebra-section-phantom]{Compléments d'algèbre}
\item \hyperref[smoothing-section-phantom]{Lissage des morphismes d'anneaux}
\item \hyperref[modules-section-phantom]{Faisceaux de modules}
\item \hyperref[sites-modules-section-phantom]{Modules sur les sites}
\item \hyperref[injectives-section-phantom]{Objets injectifs}
\item \hyperref[cohomology-section-phantom]{Cohomologie des faisceaux}
\item \hyperref[sites-cohomology-section-phantom]{Cohomologie sur les sites}
\item \hyperref[dga-section-phantom]{Algèbre différentielle graduée}
\item \hyperref[dpa-section-phantom]{Algèbre à puissances divisées}
\item \hyperref[sdga-section-phantom]{Faisceaux différentiels gradués}
\item \hyperref[hypercovering-section-phantom]{Hyperrecouvrements}
\end{enumerate}
Schémas
\begin{enumerate}
\setcounter{enumi}{25}
\item \hyperref[schemes-section-phantom]{Schémas}
\item \hyperref[constructions-section-phantom]{Constructions de schémas}
\item \hyperref[properties-section-phantom]{Propriétés des schémas}
\item \hyperref[morphisms-section-phantom]{Morphismes de schémas}
\item \hyperref[coherent-section-phantom]{Cohomologie des schémas}
\item \hyperref[divisors-section-phantom]{Diviseurs}
\item \hyperref[limits-section-phantom]{Limites de schémas}
\item \hyperref[varieties-section-phantom]{Variétés}
\item \hyperref[topologies-section-phantom]{Topologies sur les schémas}
\item \hyperref[descent-section-phantom]{Descente}
\item \hyperref[perfect-section-phantom]{Catégories dérivées de schémas}
\item \hyperref[more-morphisms-section-phantom]{Compléments sur les morphismes}
\item \hyperref[flat-section-phantom]{Compléments sur la platitude}
\item \hyperref[groupoids-section-phantom]{Schémas en groupoïdes}
\item \hyperref[more-groupoids-section-phantom]{Compléments sur les schémas en groupoïdes}
\item \hyperref[etale-section-phantom]{Morphismes étales de schémas}
\end{enumerate}
Thèmes de théorie des schémas
\begin{enumerate}
\setcounter{enumi}{41}
\item \hyperref[chow-section-phantom]{Homologie de Chow}
\item \hyperref[intersection-section-phantom]{Théorie de l'intersection}
\item \hyperref[pic-section-phantom]{Schémas de Picard des courbes}
\item \hyperref[weil-section-phantom]{Théories cohomologiques de Weil}
\item \hyperref[adequate-section-phantom]{Modules adéquats}
\item \hyperref[dualizing-section-phantom]{Complexes dualisants}
\item \hyperref[duality-section-phantom]{Dualité pour les schémas}
\item \hyperref[discriminant-section-phantom]{Discriminants et différentes}
\item \hyperref[derham-section-phantom]{Cohomologie de de Rham}
\item \hyperref[local-cohomology-section-phantom]{Cohomologie locale}
\item \hyperref[algebraization-section-phantom]{Géométrie algébrique et formelle}
\item \hyperref[curves-section-phantom]{Courbes algébriques}
\item \hyperref[resolve-section-phantom]{Résolution des surfaces}
\item \hyperref[models-section-phantom]{Réduction semi-stable}
\item \hyperref[functors-section-phantom]{Foncteurs et morphismes}
\item \hyperref[equiv-section-phantom]{Catégories dérivées de variétés}
\item \hyperref[pione-section-phantom]{Groupes fondamentaux de schémas}
\item \hyperref[etale-cohomology-section-phantom]{Cohomologie étale}
\item \hyperref[crystalline-section-phantom]{Cohomologie cristalline}
\item \hyperref[proetale-section-phantom]{Cohomologie pro-étale}
\item \hyperref[relative-cycles-section-phantom]{Cycles relatifs}
\item \hyperref[more-etale-section-phantom]{Compléments de cohomologie étale}
\item \hyperref[trace-section-phantom]{La formule des traces}
\end{enumerate}
Espaces algébriques
\begin{enumerate}
\setcounter{enumi}{64}
\item \hyperref[spaces-section-phantom]{Espaces algébriques}
\item \hyperref[spaces-properties-section-phantom]{Propriétés des espaces algébriques}
\item \hyperref[spaces-morphisms-section-phantom]{Morphismes d'espaces algébriques}
\item \hyperref[decent-spaces-section-phantom]{Espaces algébriques décents}
\item \hyperref[spaces-cohomology-section-phantom]{Cohomologie des espaces algébriques}
\item \hyperref[spaces-limits-section-phantom]{Limites d'espaces algébriques}
\item \hyperref[spaces-divisors-section-phantom]{Diviseurs sur les espaces algébriques}
\item \hyperref[spaces-over-fields-section-phantom]{Espaces algébriques sur des corps}
\item \hyperref[spaces-topologies-section-phantom]{Topologies sur les espaces algébriques}
\item \hyperref[spaces-descent-section-phantom]{Descente et espaces algébriques}
\item \hyperref[spaces-perfect-section-phantom]{Catégories dérivées d'espaces}
\item \hyperref[spaces-more-morphisms-section-phantom]{Compléments sur les morphismes d'espaces}
\item \hyperref[spaces-flat-section-phantom]{Platitude sur les espaces algébriques}
\item \hyperref[spaces-groupoids-section-phantom]{Groupoïdes dans les espaces algébriques}
\item \hyperref[spaces-more-groupoids-section-phantom]{Compléments sur les groupoïdes dans les espaces}
\item \hyperref[bootstrap-section-phantom]{Amorçage}
\item \hyperref[spaces-pushouts-section-phantom]{Sommes amalgamées d'espaces algébriques}
\end{enumerate}
Thèmes de géométrie
\begin{enumerate}
\setcounter{enumi}{81}
\item \hyperref[spaces-chow-section-phantom]{Groupes de Chow des espaces}
\item \hyperref[groupoids-quotients-section-phantom]{Quotients de groupoïdes}
\item \hyperref[spaces-more-cohomology-section-phantom]{Compléments sur la cohomologie des espaces}
\item \hyperref[spaces-simplicial-section-phantom]{Espaces simpliciaux}
\item \hyperref[spaces-duality-section-phantom]{Dualité pour les espaces}
\item \hyperref[formal-spaces-section-phantom]{Espaces algébriques formels}
\item \hyperref[restricted-section-phantom]{Algébrisation des espaces formels}
\item \hyperref[spaces-resolve-section-phantom]{Résolution des surfaces revisitée}
\end{enumerate}
Théorie des déformations
\begin{enumerate}
\setcounter{enumi}{89}
\item \hyperref[formal-defos-section-phantom]{Théorie formelle des déformations}
\item \hyperref[defos-section-phantom]{Théorie des déformations}
\item \hyperref[cotangent-section-phantom]{Le complexe cotangent}
\item \hyperref[examples-defos-section-phantom]{Problèmes de déformation}
\end{enumerate}
Champs algébriques
\begin{enumerate}
\setcounter{enumi}{93}
\item \hyperref[algebraic-section-phantom]{Champs algébriques}
\item \hyperref[examples-stacks-section-phantom]{Exemples de champs}
\item \hyperref[stacks-sheaves-section-phantom]{Faisceaux sur les champs algébriques}
\item \hyperref[criteria-section-phantom]{Critères de représentabilité}
\item \hyperref[artin-section-phantom]{Axiomes d'Artin}
\item \hyperref[quot-section-phantom]{Espaces de Quot et de Hilbert}
\item \hyperref[stacks-properties-section-phantom]{Propriétés des champs algébriques}
\item \hyperref[stacks-morphisms-section-phantom]{Morphismes de champs algébriques}
\item \hyperref[stacks-limits-section-phantom]{Limites de champs algébriques}
\item \hyperref[stacks-cohomology-section-phantom]{Cohomologie des champs algébriques}
\item \hyperref[stacks-perfect-section-phantom]{Catégories dérivées de champs}
\item \hyperref[stacks-introduction-section-phantom]{Introduction aux champs algébriques}
\item \hyperref[stacks-more-morphisms-section-phantom]{Compléments sur les morphismes de champs}
\item \hyperref[stacks-geometry-section-phantom]{La géométrie des champs}
\end{enumerate}
Thèmes de théorie des espaces de modules
\begin{enumerate}
\setcounter{enumi}{107}
\item \hyperref[moduli-section-phantom]{Champs de modules}
\item \hyperref[moduli-curves-section-phantom]{Espaces de modules de courbes}
\end{enumerate}
Divers
\begin{enumerate}
\setcounter{enumi}{109}
\item \hyperref[examples-section-phantom]{Exemples}
\item \hyperref[exercises-section-phantom]{Exercices}
\item \hyperref[guide-section-phantom]{Guide bibliographique}
\item \hyperref[desirables-section-phantom]{Desiderata}
\item \hyperref[coding-section-phantom]{Style de programmation}
\item \hyperref[obsolete-section-phantom]{Obsolète}
\item \hyperref[fdl-section-phantom]{Licence de documentation libre GNU}
\item \hyperref[index-section-phantom]{Index généré automatiquement}
\end{enumerate}
\end{multicols}


\bibliography{my}
\bibliographystyle{amsalpha}

\end{document}
